% Leçon 15 — Expression orale : recyclage et paludisme — Chinois Tle A — Cours signé OnBuch+
% Programme officiel MINESEC. Compiler avec : tectonic ZH15-oral-recyclage-et-paludisme.tex
\newcommand{\DOCMATIERE}{Chinois}
\newcommand{\DOCNIVEAU}{Tle A}
\newcommand{\DOCMODULE}{Module 2 : Environnement et santé}
\newcommand{\DOCLECON}{ZH15}
\newcommand{\DOCTITRE}{Leçon 15 — Expression orale : recyclage et paludisme}
% OnBuch+ — préambule « cours signés » ENRICHI (compilé avec tectonic / XeLaTeX)
% Système visuel « Pop » d'OnBuch+ : fond crème (jamais blanc), bordures encre
% épaisses, ombres « dures » décalées, titres Archivo Black en capitales.
% Version MATHÉMATIQUES : graphes (pgfplots/tikz), boîtes pédagogiques
% (définition, propriété, méthode, attention, le savais-tu, exercice résolu,
% exercice, TP, bilan), page de garde et signature OnBuch+ ancrée en bas.
%
% Macros attendues dans main.tex AVANT \input{preamble} :
%   \DOCMATIERE  \DOCNIVEAU  \DOCMODULE  \DOCLECON (numéro)  \DOCTITRE
\documentclass[11pt]{article}

\usepackage{fontspec}
\usepackage[french]{babel} % français = langue principale ; le chinois via \zh{..} / textebox (xeCJK)
\usepackage{xeCJK}
\usepackage{pifont} % \ding{51} \ding{55} utilisés par les modèles
\usepackage[a4paper,margin=2cm,top=3.4cm,bottom=2.8cm,headheight=1.4cm,footskip=1.3cm]{geometry}
\usepackage{amsmath,amssymb,amsfonts}
\usepackage{mathtools} % \xmapsto, \coloneqq... souvent utilisés par les modèles
\usepackage{cancel} % simplifications barrées (bilans de forces)
% \abs{z} (module, valeur absolue) et \norm{u} (norme), utilisés par les modèles
\providecommand{\abs}{}\let\abs\relax\DeclarePairedDelimiter{\abs}{\lvert}{\rvert}
\providecommand{\norm}{}\let\norm\relax\DeclarePairedDelimiter{\norm}{\lVert}{\rVert}
\usepackage[table]{xcolor} % [table] : fournit aussi \rowcolors (alternance de lignes)
\usepackage{pagecolor}
\usepackage{tikz}
\usetikzlibrary{arrows.meta,positioning,calc,shapes.geometric,shapes.misc,shapes.arrows,decorations.pathmorphing,decorations.pathreplacing,decorations.markings,patterns,shadows,mindmap,trees,fit,backgrounds,matrix,intersections,angles,quotes}
\usepackage{pgfplots}
\usepackage[version=4]{mhchem} % équations nucléaires \ce{^{235}_{92}U}
\usepackage[european,siunitx]{circuitikz} % schémas électriques
\pgfplotsset{compat=1.18}
\usepgfplotslibrary{fillbetween,groupplots} % aires entre courbes (calcul intégral)
\usepackage{enumitem}
\usepackage{fancyhdr}
% [most] charge toutes les bibliothèques tcolorbox (listings, minted, raster,
% documentation...) et gonfle inutilement la mémoire TeX sur un document long
% (~300 boîtes) : on ne charge que ce qui est réellement utilisé.
\usepackage[skins,breakable]{tcolorbox}
\usepackage{graphicx}
\usepackage{colortbl}
\usepackage{booktabs}
\usepackage{tabularx}
\usepackage{diagbox} % case d'en-tête diagonale des tableaux à double entrée
% Colonnes extensibles centrée / gauche / droite (C, L, R), souvent utilisées par les modèles
\newcolumntype{C}{>{\centering\arraybackslash}X}
\newcolumntype{L}{>{\raggedright\arraybackslash}X}
\newcolumntype{R}{>{\raggedleft\arraybackslash}X}
\usepackage{array}
\usepackage{multicol}
\usepackage{multirow}
\usepackage{siunitx}
% parse-numbers=false : accepte aussi \SI{2,5 \times 10^{-3}}{...}
\sisetup{locale=FR,output-decimal-marker={,},group-minimum-digits=4,parse-numbers=false}
\DeclareSIUnit\ohm{\ensuremath{\symup{\Omega}}} % U+2126 absent de la police
\DeclareSIUnit\division{div} % réglages d'oscilloscope (ms/div, V/div)
\DeclareSIUnit\tr{tr} % tours (vitesse de rotation, ex. \SI{1500}{\tr\per\minute})
\DeclareSIUnit\molar{M} % concentration molaire (ex. \SI{3}{\molar}, \SI{1}{\micro\molar})
\DeclareSIUnit\an{an} % année (le modèle confond parfois avec \year, un compteur TeX) — ex. \SI{5}{\kilo\meter\per\an}
\DeclareSIUnit\FCFA{FCFA} % franc CFA (ex. \SI{500}{\FCFA\per\kilo\gram})
\DeclareSIUnit\degreeBrix{\textdegree Brix} % teneur en sucre d'un sirop/jus (réfractométrie)
\DeclareSIUnit\month{mois} % le modèle utilise parfois \month, un compteur TeX, comme unité de durée
\usepackage{qrcode}
% \degree : commande fréquemment utilisée par les modèles pour un angle ou une
% température en dehors de \SI{}{} (non fournie par les paquets chargés).
\providecommand{\degree}{^{\circ}}
% \qed (fin de démonstration), utilisé par les modèles sans amsthm
\providecommand{\qed}{\hfill$\square$}
% \barre : double barre des valeurs interdites dans un tableau de variations (tkz-tab)
\providecommand{\barre}{\ensuremath{\|}}
% environnement proof (amsthm non chargé) utilisé par les modèles
\ifdefined\proof\else\newenvironment{proof}[1][Démonstration]{\par\noindent\textit{#1.}\ }{\qed\par}\fi
\usepackage{lastpage}
\usepackage{hyperref}
\hypersetup{hidelinks}

% ============================================================
% Palette « Pop » OnBuch+ — mêmes hex que la classe P (plus_ui.dart)
% ============================================================
\definecolor{poporange}{HTML}{F59321}
\definecolor{popdark}{HTML}{E07A0C}
\definecolor{popcream}{HTML}{FFF6E8}
\definecolor{popink}{HTML}{1C1714}
\definecolor{poppurple}{HTML}{7A5AE0}
\definecolor{popgreen}{HTML}{2E9E6A}
\definecolor{poppink}{HTML}{E0457B}
\definecolor{popblue}{HTML}{2F7DE1}
\definecolor{popgold}{HTML}{C98A12}
\definecolor{popmuted}{HTML}{8A7F76}
\definecolor{popline}{HTML}{EADFD2}
\definecolor{popgray}{HTML}{8A7F76}
\colorlet{popgrey}{popgray}
% Alias tolérants (le modèle nomme parfois les couleurs différemment)
\colorlet{popwhite}{white}
\colorlet{popblack}{popink}
\colorlet{popred}{poppink}
\colorlet{popyellow}{popgold}
% Teintes claires (fonds de tableaux/encadrés) — le modèle nomme parfois
% les couleurs avec un suffixe L (ex. popblueL) sans les avoir définies.
\colorlet{poporangeL}{poporange!20}
\colorlet{popdarkL}{popdark!20}
\colorlet{poppurpleL}{poppurple!20}
\colorlet{popgreenL}{popgreen!20}
\colorlet{poppinkL}{poppink!20}
\colorlet{popblueL}{popblue!20}
\colorlet{popgoldL}{popgold!20}
\colorlet{popredL}{popred!20}
\colorlet{popgrayL}{popgray!20}
% Teintes claires pour fonds de tableaux / zones
\colorlet{popblueL}{popblue!10!white}
\colorlet{poporangeL}{poporange!14!white}
\colorlet{popgreenL}{popgreen!12!white}
\colorlet{poppurpleL}{poppurple!12!white}
\colorlet{poppinkL}{poppink!12!white}
\colorlet{popgoldL}{popgold!14!white}

\pagecolor{popcream}

% ============================================================
% Typographie
% ============================================================
\defaultfontfeatures{Ligatures=TeX}
\setmainfont{PlusJakartaSans-Medium.ttf}[Path=../../../fonts/,BoldFont=PlusJakartaSans-Bold.ttf]
% Symboles, API et emojis absents de la police principale : repli sur DejaVu Sans (caractères actifs)
\newfontface\symfont{DejaVuSans.ttf}[Path=../../../fonts/]
\makeatletter
\newcommand\symactive[1]{\catcode#1=13 \begingroup\lccode`\~=#1 \lowercase{\endgroup\def~}{{\symfont\char#1}}}
\newcommand\symblank[1]{\catcode#1=13 \begingroup\lccode`\~=#1 \lowercase{\endgroup\def~}{}}
\newcommand\symhyph[1]{\catcode#1=13 \begingroup\lccode`\~=#1 \lowercase{\endgroup\def~}{\mbox{-}}}
\newcount\symc
\newcommand\symrange[2]{\symc=#1 \@whilenum\symc<#2 \do{\expandafter\symactive\expandafter{\the\symc}\advance\symc by 1 }}
\newcommand\symblankrange[2]{\symc=#1 \@whilenum\symc<#2 \do{\expandafter\symblank\expandafter{\the\symc}\advance\symc by 1 }}
\makeatother
\symrange{"0250}{"02FF}\symactive{"2713}\symactive{"2714}\symactive{"2717}\symactive{"2718}\symactive{"2610}\symactive{"2611}\symactive{"2612}
\symrange{"2600}{"27BF}
\catcode"2705=13 \begingroup\lccode`\~="2705 \lowercase{\endgroup\def~}{{\symfont\char"2713}}\symblank{"26BD}\symblank{"26A1}
\symrange{"2460}{"257F}\symactive{"223C}\symactive{"2248}\symactive{"2260}
\symhyph{"2011}\symhyph{"2E1A}\symblankrange{"1F300}{"1FAFF}\symblank{"FE0F}
% Caractères chinois : WenQuanYi Zen Hei (sous-ensemble GB2312 + pinyin), gras/italique simulés
\setCJKmainfont{WenQuanYiZenHei-subset.ttf}[Path=../../../fonts/,AutoFakeBold=3,AutoFakeSlant=0.2]
\xeCJKsetup{CJKspace=false}
% Pinyin : voyelles à caron (ǎ ǐ ǒ ǔ ǖ ǘ ǚ ǜ) absentes de la police principale -> caractères actifs
\newfontface\pyfont{WenQuanYiZenHei-subset.ttf}[Path=../../../fonts/]
\newcommand\pyactive[1]{\catcode#1=13 \begingroup\lccode`\~=#1 \lowercase{\endgroup\def~}{{\pyfont\char#1}}}
\pyactive{"01CD}\pyactive{"01CE}\pyactive{"01CF}\pyactive{"01D0}\pyactive{"01D1}\pyactive{"01D2}\pyactive{"01D3}\pyactive{"01D4}\pyactive{"01D5}\pyactive{"01D6}\pyactive{"01D7}\pyactive{"01D8}\pyactive{"01D9}\pyactive{"01DA}\pyactive{"01DB}\pyactive{"01DC}
% Maths en sans-serif (Fira Math) avec chiffres et lettres droites de Plus
% Jakarta, pour que les formules se fondent dans le texte.
\usepackage[math-style=ISO,bold-style=ISO]{unicode-math}
\setmathfont{FiraMath-Regular.otf}[Path=../../../fonts/]
\setmathfont{PlusJakartaSans-Medium.ttf}[Path=../../../fonts/,range={up/num,up/{Latin,latin},\mathup}]
\usepackage{icomma} % virgule décimale sans espace en mode math
% Caractères mathématiques Unicode absents des polices de texte (Plus Jakarta,
% Archivo Black) : rendus par leur équivalent mathématique, en texte comme en
% formule — sinon ils s'affichent en carré vide (ex. « Arithmétique dans ℤ »).
\usepackage{newunicodechar}
\newunicodechar{ℝ}{\ensuremath{\mathbb{R}}}
\newunicodechar{ℤ}{\ensuremath{\mathbb{Z}}}
\newunicodechar{ℂ}{\ensuremath{\mathbb{C}}}
\newunicodechar{ℕ}{\ensuremath{\mathbb{N}}}
\newunicodechar{ℚ}{\ensuremath{\mathbb{Q}}}
\newunicodechar{𝔻}{\ensuremath{\mathbb{D}}}
\newunicodechar{α}{\ensuremath{\alpha}}
\newunicodechar{β}{\ensuremath{\beta}}
\newunicodechar{γ}{\ensuremath{\gamma}}
\newunicodechar{δ}{\ensuremath{\delta}}
\newunicodechar{ε}{\ensuremath{\varepsilon}}
\newunicodechar{θ}{\ensuremath{\theta}}
\newunicodechar{λ}{\ensuremath{\lambda}}
\newunicodechar{μ}{\ensuremath{\mu}}
\newunicodechar{σ}{\ensuremath{\sigma}}
\newunicodechar{τ}{\ensuremath{\tau}}
\newunicodechar{φ}{\ensuremath{\varphi}}
\newunicodechar{ψ}{\ensuremath{\psi}}
\newunicodechar{ω}{\ensuremath{\omega}}
\newunicodechar{Σ}{\ensuremath{\Sigma}}
% majuscules produites par \MakeUppercase dans les titres (α-aminés -> Α)
\newunicodechar{Α}{\ensuremath{\alpha}}
\newunicodechar{Β}{\ensuremath{\beta}}
\newunicodechar{Γ}{\ensuremath{\Gamma}}
\newunicodechar{Δ}{\ensuremath{\Delta}}
\newunicodechar{Ε}{\ensuremath{\varepsilon}}
\newunicodechar{Θ}{\ensuremath{\Theta}}
\newunicodechar{Λ}{\ensuremath{\Lambda}}
\newunicodechar{Μ}{\ensuremath{\mu}}
\newunicodechar{Τ}{\ensuremath{\tau}}
\newunicodechar{Φ}{\ensuremath{\Phi}}
\newunicodechar{Ψ}{\ensuremath{\Psi}}
\newunicodechar{Ω}{\ensuremath{\Omega}}
\newunicodechar{↦}{\ensuremath{\mapsto}}
\newunicodechar{∈}{\ensuremath{\in}}
\newunicodechar{∉}{\ensuremath{\notin}}
\newunicodechar{∧}{\ensuremath{\wedge}}
\newunicodechar{⇔}{\ensuremath{\Leftrightarrow}}
\newunicodechar{⟺}{\ensuremath{\Longleftrightarrow}}
\newunicodechar{⇒}{\ensuremath{\Rightarrow}}
\newunicodechar{⟹}{\ensuremath{\Longrightarrow}}
\newunicodechar{∘}{\ensuremath{\circ}}
\newunicodechar{∪}{\ensuremath{\cup}}
\newunicodechar{∩}{\ensuremath{\cap}}
\newunicodechar{≡}{\ensuremath{\equiv}}
\newunicodechar{⊂}{\ensuremath{\subset}}
\newunicodechar{⊥}{\ensuremath{\perp}}
\newunicodechar{∃}{\ensuremath{\exists}}
\newunicodechar{∀}{\ensuremath{\forall}}
\newunicodechar{∛}{\ensuremath{\sqrt[3]{\,}}}
\newunicodechar{∜}{\ensuremath{\sqrt[4]{\,}}}
\newunicodechar{⃗}{\ensuremath{{}^{\to}}}
\newunicodechar{ⁿ}{\ifmmode^{n}\else\textsuperscript{\ensuremath{n}}\fi}
\newunicodechar{ˣ}{\ifmmode^{x}\else\textsuperscript{\ensuremath{x}}\fi}
\newunicodechar{ᵇ}{\ifmmode^{b}\else\textsuperscript{\ensuremath{b}}\fi}
\newunicodechar{ʳ}{\ifmmode^{r}\else\textsuperscript{\ensuremath{r}}\fi}
\newunicodechar{ᵘ}{\ifmmode^{u}\else\textsuperscript{\ensuremath{u}}\fi}
\newunicodechar{ʸ}{\ifmmode^{y}\else\textsuperscript{\ensuremath{y}}\fi}
\newunicodechar{ᵃ}{\ifmmode^{a}\else\textsuperscript{\ensuremath{a}}\fi}
\newunicodechar{ᵉ}{\ifmmode^{e}\else\textsuperscript{\ensuremath{e}}\fi}
\newunicodechar{ᵏ}{\ifmmode^{k}\else\textsuperscript{\ensuremath{k}}\fi}
\newunicodechar{ᵖ}{\ifmmode^{p}\else\textsuperscript{\ensuremath{p}}\fi}
\newunicodechar{ᴺ}{\ifmmode^{N}\else\textsuperscript{\ensuremath{N}}\fi}
\newunicodechar{ᶜ}{\ifmmode^{c}\else\textsuperscript{\ensuremath{c}}\fi}
\newunicodechar{ʰ}{\ifmmode^{h}\else\textsuperscript{\ensuremath{h}}\fi}
\newunicodechar{ᵠ}{\ifmmode^{q}\else\textsuperscript{\ensuremath{q}}\fi}
\newunicodechar{ₙ}{\ifmmode_{n}\else\textsubscript{\ensuremath{n}}\fi}
\newunicodechar{ₐ}{\ifmmode_{a}\else\textsubscript{\ensuremath{a}}\fi}
\newunicodechar{ᵢ}{\ifmmode_{i}\else\textsubscript{\ensuremath{i}}\fi}
\newunicodechar{ᵣ}{\ifmmode_{r}\else\textsubscript{\ensuremath{r}}\fi}
\newunicodechar{ₖ}{\ifmmode_{k}\else\textsubscript{\ensuremath{k}}\fi}
\newunicodechar{ᵦ}{\ifmmode_{\beta}\else\textsubscript{\ensuremath{\beta}}\fi}
\newunicodechar{ₓ}{\ifmmode_{x}\else\textsubscript{\ensuremath{x}}\fi}
\newfontfamily\popdisplay{ArchivoBlack-Regular.ttf}[Path=../../../fonts/]
\newfontfamily\poplogo{SpaceGrotesk-Bold.ttf}[Path=../../../fonts/]
\newfontfamily\popsemi{PlusJakartaSans-SemiBold.ttf}[Path=../../../fonts/]
\newfontfamily\popxbold{PlusJakartaSans-ExtraBold.ttf}[Path=../../../fonts/]
\newfontfamily\popxboldls{PlusJakartaSans-ExtraBold.ttf}[Path=../../../fonts/,LetterSpace=3.0]

\setlist[enumerate]{leftmargin=1.7em,itemsep=5pt,topsep=4pt}
\setlist[itemize]{leftmargin=1.7em,itemsep=4pt,topsep=4pt,label={\color{poporange}\textbullet}}
\setlength{\parindent}{0pt}
\setlength{\parskip}{5pt}
\renewcommand{\arraystretch}{1.25}

% pgfplots : style Pop par défaut
\pgfplotsset{
  every axis/.append style={axis line style={popink,line width=1pt},tick style={popink},
    grid style={popline,line width=0.5pt},label style={font=\small},tick label style={font=\footnotesize}},
  popcurve/.style={poporange,line width=1.8pt,no markers,smooth},
}
% Même style disponible dans un \draw TikZ simple (hors pgfplots)
\tikzset{popcurve/.style={poporange,line width=1.8pt}}
\tikzset{popgrid/.style={popline,very thin}}
\colorlet{popcurve}{poporange} % le nom du style est parfois utilisé comme couleur

% ============================================================
% Pill / squiggle
% ============================================================
\newcommand{\poppill}[3]{%
  \tikz[baseline=-0.55ex]\node[draw=popink,line width=1.2pt,rounded corners=2.6pt,%
    fill=#2,inner xsep=6.5pt,inner ysep=3pt]{%
    \popxboldls\fontsize{7}{7}\selectfont\color{#3}\MakeUppercase{#1}};%
}
\newcommand{\popsquiggle}[1]{%
  \tikz[baseline=-0.4ex]{
    \draw[#1,line width=2.6pt,line cap=round]
      plot [smooth] coordinates {(0,0.09) (0.34,0.01) (0.68,0.21) (1.02,0.01) (1.36,0.10)};
  }%
}
% Mot-clé surligné (marqueur orange sous le texte)
\newcommand{\cle}[1]{{\popxbold\color{popdark}#1}}

% ============================================================
% En-tête / pied
% ============================================================
\pagestyle{fancy}
\fancyhf{}
\renewcommand{\headrulewidth}{0pt}
\renewcommand{\footrulewidth}{0pt}
\lhead{\raisebox{-0.15ex}{\poplogo\fontsize{13}{13}\selectfont\color{popink}OnBuch\color{poporange}+}}
\rhead{\poppill{\DOCMATIERE\ \textbullet\ \DOCNIVEAU\ \textbullet\ Leçon \DOCLECON}{white}{popink}}
\lfoot{\popxbold\fontsize{7.6}{7.6}\selectfont\color{popmuted}\MakeUppercase{Cours signé OnBuch+}\ \textbullet\ {\popsemi\DOCTITRE}}
\rfoot{\tikz[baseline=-0.6ex]\node[rounded corners=8.5pt,fill=popink,minimum height=17pt,minimum width=17pt,inner xsep=5pt,inner sep=0]{\popxbold\fontsize{8}{8}\selectfont\color{popcream}\thepage\,/\,\pageref*{LastPage}};}

% ============================================================
% Titres
% ============================================================
\newcommand{\chaptitre}[2]{%
  \begin{tcolorbox}[enhanced,colback=poporange,colframe=popink,boxrule=2.6pt,arc=11pt,
      drop shadow=popink,left=15pt,right=15pt,top=12pt,bottom=11pt,before skip=3pt,after skip=18pt]
    \poppill{#2}{popink}{popcream}\par\vspace{9pt}
    {\raggedright\hyphenpenalty=10000\exhyphenpenalty=10000\popdisplay\fontsize{22}{24}\selectfont\color{popcream}\MakeUppercase{#1}\par}
  \end{tcolorbox}
}
% Section numérotée : pastille orange avec le numéro + titre Archivo
\newcounter{popsec}
\newcommand{\coursec}[1]{%
  \stepcounter{popsec}\setcounter{popsub}{0}%
  \par\vspace{16pt}\noindent
  \begin{minipage}{\linewidth}
  \tikz[baseline=-0.7ex]\node[draw=popink,line width=1.6pt,fill=poporange,rounded corners=4pt,minimum size=22pt,inner sep=0]{\popdisplay\fontsize{12}{12}\selectfont\color{popcream}\thepopsec};%
  \hspace{8pt}{\popdisplay\fontsize{15}{17}\selectfont\color{popink}\MakeUppercase{#1}}\par
  \vspace{4pt}\noindent{\color{poporange}\rule{36pt}{3.6pt}}
  \end{minipage}\par\nopagebreak\vspace{8pt}
}
\newcounter{popsub}
\newcommand{\courssub}[1]{%
  \stepcounter{popsub}%
  \par\vspace{9pt}\noindent{\popxbold\fontsize{11.5}{13}\selectfont\color{popink}{\color{poporange}\thepopsec.\thepopsub}\hspace{6pt}#1}\par\nopagebreak\vspace{4pt}
}

% ============================================================
% Boîtes pédagogiques — même recette : fond blanc, bordure épaisse colorée,
% ombre dure encre, badge plein. Titre optionnel [#1].
% ============================================================
\tcbset{popbase/.style={breakable,enhanced,colback=white,boxrule=2.3pt,arc=9pt,
  shadow={3pt}{-3pt}{0pt}{popink},left=14pt,right=14pt,top=11pt,bottom=11pt,before skip=11pt,after skip=15pt,
  fontupper=\popsemi\fontsize{10.6}{14.5}\selectfont\color{popink},
  fontlower=\popsemi\fontsize{10.6}{14.5}\selectfont\color{popink}}}
% Coche dessinée (le glyphe ✓ n'existe pas dans les polices du document)
\newcommand{\popcheck}[1]{\tikz[baseline=-0.5ex]\draw[#1,line width=1.6pt,line cap=round,line join=round](-0.13,0.0)--(-0.04,-0.09)--(0.13,0.1);}
\providecommand{\coche}{\popcheck{popgreen}}
\providecommand{\resultat}[1]{\par\smallskip\noindent\fcolorbox{popgreen}{popgreenL}{\parbox{\dimexpr\linewidth-2\fboxsep-2\fboxrule\relax}{\textbf{Résultat.} #1}}\par\smallskip}
\newcommand{\popboxhead}[3]{\poppill{#1}{#2}{white}\ifx\relax#3\relax\else\hspace{6pt}{\popxbold\fontsize{10.6}{12}\selectfont\color{popink}#3}\fi\par\vspace{7pt}}

\newtcolorbox{aretenir}[1][]{popbase,colframe=popgreen,before upper={\popboxhead{À retenir}{popgreen}{#1}}}
% Texte en chinois (document de lecture, dialogue, extrait) : cadre
% d'étude. Un titre optionnel (source, auteur) s'affiche dans l'en-tête.
\newtcolorbox{textebox}[1][]{popbase,colframe=popblue,before upper={\popboxhead{文本 · Texte}{popblue}{#1}},after upper={}}
% Marques ✓ / ✗ (forme correcte / fautive) souvent utilisées par les modèles
\providecommand{\cross}{\textcolor{poppink}{\ensuremath{\times}}}
\providecommand{\xmark}{\cross}\providecommand{\crossmark}{\cross}\providecommand{\cmark}{\ensuremath{\checkmark}}
% Ligne de réponse / justification à compléter (inventée par les modèles)
\providecommand{\justification}[1]{\par\noindent\textit{Justification~:} \dotfill\par}
% \textipa (package tipa non chargé) : la police ne couvre pas l'API -> simple passage
\providecommand{\textipa}[1]{#1}
% Soulignement (ulem non chargé) : \uline, \ul
\providecommand{\uline}[1]{\underline{#1}}\providecommand{\ul}[1]{\underline{#1}}
% \glossaire{\item ...} inventée par les modèles : liste de glossaire
\providecommand{\glossaire}[1]{\begin{itemize}#1\end{itemize}}
% \alert (beamer) inventée par les modèles : texte mis en évidence
\providecommand{\alert}[1]{\textbf{\textcolor{poppink}{#1}}}
% variantes ASCII d'ordinaux écrites en macro
\providecommand{\iere}{\textsuperscript{re}}\providecommand{\ieme}{\textsuperscript{e}}\providecommand{\ere}{\textsuperscript{re}}\providecommand{\eme}{\textsuperscript{e}}\providecommand{\er}{\textsuperscript{er}}\providecommand{\ie}{\textsuperscript{e}}
% Ordinaux français écrits en macro par les modèles : 1\ière, 3\ième
\newcommand{\ière}{\textsuperscript{re}}\newcommand{\ième}{\textsuperscript{e}}
% \exemple (inventée par les modèles) : étiquette « Exemple : »
\providecommand{\exemple}{\textbf{Exemple~:}\ }
% \square utilisable aussi hors mode math (cases à cocher)
\AtBeginDocument{\let\squareMath\square\renewcommand{\square}{\ensuremath{\squareMath}}}
% Mot ou phrase en chinois à l'intérieur d'un paragraphe français
\newcommand{\zh}[1]{\ifmmode\text{#1}\else{#1}\fi}
\let\eng\zh \let\esp\zh \let\all\zh % alias tolérés
% flèches d'intonation inventées par les modèles
\providecommand{\textsearrow}{}\renewcommand{\textsearrow}{\ensuremath{\searrow}}\providecommand{\textnearrow}{}\renewcommand{\textnearrow}{\ensuremath{\nearrow}}
% \fr{..} : mot français (inventé par les modèles)
\providecommand{\fr}[1]{\textit{#1}}
% zhbox : encadré de texte chinois inventé par les modèles
\newenvironment{zhbox}{\begin{quote}}{\end{quote}}
% \vs : « versus » inventé par les modèles
\providecommand{\vs}{\textit{vs}\ }
% \blank : case à compléter inventée par les modèles
\providecommand{\blank}[1][]{\underline{\hspace{1.6em}}}
\newtcolorbox{exemplebox}[1][]{popbase,colframe=poporange,before upper={\popboxhead{Exemple}{poporange}{#1}}}
\newtcolorbox{definition}[1][]{popbase,colframe=popblue,before upper={\popboxhead{Définition}{popblue}{#1}}}
\newtcolorbox{propriete}[1][]{popbase,colframe=poppurple,before upper={\popboxhead{Propriété}{poppurple}{#1}}}
\newtcolorbox{methode}[1][]{popbase,colframe=popgold,before upper={\popboxhead{Méthode}{popgold}{#1}}}
\newtcolorbox{attention}[1][]{popbase,colframe=poppink,before upper={\popboxhead{Attention}{poppink}{#1}}}
\newtcolorbox{experience}[1][]{popbase,colframe=popblue,colback=popblueL,before upper={\popboxhead{Expérience}{popblue}{#1}}}
% « Le savais-tu ? » : carte violette pleine (contexte, culture, Cameroun)
\newtcolorbox{savaistu}[1][]{popbase,colback=poppurple,colframe=popink,
  fontupper=\popsemi\fontsize{10.4}{14.2}\selectfont\color{white},
  before upper={\poppill{Le savais-tu\,?}{white}{poppurple}\ifx\relax#1\relax\else\hspace{6pt}{\popxbold\color{white}#1}\fi\par\vspace{7pt}}}
% Exercice résolu : énoncé en haut, \tcblower puis la solution sur fond crème
\newcounter{popexo}\newcounter{popexr}
\newtcolorbox{exoresolu}[1][]{popbase,colframe=poporange,colbacklower=popcream,
  segmentation style={popink,line width=1.2pt,dashed},
  before upper={\stepcounter{popexr}\popboxhead{Exercice résolu \thepopexr}{poporange}{#1}},
  before lower={\poppill{Solution}{popink}{popcream}\par\vspace{6pt}}}
% Exercice (non résolu en place) — la correction est dans la partie Corrigés
\newtcolorbox{exercice}[1][]{popbase,colframe=popink,
  before upper={\stepcounter{popexo}\popboxhead{Exercice \thepopexo}{popink}{#1}}}
% Corrigé
\newtcolorbox{corrige}[1][]{popbase,colframe=popgreen,colback=popgreenL,
  before upper={\popboxhead{Corrigé}{popgreen}{#1}}}
% Objectifs / prérequis / bilan
\newtcolorbox{objectifs}{popbase,colframe=popink,colback=poporangeL,
  before upper={\popboxhead{Objectifs de la leçon}{popink}{}}}
\newtcolorbox{prerequis}{popbase,colframe=popink,colback=popblueL,
  before upper={\popboxhead{Prérequis}{popblue}{}}}
\newtcolorbox{bilan}{popbase,colframe=popink,colback=white,boxrule=2.8pt,
  before upper={\popboxhead{Fiche bilan}{poporange}{L'essentiel à connaître pour l'examen}}}
% Figure encadrée avec légende
\newtcolorbox{popfigure}[1][]{enhanced,breakable=false,colback=white,colframe=popline,boxrule=1.4pt,arc=8pt,
  left=10pt,right=10pt,top=10pt,bottom=8pt,before skip=10pt,after skip=12pt,halign=center,
  lower separated=false,halign lower=center,
  fontlower=\popsemi\fontsize{9}{11.5}\selectfont\color{popmuted},
  #1}
\newcommand{\legende}[1]{\par\vspace{6pt}{\popsemi\fontsize{9}{11.5}\selectfont\color{popmuted}#1}}

% ============================================================
% Signature OnBuch+
% ============================================================
\newcommand{\poplogoinline}[1]{{\poplogo\fontsize{#1}{#1}\selectfont\color{popink}OnBuch\color{poporange}+}}
\newcommand{\signatureonbuch}{%
  \begin{tcolorbox}[enhanced,colback=popink,colframe=popink,boxrule=2.2pt,arc=10pt,
      drop shadow=poporange,left=14pt,right=14pt,top=10pt,bottom=10pt,before skip=0pt,after skip=0pt]
    \tikz[baseline=-0.7ex]\node[circle,fill=popgreen,draw=popcream,line width=1.3pt,minimum size=22pt,inner sep=0]{\popcheck{white}};%
    \hspace{9pt}\begin{minipage}[c]{0.62\linewidth}
      {\popxbold\fontsize{12}{13}\selectfont\color{popcream}Signé\ \textbullet\ {\poplogo OnBuch\color{poporange}+}}\par\vspace{2pt}
      {\raggedright\popsemi\fontsize{8}{10.5}\selectfont\color{popline}\MakeUppercase{Équipe pédagogique OnBuch+}\\\MakeUppercase{Conforme au programme officiel MINESEC}\par}
    \end{minipage}\hfill
    {\poplogo\fontsize{22}{22}\selectfont\color{popcream}OnBuch\color{poporange}+}
  \end{tcolorbox}%
}

% ============================================================
% Page de garde
% #1 titre · #2 matière · #3 niveau · #4 module · #5 n° leçon · #6 durée ·
% #7 description / objectifs (texte ou itemize)
% ============================================================
\newcommand{\pagegarde}[7]{%
  \thispagestyle{empty}
  \newgeometry{a4paper,margin=1.7cm,top=1.6cm,bottom=1.4cm}%
  \begin{tikzpicture}[remember picture,overlay]
    \fill[poporange!18] ([xshift=-3.2cm,yshift=-3.6cm]current page.north east) circle (4.6cm);
    \fill[poppurple!14] ([xshift=2.4cm,yshift=7.6cm]current page.south west) circle (3.8cm);
  \end{tikzpicture}%
  {\poplogo\fontsize{30}{30}\selectfont\color{popink}OnBuch\color{poporange}+}\hfill
  \raisebox{0.6ex}{\poppill{Cours signé \textbullet\ Premium}{popink}{popcream}}\par
  \vspace{1pt}\popsquiggle{poppurple}\par
  \vspace{8pt}
  \begin{tcolorbox}[enhanced,colback=poporange,colframe=popink,boxrule=2.8pt,arc=13pt,
      drop shadow=popink,left=18pt,right=18pt,top=12pt,bottom=13pt,after skip=10pt]
    \poppill{#2 \textbullet\ #3}{popink}{popcream}\hspace{5pt}\poppill{#4}{white}{popink}\par\vspace{8pt}
    {\popdisplay\fontsize{14}{14}\selectfont\color{popink}LEÇON #5}\par\vspace{4pt}
    {\raggedright\hyphenpenalty=10000\exhyphenpenalty=10000\popdisplay\fontsize{25}{27}\selectfont\color{popcream}\MakeUppercase{#1}\par}
    \vspace{8pt}
    {\popxbold\fontsize{9.5}{9.5}\selectfont\color{popink}\MakeUppercase{Durée conseillée : #6}}
  \end{tcolorbox}
  \begin{tcolorbox}[enhanced,colback=white,colframe=popink,boxrule=2.2pt,arc=10pt,
      drop shadow=popink,left=16pt,right=16pt,top=10pt,bottom=10pt,after skip=8pt]
    \poppill{Dans cette leçon}{poppurple}{white}\par\vspace{5pt}
    \popsemi\fontsize{9.3}{12.6}\selectfont\color{popink}#7
  \end{tcolorbox}
  \noindent\begin{minipage}[t]{0.487\linewidth}
    \begin{tcolorbox}[enhanced,colback=popgreen,colframe=popink,boxrule=2.2pt,arc=10pt,
        drop shadow=popink,left=11pt,right=11pt,top=10pt,bottom=10pt,height=3.1cm,valign=center]
      \tcbox[colback=white,colframe=popink,boxrule=1.3pt,arc=5pt,boxsep=0pt,nobeforeafter,
        left=3pt,right=3pt,top=3pt,bottom=3pt,valign=center]{\qrcode[height=1.9cm]{https://play.google.com/store/apps/details?id=com.luvvix.onbuch&hl=fr}}%
      \hspace{8pt}\begin{minipage}[c]{0.5\linewidth}
        {\popdisplay\fontsize{11}{12.5}\selectfont\color{white}\MakeUppercase{Télécharge OnBuch}}\par\vspace{4pt}
        {\raggedright\popsemi\fontsize{8.4}{11}\selectfont\color{white}Gratuit \textbullet\ Google Play\\Tuteur IA, annales, concours, résultats\par}
      \end{minipage}
    \end{tcolorbox}
  \end{minipage}\hfill
  \begin{minipage}[t]{0.487\linewidth}
    \begin{tcolorbox}[enhanced,colback=poppurple,colframe=popink,boxrule=2.2pt,arc=10pt,
        drop shadow=popink,left=11pt,right=11pt,top=10pt,bottom=10pt,height=3.1cm,valign=center]
      \tikz[baseline=-0.7ex]\node[circle,fill=white,draw=popink,line width=1.3pt,minimum size=1.6cm,inner sep=0]{\popdisplay\fontsize{24}{24}\selectfont\color{poppurple}+};%
      \hspace{8pt}\begin{minipage}[c]{0.6\linewidth}
        {\popdisplay\fontsize{11}{12.5}\selectfont\color{white}\MakeUppercase{Passe à OnBuch+}}\par\vspace{4pt}
        {\raggedright\popsemi\fontsize{8.4}{11}\selectfont\color{white}Tous les cours signés, Tuteur IA illimité, fiches PDF — onglet OnBuch+ de l'app\par}
      \end{minipage}
    \end{tcolorbox}
  \end{minipage}
  \vspace{8pt}
  \popcontenu
  \vfill
  \signatureonbuch
  \restoregeometry
  \newpage
}

% Rangée « Ce cours contient » (page de garde)
\newcommand{\poptile}[3]{% #1 couleur #2 gros texte #3 légende
  \begin{tcolorbox}[enhanced,colback=#1,colframe=popink,boxrule=2pt,arc=9pt,shadow={2.5pt}{-2.5pt}{0pt}{popink},
      left=6pt,right=6pt,top=9pt,bottom=9pt,halign=center,valign=center,height=1.75cm,width=\linewidth,nobeforeafter]
    {\popdisplay\fontsize{12.5}{13}\selectfont\color{white}\MakeUppercase{#2}}\par\vspace{3pt}
    {\centering\popsemi\fontsize{7.8}{9.5}\selectfont\color{white}#3\par}
  \end{tcolorbox}}
\newcommand{\popcontenu}{%
  \par\noindent{\popxboldls\fontsize{7.5}{7.5}\selectfont\color{popmuted}CE COURS SIGNÉ CONTIENT}\par\vspace{7pt}
  \noindent\begin{minipage}[t]{0.235\linewidth}\poptile{popblue}{Cours}{complet, illustré, pas à pas}\end{minipage}\hfill
  \begin{minipage}[t]{0.235\linewidth}\poptile{poporange}{Méthodes}{et exercices résolus}\end{minipage}\hfill
  \begin{minipage}[t]{0.235\linewidth}\poptile{poppink}{Exercices}{type examen + corrigés}\end{minipage}\hfill
  \begin{minipage}[t]{0.235\linewidth}\poptile{popgreen}{Fiche bilan}{l'essentiel à retenir}\end{minipage}\par}

% Sommaire
\newtcolorbox{sommaire}{enhanced,colback=white,colframe=popink,boxrule=2.2pt,arc=10pt,
  drop shadow=popink,left=18pt,right=18pt,top=13pt,bottom=13pt,before skip=6pt,after skip=20pt,
  before upper={\poppill{Sommaire}{poppurple}{white}\par\vspace{9pt}\popsemi\fontsize{10.6}{14.5}\selectfont\color{popink}}}

% Signature de fin de cours — poussée tout en bas de la dernière page
\newcommand{\findecours}{%
  \par\vfill\nopagebreak
  \begin{center}\popsquiggle{poporange}\end{center}\vspace{4pt}
  \signatureonbuch
}
\AtBeginDocument{\def\llbracket{\mathopen{[\mkern-3.5mu[}}\def\rrbracket{\mathclose{]\mkern-3.5mu]}}} % intervalles d'entiers (glyphe ⟦ absent de la police)
% Glyphes absents de Fira Math : redéfinis à partir de glyphes disponibles
\AtBeginDocument{%
  \def\setminus{\mathbin{\backslash}}\let\smallsetminus\setminus
  \def\vdots{\mathord{\vbox{\baselineskip3.2pt\lineskiplimit0pt\kern4pt\hbox{.}\hbox{.}\hbox{.}}}}%
  \def\ddots{\mathinner{\mkern1mu\raise6.5pt\hbox{.}\mkern2mu\raise3.6pt\hbox{.}\mkern2mu\raise0.7pt\hbox{.}\mkern1mu}}%
  \def\triangle{\mathord{\scalebox{0.9}{$\symup{\Delta}$}}}%
  \def\nabla{\mathord{\raisebox{\depth}{\scalebox{1}[-1]{$\symup{\Delta}$}}}}%
  \def\checkmark{\popcheck{popdark}}%
  \def\star{\mathbin{\ast}}\def\bigstar{\mathord{\ast}}%
  \def\diamond{\mathbin{\scalebox{0.6}{$\Diamond$}}}%
  \def\bigcup{\mathop{\mathchoice{\raisebox{-0.3em}{\scalebox{1.7}{$\cup$}}}{\scalebox{1.25}{$\cup$}}{\cup}{\cup}}\displaylimits}%
  \def\bigcap{\mathop{\mathchoice{\raisebox{-0.3em}{\scalebox{1.7}{$\cap$}}}{\scalebox{1.25}{$\cap$}}{\cap}{\cap}}\displaylimits}%
  \def\lesseqqgtr{\mathrel{\lessgtr}}\def\bigoplus{\mathop{\oplus}\displaylimits}%
}
\providecommand{\ssi}{si et seulement si} % abréviation fréquente
\AtBeginDocument{\providecommand{\wideparen}{\overparen}} % arc de cercle

%% FIGFIX-BEGIN v5 — correctifs globaux des figures (docs/onbuch-generation/tools/figfix)
\usepackage{adjustbox}\usepackage{etoolbox}\tcbuselibrary{raster}
% toute figure TikZ/pgfplots plus large que la ligne est réduite à la largeur disponible
\BeforeBeginEnvironment{tikzpicture}{\begin{adjustbox}{max width=\linewidth}}
\AfterEndEnvironment{tikzpicture}{\end{adjustbox}}
% légendes pgfplots lisibles par-dessus les courbes
\pgfplotsset{every axis/.append style={legend style={fill=white,fill opacity=0.92,text opacity=1,draw=black!15,inner sep=3pt}}}
% étiquettes d'arêtes (syntaxe edge["..."]) avec fond blanc
\tikzset{every edge quotes/.append style={fill=white,inner sep=1.2pt}}
%% FIGFIX-END
\begin{document}

\pagegarde{Leçon 15 — Expression orale : recyclage et paludisme}{Chinois}{Tle A}{Module 2 : Environnement et santé}{ZH15}{2 h}{Cette leçon d'expression orale entraîne les élèves à présenter un exposé simple et à participer à un débat guidé en chinois sur deux thèmes de santé et d'environnement : le recyclage et le paludisme. Elle fournit les formules fonctionnelles, les dialogues modèles et le travail des tons nécessaires à une prise de parole claire et correcte.
\vspace{4pt}
\begin{multicols}{2}\small
\begin{itemize}
\item À la fin de cette leçon, je sais présenter oralement un exposé simple et structuré en chinois sur le recyclage ou le paludisme
\item je sais utiliser les formules fonctionnelles pour ouvrir un exposé, donner mon avis et conclure
\item je sais relancer la parole et poser des questions à mes camarades pendant un débat
\item je sais exprimer l'accord et le désaccord poliment en chinois
\item je sais prononcer correctement les tons des mots clés du thème santé-environnement
\item je sais participer à un jeu de rôle (journaliste, expert, élève) sur ces thèmes
\end{itemize}\end{multicols}}

\begin{sommaire}
\begin{multicols}{2}
\begin{enumerate}
\item Mise en route et lexique oral du thème
\item Dialogue modèle 1 : exposé sur le recyclage
\item Dialogue modèle 2 : débat sur le paludisme
\item Formules fonctionnelles de l'exposé et du débat
\item Travail des tons et de la prononciation
\item Jeux de rôle et débat guidé
\item Activité d'intégration
\item Méthodes et exercices résolus
\item Exercices
\item Corrigés des exercices
\item Fiche bilan
\end{enumerate}
\end{multicols}
\end{sommaire}

\begin{objectifs}
\begin{itemize}
\item À la fin de cette leçon, je sais présenter oralement un exposé simple et structuré en chinois sur le recyclage ou le paludisme
\item je sais utiliser les formules fonctionnelles pour ouvrir un exposé, donner mon avis et conclure
\item je sais relancer la parole et poser des questions à mes camarades pendant un débat
\item je sais exprimer l'accord et le désaccord poliment en chinois
\item je sais prononcer correctement les tons des mots clés du thème santé-environnement
\item je sais participer à un jeu de rôle (journaliste, expert, élève) sur ces thèmes
\item je sais réagir spontanément avec des connecteurs simples (因为, 所以, 但是, 而且)
\item je sais évaluer la prise de parole d'un camarade à l'aide d'une grille simple
\end{itemize}
\end{objectifs}

\begin{prerequis}
\begin{itemize}
\item Vocabulaire de l'environnement et de la santé vu en Première (环境, 健康, 生病)
\item Structures de base : 因为...所以..., 我觉得..., 应该
\item Nombres, dates et pourcentages simples pour citer des données
\item Maîtrise des quatre tons et du ton neutre
\item Présentation personnelle simple (我叫..., 我是...)
\end{itemize}
\end{prerequis}

\newpage

\coursec{Mise en route et lexique oral du thème}

\courssub{Situation d'accroche et problématique}

À Yaoundé comme à Douala, les bouteilles en plastique s'accumulent dans les rues après les pluies, et l'eau stagnante qu'elles retiennent devient un gîte idéal pour les moustiques. Au marché Mokolo, une jeune élève propose à ses camarades de recycler ces bouteilles pour lutter à la fois contre la pollution et contre le paludisme. Comment présenter cette idée en chinois devant la classe ? Comment donner son avis, relancer la parole et débattre ? C'est tout l'enjeu de cette leçon d'expression orale.

\begin{textebox}[Situation de communication : Au marché Mokolo]
我们在莫科洛市场看到很多塑料瓶子。雨后，瓶子里有水，蚊子喜欢在水里产卵。我们应该回收瓶子，保护环境，预防疟疾。
Wǒmen zài Mòkēluó shìchǎng kàndào hěnduō sùliào píngzi. Yǔhòu, píngzi lǐ yǒu shuǐ, wénzi xǐhuan zài shuǐ lǐ chǎnluǎn. Wǒmen yīnggāi huíshōu píngzi, bǎohù huánjìng, yùfáng nüèjí.
Nous voyons beaucoup de bouteilles en plastique au marché Mokolo. Après la pluie, il y a de l'eau dans les bouteilles, les moustiques aiment pondre leurs œufs dans l'eau. Nous devrions recycler les bouteilles, protéger l'environnement, prévenir le paludisme.
\end{textebox}

\courssub{Brainstorming en classe : déchets et moustiques}

Avant d'aborder le lexique précis, activons vos connaissances. Observez votre environnement immédiat : la cour du lycée, les abords de la cantine, les bords de route.

\begin{experience}[Brainstorming guidé : « Autour de moi »]
\textbf{Consigne :} En groupe de 4, dressez deux listes en 3 minutes.
\begin{enumerate}
    \item \textbf{Déchets visibles :} Quels types de déchets (\zh{垃圾} lājī) voyez-vous le plus souvent ? (Ex: bouteilles, sachets, canettes, papiers…)
    \item \textbf{Gîtes larvaires :} Où l'eau stagnante peut-elle s'accumuler pour devenir un nid à moustiques (\zh{蚊子} wénzi) ? (Ex: pneus usagés, bocaux, gouttières bouchées, soucoupes de pots de fleurs…)
\end{enumerate}
\textbf{Mise en commun :} Un secrétaire par groupe note au tableau sous deux colonnes : \zh{垃圾} et \zh{蚊子的家}. Le professeur valide et complète en chinois.
\end{experience}

\courssub{Lexique thématique : Recyclage (回收)}

Observons le texte de la situation d'accroche et les mots nouveaux. La structure \textbf{Sujet + 应该 + Verbe + Objet} (devrait) est essentielle pour exprimer une obligation morale ou un conseil.

\begin{definition}[Vocabulaire de base : Le recyclage]
\begin{center}
\begin{tabularx}{\linewidth}{|X|X|X|X|}
\hline
\rowcolor{popblueL}
\textbf{Caractères} & \textbf{Pinyin} & \textbf{Nature} & \textbf{Traduction} \\
\hline
回收 & huíshōu & verbe & recycler, récupérer \\
\hline
垃圾 & lājī & nom & déchets, ordures \\
\hline
瓶子 & píngzi & nom & bouteille \\
\hline
塑料 & sùliào & nom & plastique \\
\hline
环境 & huánjìng & nom & environnement \\
\hline
污染 & wūrǎn & verbe/nom & polluer / pollution \\
\hline
保护 & bǎohù & verbe & protéger \\
\hline
\end{tabularx}
\end{center}
\end{definition}

\begin{propriete}[Structure clé : Exprimer un devoir / un conseil avec \cle{应该}]
\textbf{Forme :} \zh{主语 + 应该 + 动词 (+ 对象)} \\
\textbf{Emploi :} Exprime une obligation morale, un conseil fort, une recommandation. Plus fort que \zh{可以} (pouvoir), moins contraignant que \zh{必须} (devoir absolu). \\
\textbf{Négation :} \zh{不应该} (ne devrait pas). \\
\textbf{Exemples :}
\begin{center}
\begin{tabularx}{\linewidth}{|X|X|X|}
\hline
\rowcolor{popgreenL}
\textbf{Chinois (Caractères + Pinyin)} & \textbf{Français} & \textbf{Contexte} \\
\hline
我们 \textbf{应该} 回收 瓶子。 Wǒmen \textbf{yīnggāi} huíshōu píngzi. & Nous \textbf{devrions} recycler les bouteilles. & Conseil collectif \\
\hline
学生 \textbf{不应该} 乱扔 垃圾。 Xuésheng \textbf{bù yīnggāi} luàn rēng lājī. & Les élèves \textbf{ne devraient pas} jeter les déchets n'importe où. & Règle de vie \\
\hline
为了 环境，大 家 \textbf{应该} 少用 塑料 袋。 Wèile huánjìng, dàjiā \textbf{yīnggāi} shǎo yòng sùliào dài. & Pour l'environnement, tout le monde \textbf{devrait} moins utiliser de sacs en plastique. & Argumentation \\
\hline
\end{tabularx}
\end{center}
\end{propriete}

\begin{exemplebox}[Phrases modèles : Recyclage et environnement]
\begin{itemize}
    \item \zh{垃圾 会 污染 环境。} Lājī huì wūrǎn huánjìng. Les déchets polluent l'environnement. (Prédiction/constat avec \cle{会})
    \item \zh{塑料 瓶子 可以 回收。} Sùliào píngzi kěyǐ huíshōu. Les bouteilles en plastique peuvent être recyclées. (Possibilité avec \cle{可以})
    \item \zh{保护 环境 是 我们 的 责任。} Bǎohù huánjìng shì wǒmen de zérèn. Protéger l'environnement est notre responsabilité. (Structure \cle{是...的} pour insister sur le fait/identité)
\end{itemize}
\end{exemplebox}

\courssub{Lexique thématique : Paludisme (疟疾)}

Le paludisme est un fléau majeur au Cameroun. Le vocabulaire médical de base permet de décrire les symptômes, le lieu de soin et la prévention.

\begin{definition}[Vocabulaire de base : Le paludisme]
\begin{center}
\begin{tabularx}{\linewidth}{|X|X|X|X|}
\hline
\rowcolor{poporangeL}
\textbf{Caractères} & \textbf{Pinyin} & \textbf{Nature} & \textbf{Traduction} \\
\hline
疟疾 & nüèjí & nom & paludisme (malaria) \\
\hline
蚊子 & wénzi & nom & moustique \\
\hline
发烧 & fāshāo & verbe & avoir de la fièvre \\
\hline
医院 & yīyuàn & nom & hôpital \\
\hline
预防 & yùfáng & verbe & prévenir, prévention \\
\hline
蚊帐 & wénzhàng & nom & moustiquaire \\
\hline
症状 & zhèngzhuàng & nom & symptôme \\
\hline
药 & yào & nom & médicament \\
\hline
\end{tabularx}
\end{center}
\end{definition}

\begin{attention}[Piège de prononciation : \cle{蚊子} vs \cle{文字}]
Ne confondez pas ces deux mots très fréquents, seul le ton du second caractère change :
\begin{itemize}
    \item \zh{蚊子} \textbf{wén\emph{zi}} (ton 2 + ton neutre) = \textbf{moustique}. Clé \zh{虫} (insecte).
    \item \zh{文字} \textbf{wén\emph{zì}} (ton 2 + ton 4) = \textbf{écriture, caractère, texte}. Clé \zh{文} (culture/écriture).
\end{itemize}
\textbf{Astuce mnémotechnique :} Le moustique (\zh{子} = enfant/bébé) pique et fait « zi » (ton neutre, court, comme une piqûre rapide). L'écriture (\zh{字} = caractère) est structurée, le ton 4 tombe net comme un trait de pinceau final.
\end{attention}

\begin{propriete}[Structure clé : Exprimer la nécessité / l'obligation pratique avec \cle{要}]
\textbf{Forme :} \zh{主语 + 要 + 动词 (+ 对象)} \\
\textbf{Emploi :} « Il faut / devoir » (nécessité objective, plan, contrainte pratique). Différent de \zh{应该} (devoir moral). \\
\textbf{Exemples :}
\begin{center}
\begin{tabularx}{\linewidth}{|X|X|}
\hline
\rowcolor{poppurpleL}
\textbf{Chinois + Pinyin} & \textbf{Français} \\
\hline
睡觉 \textbf{要} 用 蚊帐。 Shuìjiào \textbf{yào} yòng wénzhàng. & Pour dormir, \textbf{il faut} utiliser une moustiquaire. \\
\hline
发烧 了 \textbf{要} 去 医院。 Fāshāo le \textbf{yào} qù yīyuàn. & Si on a de la fièvre, \textbf{il faut} aller à l'hôpital. \\
\hline
预防 疟疾 \textbf{要} 清理 积水。 Yùfáng nüèjí \textbf{yào} qīnglǐ jīshuǐ. & Pour prévenir le paludisme, \textbf{il faut} éliminer l'eau stagnante. \\
\hline
\end{tabularx}
\end{center}
\end{propriete}

\begin{exemplebox}[Phrases modèles : Paludisme et prévention]
\begin{itemize}
    \item \zh{蚊子 传播 疟疾。} Wénzi chuánbō nüèjí. Les moustiques transmettent le paludisme.
    \item \zh{他 发烧 了，去 医院 看 病。} Tā fāshāo le, qù yīyuàn kàn bìng. Il a de la fièvre, il est allé à l'hôpital se faire soigner. (\cle{了} marque le changement d'état / aspect accompli).
    \item \zh{预防 胜于 治疗。} Yùfáng shèngyú zhìliáo. La prévention vaut mieux que le traitement. (Proverbe / Chengyu simplifié).
\end{itemize}
\end{exemplebox}

\begin{aretenir}[Bilan lexical de la section 1 — Mots-outils à mémoriser pour l'oral]
\begin{center}
\begin{tabularx}{\linewidth}{|X|X|X|X|}
\hline
\rowcolor{popgoldL}
\textbf{Thème} & \textbf{Mot chinois} & \textbf{Pinyin} & \textbf{Fonction orale} \\
\hline
Recyclage & 回收 & huíshōu & Verbe d'action principal \\
Recyclage & 垃圾 & lājī & Sujet/Objet principal \\
Recyclage & 保护 & bǎohù & Verbe de finalité \\
Recyclage & 环境 & huánjìng & Objet de \zh{保护} \\
Paludisme & 疟疾 & nüèjí & Sujet de santé publique \\
Paludisme & 蚊子 & wénzi & Agent causal (attention ton !) \\
Paludisme & 预防 & yùfáng & Verbe de stratégie \\
Paludisme & 蚊帐 & wénzhàng & Outil de prévention \\
Grammaire & 应该 & yīnggāi & Conseil / Devoir moral \\
Grammaire & 要 & yào & Nécessité / Obligation pratique \\
\hline
\end{tabularx}
\end{center}
\end{aretenir}

\coursec{Dialogue modèle 1 : Exposé sur le recyclage}

\courssub{Présentation orale structurée : « Mon projet pour le lycée »}

Ce dialogue modèle est un \textbf{monologue exposé} (独白 dúbái) typique d'un élève représentant sa classe. Il suit la structure classique : Salutation $\rightarrow$ Annonce du sujet $\rightarrow$ Arguments (Premièrement/Deuxièmement/Enfin) $\rightarrow$ Proposition concrète $\rightarrow$ Conclusion/Remerciements.

\begin{textebox}[Dialogue modèle 1 : Exposé sur le recyclage au lycée]
\textbf{Élève (Exposé) :}
大家 好！今天 我 想 谈谈 我们 学校 的 垃圾 问题。\\
Dàjiā hǎo! Jīntiān wǒ xiǎng tántan wǒmen xuéxiào de lājī wèntí.\\
Bonjour à tous ! Aujourd'hui, je veux parler du problème des déchets dans notre école.

首先，校园 里 有 很多 塑料 瓶子 和 塑料 袋。\\
Shǒuxiān, xiàoyuán lǐ yǒu hěnduō sùliào píngzi hé sùliào dài.\\
Premièrement, il y a beaucoup de bouteilles et de sacs en plastique dans l'enceinte de l'école.

其次，雨后 这些 垃圾 积水，容易 养 蚊子，传播 疟疾。\\
Qícì, yǔhòu zhèxiē lājī jīshuǐ, róngyì yǎng wénzi, chuánbō nüèjí.\\
Deuxièmement, après la pluie, ces déchets retiennent l'eau, favorisent les moustiques et transmettent le paludisme.

所以，我 建议 我们 设立 回收 站，分类 投放 垃圾。\\
Suǒyǐ, wǒ jiànyì wǒmen shèlì huíshōu zhàn, fēnlèi tóufàng lājī.\\
C'est pourquoi, je propose qu'on installe des points de tri et qu'on jette les déchets par catégorie.

最后，保护 环境 是 我们 共同 的 责任。谢谢 大家！\\
Zuìhòu, bǎohù huánjìng shì wǒmen gòngtóng de zérèn. Xièxie dàjiā!\\
Enfin, protéger l'environnement est notre responsabilité commune. Merci à tous !
\end{textebox}

\begin{exemplebox}[Analyse de la structure de l'exposé (Modèle réutilisable)]
\begin{center}
\begin{tabularx}{\linewidth}{|X|X|X|}
\hline
\rowcolor{popblueL}
\textbf{Étape} & \textbf{Formule chinoise (Pinyin)} & \textbf{Fonction} \\
\hline
Ouverture & \zh{大家好！今天我想谈谈……} (Dàjiā hǎo! Jīntiān wǒ xiǎng tántan…) & Saluer + Annoncer le sujet \\
\hline
Argument 1 & \zh{首先，……} (Shǒuxiān…) & Premier argument (faits) \\
\hline
Argument 2 & \zh{其次，……} (Qícì…) & Deuxième argument (conséquences) \\
\hline
Proposition & \zh{所以，我建议……} (Suǒyǐ, wǒ jiànyì…) & Conclusion logique + Proposition (\cle{建议}) \\
\hline
Conclusion & \zh{最后，…… 谢谢大家！} (Zuìhòu… Xièxie dàjiā!) & Valeur morale + Remerciements \\
\hline
\end{tabularx}
\end{center}
\end{exemplebox}

\coursec{Dialogue modèle 2 : Débat sur le paludisme}

\courssub{Échange interactif : « Quelle est la meilleure prévention ? »}

Ce dialogue modèle met en scène un \textbf{échange à deux} (对话 duìhuà) où les interlocuteurs expriment leur avis, s'accordent, nuancent et relancent. C'est le cœur de l'évaluation orale : l'interaction.

\begin{textebox}[Dialogue modèle 2 : Débat prévention paludisme]
\textbf{A (Élève 1) :} 你 觉得 预防 疟疾，什么 办法 最 好？\\
Nǐ juéde yùfáng nüèjí, shénme bànfǎ zuì hǎo?\\
Selon toi, quelle est la meilleure méthode pour prévenir le paludisme ?

\textbf{B (Élève 2) :} 我 觉得 \textbf{最重要} 的 是 \textbf{睡觉 用 蚊帐}。\\
Wǒ juéde \textbf{zuì zhòngyào} de shì \textbf{shuìjiào yòng wénzhàng}.\\
Je pense que le plus important, c'est de dormir sous moustiquaire.

\textbf{A :} \textbf{我 同意}。蚊帐 能 挡住 蚊子。\textbf{但是}……\\
Wǒ \textbf{tóngyì}. Wénzhàng néng dǎngzhù wénzi. \textbf{Dànshì}…\\
Je suis d'accord. La moustiquaire bloque les moustiques. Mais…

\textbf{B :} 但是 什么？\\
Dànshì shénme?\\
Mais quoi ?

\textbf{A :} \textbf{光 用 蚊帐 不够}。白天 蚊子 也 咬 人。\textbf{还要} 清理 积水。\\
\textbf{Guāng yòng wénzhàng bùgòu}. Báitiān wénzi yě yǎo rén. \textbf{Hái yào} qīnglǐ jīshuǐ.\\
\textbf{Seulement} la moustiquaire ne suffit pas. Le jour aussi les moustiques piquent. \textbf{Il faut aussi} éliminer l'eau stagnante.

\textbf{B :} \textbf{有 道理}。环境 卫生 是 根本。我们 \textbf{一起} 动员 同学 大扫除 吧？\\
\textbf{Yǒu dàolǐ}. Huánjìng wèishēng shì gēnběn. Wǒmen \textbf{yīqǐ} dòngyuán tóngxué dàsǎochú ba?\\
C'est vrai / Tu as raison. L'hygiène de l'environnement est la racine. On mobilise les camarades pour un grand nettoyage ?

\textbf{A :} \textbf{好主意}！\textbf{那就 这么 办}。\\
\textbf{Hǎo zhǔyi}! \textbf{Nà jiù zhème bàn}.\\
Bonne idée ! C'est décidé / Faisons comme ça.
\end{textebox}

\begin{exemplebox}[Marqueurs de l'interaction orale (Connecteurs logiques)]
\begin{itemize}
    \item \zh{我觉得 / 我认为} (Wǒ juéde / Wǒ rènwéi) — « Je pense que / À mon avis » (Donner son avis).
    \item \zh{我同意 / 我赞成} (Wǒ tóngyì / Wǒ zànchéng) — « Je suis d'accord / J'approuve » (Approuver).
    \item \zh{有道理} (Yǒu dàolǐ) — « C'est logique / Tu as raison » (Valider l'argument de l'autre).
    \item \zh{但是 / 可是} (Dànshì / Kěshì) — « Mais / Cependant » (Nuancer, contredire poliment).
    \item \zh{还有 / 另外} (Háiyǒu / Lìngwài) — « En plus / De plus » (Ajouter un argument).
    \item \zh{所以 / 因此} (Suǒyǐ / Yīncǐ) — « Donc / Par conséquent » (Conclure).
    \item \zh{那 / 那就……吧} (Nà / Nà jiù… ba) — « Alors / Dans ce cas… » (Prendre une décision commune).
\end{itemize}
\end{exemplebox}

\coursec{Formules fonctionnelles de l'exposé et du débat}

\courssub{Boîte à outils orale : Présenter, Argumenter, Interagir}

Ces formules sont des \textbf{blocs préfabriqués} (chunks) à mémoriser par cœur pour fluidifier la prise de parole en continu (exposé) et en interaction (débat). Elles couvrent les actes de parole attendus au programme de Terminale.

\begin{definition}[Formules pour l'Exposé (Expression orale en continu)]
\begin{center}
\begin{tabularx}{\linewidth}{|X|X|X|}
\hline
\rowcolor{popblueL}
\textbf{Fonction} & \textbf{Chinois (Caractères + Pinyin)} & \textbf{Français} \\
\hline
Accroche / Sujet & \zh{今天我的话题是……} (Jīntiān wǒ de huàtí shì…) & Mon sujet aujourd'hui est… \\
\hline
Accroche / Sujet & \zh{我想跟大家分享……} (Wǒ xiǎng gēn dàjiā fēnxiǎng…) & Je veux partager avec vous… \\
\hline
Structurer (1) & \zh{第一，…… 第二，…… 第三，……} (Dì yī… Dì èr… Dì sān…) & Premièrement… Deuxièmement… Troisièmement… \\
\hline
Structurer (2) & \zh{一方面……，另一方面……} (Yī fāngmiàn… lìng yī fāngmiàn…) & D'un côté… de l'autre côté… \\
\hline
Donner exemple & \zh{比如说…… / 例如……} (Bǐrúshuō… / Lìrú…) & Par exemple… / Par exemple… \\
\hline
Cause / Conséquence & \zh{因为……，所以……} (Yīnwèi… suǒyǐ…) & Parce que…, donc… \\
\hline
Proposer & \zh{我建议大家……} (Wǒ jiànyì dàjiā…) & Je suggère à tous de… \\
\hline
Conclure & \zh{总之，……} (Zǒngzhī…) & En résumé… / Bref… \\
\hline
Conclure & \zh{希望大家都能……} (Xīwàng dàjiā dōu néng…) & J'espère que tout le monde pourra… \\
\hline
\end{tabularx}
\end{center}
\end{definition}

\begin{definition}[Formules pour le Débat (Expression orale en interaction)]
\begin{center}
\begin{tabularx}{\linewidth}{|X|X|X|}
\hline
\rowcolor{poporangeL}
\textbf{Fonction} & \textbf{Chinois (Caractères + Pinyin)} & \textbf{Français} \\
\hline
Demander l'avis & \zh{你怎么看？ / 你觉得呢？} (Nǐ zěnme kàn? / Nǐ juéde ne?) & Qu'en penses-tu ? / Ton avis ? \\
\hline
Exprimer accord & \zh{我完全同意。/ 赞成！} (Wǒ wánquán tóngyì. / Zànchéng!) & Je suis tout à fait d'accord. / Approuvé ! \\
\hline
Exprimer accord partiel & \zh{你说得对，但是……} (Nǐ shuō duì le, dànshì…) & Tu as raison, mais… \\
\hline
Exprimer désaccord & \zh{我怕不太对。/ 我不太同意。} (Wǒ pà bù tài duì. / Wǒ bù tài tóngyì.) & Je crains que ce ne soit pas juste. / Je ne suis pas trop d'accord. \\
\hline
Relancer / Creuser & \zh{为什么？ / 能不能具体说说？} (Wèishénme? / Néng bù néng jùtǐ shuōshuo?) & Pourquoi ? / Peux-tu être plus précis ? \\
\hline
Demander précision & \zh{这是什么意思？} (Zhè shì shénme yìsi?) & Qu'est-ce que cela veut dire ? \\
\hline
Prendre décision & \zh{那就这么说定了。} (Nà jiù zhème shuō dìng le.) & C'est dit / C'est décidé comme ça. \\
\hline
Politesse / Tour de parole & \zh{请说。/ 你先说。} (Qǐng shuō. / Nǐ xiān shuō.) & Parlez. / Vous d'abord. \\
\hline
\end{tabularx}
\end{center}
\end{definition}

\coursec{Travail des tons et de la prononciation}

\courssub{Protocole « 3-2-1 » et carte mentale lexicale}

La maîtrise des tons est la condition \textbf{sine qua non} de l'intelligibilité. Nous travaillons d'abord en chorale (groupe classe), puis individuellement (échantillon de 4-5 élèves), enfin en binômes.

\begin{methode}[Protocole de répétition « 3-2-1 »]
\begin{enumerate}
    \item \textbf{Écoute-Modèle (x3) :} Le professeur (ou l'audio) prononce le mot isolé, puis dans une phrase courte. Les élèves écoutent sans parler, visualisent la courbe de ton (main en l'air).
    \item \textbf{Répétition Chorale (x2) :} Toute la classe répète ensemble. Le professeur corrige la justesse globale (ex: « Attention au 3e ton de \zh{回收}, ça descend puis remonte ! »).
    \item \textbf{Relais Individuel (x1 par élève) :} Le professeur désigne un élève : « \zh{垃圾} ». L'élève dit le mot + une phrase : « \zh{垃圾会污染环境} ». Passage rapide pour garder le rythme.
    \item \textbf{Binôme « Miroir » :} Face à face, A dit le mot, B répète et donne le sens en français. Inversion des rôles au signal.
\end{enumerate}
\end{methode}

\begin{popfigure}
\begin{center}\tcbox[colback=popblue,colframe=popblue,coltext=white,arc=9pt,boxsep=4pt,left=6pt,right=6pt]{\bfseries\small Thème Environnement \& Santé}\end{center}
\vspace{-2pt}
\begin{tcbraster}[raster columns=2,raster equal height=rows,raster column skip=6pt,raster row skip=6pt,enhanced,blankest]
\begin{tcolorbox}[enhanced,colback=white,colframe=popgreen,boxrule=0.9pt,arc=3pt,title={回收 Recyclage},fonttitle=\bfseries\scriptsize,coltitle=white,colbacktitle=popgreen,left=4pt,right=4pt,top=2pt,bottom=2pt,boxsep=1pt,toptitle=1.5pt,bottomtitle=1.5pt,halign=left]
\begin{itemize}[nosep,leftmargin=*,itemsep=1pt,font=\scriptsize]
\item 回收站 Centre de tri
\item 瓶子 Bouteille
\item 塑料 Plastique
\item 保护环境 Protéger l'env.
\end{itemize}
\end{tcolorbox}
\begin{tcolorbox}[enhanced,colback=white,colframe=poporange,boxrule=0.9pt,arc=3pt,title={疟疾 Paludisme},fonttitle=\bfseries\scriptsize,coltitle=white,colbacktitle=poporange,left=4pt,right=4pt,top=2pt,bottom=2pt,boxsep=1pt,toptitle=1.5pt,bottomtitle=1.5pt,halign=left]
\begin{itemize}[nosep,leftmargin=*,itemsep=1pt,font=\scriptsize]
\item 蚊子 Moustique
\item 发烧 Fièvre
\item 蚊帐 Moustiquaire
\item 预防 Prévention
\end{itemize}
\end{tcolorbox}
\end{tcbraster}

\legende{Carte mentale lexicale : les deux pôles du thème (Recyclage / Paludisme) et leurs satellites lexicaux HSK 2-3.}
\end{popfigure}

\courssub{Exercices ciblés sur les tons difficiles}

\begin{experience}[Discrimination auditive : Paires minimales tonales]
\textbf{Consigne :} Le professeur lit une colonne (A ou B), les élèves lèvent la carte « A » ou « B ». Puis inversion : élèves dictent au prof.
\begin{center}
\begin{tabular}{ccc}
\hline
\textbf{A (Ton 2 / Ton neutre)} & \textbf{vs} & \textbf{B (Ton 2 / Ton 4)} \\
\hline
\zh{蚊子} (wénzi) — moustique & & \zh{文字} (wénzì) — écriture \\
\zh{瓶子} (píngzi) — bouteille & & \zh{笔子} (bǐzi) — stylo (argot) / \zh{鼻子} (bízi) — nez \\
\zh{桌子} (zhuōzi) — table & & \zh{着急} (zháojí) — anxieux \\
\hline
\end{tabular}
\end{center}
\textbf{Focus Terminale :} Le ton neutre (轻声 qīngshēng) sur les suffixes \zh{子}, \zh{们}, \zh{了}, \zh{的} est court, léger, sans contour de ton propre. Il prend sa hauteur du ton précédent.
\end{experience}

\begin{experience}[Chaînes tonales : Enchaînement 3-3 et 4-4]
\textbf{Consigne :} Répéter en chorale ces séquences fréquentes dans le thème, en respectant la règle du changement de ton (3+3 $\rightarrow$ 2+3) et la chute du 4e ton.
\begin{itemize}
    \item \zh{回收 站} (huíshōu zhàn) $\rightarrow$ 2-1-4 (pas de changement).
    \item \zh{保护 环境} (bǎohù huánjìng) $\rightarrow$ 3-4-2-4 (3e ton \zh{保} devient 2e ton devant \zh{护} 4e ton ? Non, \zh{保} est 3, \zh{护} est 4. Pas de sandhi 3-3 ici. \zh{环} 2, \zh{境} 4. Correct).
    \item \zh{预防 疟疾} (yùfáng nüèjí) $\rightarrow$ 4-2-4-2 (Enchaînement 4-4 : \zh{防} 2, \zh{疟} 4. \zh{疟} 4, \zh{疾} 2. Attention à ne pas « chanter » le 4e ton).
    \item \zh{清理 积水} (qīnglǐ jīshuǐ) $\rightarrow$ 1-3-1-3.
    \item \zh{好主意} (hǎo zhǔyi) $\rightarrow$ 3-3-4 $\rightarrow$ \textbf{2-3-4} (Règle du 3e ton : \zh{好} 3 + \zh{主} 3 $\rightarrow$ \zh{好} se prononce 2e ton).
\end{itemize}
\end{experience}

\coursec{Jeux de rôle et débat guidé}

\courssub{Jeu de rôle 1 : « Le délégué écolo » (Situation formelle)}

\begin{experience}[Jeu de rôle : Présenter un projet au proviseur]
\textbf{Contexte :} Vous êtes délégué(e) de classe. Vous avez 2 minutes pour convaincre le proviseur (joué par un camarade ou le prof) d'installer des poubelles de tri sélectif (垃圾分类桶) dans la cour.
\\
\textbf{Rôles :}
\begin{itemize}
    \item \textbf{Élève A (Délégué) :} Utilise le modèle d'exposé (Dialogue 1) + formules « Proposition » + « Conclusion ».
    \item \textbf{Élève B (Proviseur) :} Pose des questions : \zh{多少钱？} (Combien ça coûte ?), \zh{谁管理？} (Qui gère ?), \zh{有什么困难？} (Quelles difficultés ?). Utilise formules « Relancer », « Demander précision ».
\end{itemize}
\textbf{Contraintes linguistiques obligatoires :}
\begin{enumerate}
    \item Utiliser \zh{应该} au moins 1 fois (devoir moral).
    \item Utiliser \zh{要} au moins 1 fois (nécessité pratique).
    \item Utiliser \zh{建议} (proposer) ou \zh{希望} (espérer).
    \item Minimum 5 phrases complètes.
\end{enumerate}
\textbf{Durée :} Préparation 5 min (notes autorisées en pinyin/mots-clés) + Passage 2 min $\times$ 6 binômes = 12 min.
\end{experience}

\courssub{Jeu de rôle 2 : « Au dispensaire » (Situation semi-formelle)}

\begin{experience}[Jeu de rôle : Décrire des symptômes et comprendre les consignes]
\textbf{Contexte :} Un élève a de la fièvre depuis hier. Il va au dispensaire scolaire (ou hôpital). L'infirmier(ère) l'interroge et donne des consignes.
\\
\textbf{Rôles :}
\begin{itemize}
    \item \textbf{Élève A (Malade) :} Décrit symptômes : \zh{我发烧了，头疼，浑身没劲。} (Fièvre, mal à la tête, fatigue). Répond aux questions.
    \item \textbf{Élève B (Infirmier) :} Interroge : \zh{几度？} (Combien de degrés ?), \zh{吃药了吗？} (Tu as pris des médocs ?), \zh{睡蚊帐吗？} (Tu dors sous moustiquaire ?). Donne consignes avec \zh{要} : \zh{要多喝水，要按时吃药，要复查。}
\end{itemize}
\textbf{Lexique spécifique à fournir au tableau :} \zh{体温} tǐwēn (température), \zh{度} dù (degré), \zh{头疼} tóuténg (mal à la tête), \zh{浑身没劲} húnshēn méijìn (fatigue générale), \zh{验血} yànxiě (analyse de sang), \zh{按时} ànshí (à heures fixes).
\end{experience}

\courssub{Débat guidé : « Recyclage vs Paludisme : quelle priorité ? »}

\begin{experience}[Débat structuré : « Le grand oral »]
\textbf{Sujet :} \zh{“为了预防疟疾，回收塑料瓶子比发蚊帐更重要。”} (« Pour prévenir le paludisme, recycler les bouteilles est plus important que distribuer des moustiquaires. »)
\\
\textbf{Organisation :} 2 équipes de 4-5 élèves (Pour / Contre) + 1 jury de 3 élèves (évalue : arguments, chinois, interaction, prononciation).
\\
\textbf{Déroulement (15 min) :}
\begin{enumerate}
    \item \textbf{Préparation (5 min) :} Chaque équipe prépare 3 arguments structurés (\zh{第一/第二/第三}, \zh{因为/所以}, \zh{比如}). Le jury prépare la grille d'évaluation.
    \item \textbf{Exposés initiaux (2 $\times$ 1 min) :} Un orateur par équipe présente la position.
    \item \textbf{Libre échange / Réfutation (6 min) :} Les équipes s'interpellent. Le professeur modère, note les belles phrases / erreurs récurrentes au tableau pour correction différée.
    \item \textbf{Synthèse du jury (2 min) :} Le jury donne son avis en chinois : \zh{我们觉得……队赢了，因为……} (Nous pensons que l'équipe… a gagné, car…).
\end{enumerate}
\textbf{Arguments types « Pour » (Recyclage prioritaire) :}
\begin{itemize}
    \item \zh{塑料瓶子是蚊子的家。} (Les bouteilles sont la maison des moustiques — source).
    \item \zh{回收解决根本问题，蚊帐只保护个人。} (Le recyclage règle le problème à la source, la moustiquaire ne protège que l'individu).
    \item \zh{保护环境，长远利益大。} (Protéger l'environnement, bénéfice à long terme).
\end{itemize}
\textbf{Arguments types « Contre » (Moustiquaire prioritaire) :}
\begin{itemize}
    \item \zh{蚊帐见效快，直接救命。} (Moustiquaire effet rapide, sauve des vies directement).
    \item \zh{回收要时间、设备、钱，农村难做。} (Recyclage prend temps, équipement, argent, difficile en zone rurale).
    \item \zh{积水不只有瓶子，轮胎、竹筒也有。} (Eau stagnante pas que bouteilles : pneus, bambous aussi).
\end{itemize}
\end{experience}

\begin{exoresolu}[Modélisation : Préparer sa fiche de débat (Côté « Pour »)]
\textbf{Énoncé :} Rédigez les notes (mots-clés + structures) pour l'orateur de l'équipe « Pour » (Recyclage prioritaire). Utilisez le vocabulaire et les formules de la leçon.
\\
\textbf{Solution détaillée :}
\begin{enumerate}
    \item \textbf{Ouverture :} \zh{大家好。我方观点：回收塑料瓶子更重要。} (Bonjour. Notre avis : recycler bouteilles plus important.)
    \item \textbf{Argument 1 (Source) :} \zh{第一，塑料瓶子积水是蚊子繁殖地。回收瓶子 = 消灭蚊子的家。} (Premièrement, bouteilles eau stagnante = gîte larvaire. Recycler = détruire maison moustiques.) \textit{Structure : Sujet + 是 + Complément. Verbe + Objet = Résultat.}
    \item \textbf{Argument 2 (Collectif vs Individuel) :} \zh{第二，蚊帐只保护一个人，回收保护整个社区。我们应该为大家着想。} (Deuxièmement, moustiquaire protège un seul, recyclage protège communauté. Nous devrions penser à tous.) \textit{Structure : \cle{应该} + Verbe. Comparaison \cle{比} implicite.}
    \item \textbf{Argument 3 (Durable) :} \zh{第三，回收还能减少污染，保护环境，长远利益大。} (Troisièmement, recyclage réduit pollution, protège environnement, grand bénéfice long terme.) \textit{Vocabulaire : \cle{减少} jiǎnshǎo (réduire), \cle{长远} chángyuǎn (long terme).}
    \item \textbf{Conclusion :} \zh{所以，请支持回收！谢谢。} (Donc, soutenez le recyclage ! Merci.)
\end{enumerate}
\textbf{Note du correcteur :} L'orateur doit prononcer distinctement \zh{繁殖地} (fánzhí dì), \zh{社区} (shèqū), \zh{长远} (chángyuǎn). Attention au ton de \zh{为} (wèi / wéi) : ici \zh{为大家着想} (wèi dàjiā zhe xiǎng) = « pour le bien de tous », \zh{为} se lit 4e ton (wèi).
\end{exoresolu}

\begin{savaistu}[Le tri obligatoire en Chine : \cle{垃圾分类}]
Depuis le 1er juillet 2019, Shanghai a mis en œuvre la réglementation la plus stricte de Chine en matière de \zh{垃圾分类} (tri des déchets). Les déchets sont classés en 4 catégories codées par couleur :
\begin{itemize}
    \item \zh{可回收物} (kě huíshōu wù) — \textbf{Bleu} : Papier, plastique, verre, métal, textile (nos \zh{瓶子}, \zh{塑料}).
    \item \zh{有害垃圾} (yǒuhài lājī) — \textbf{Rouge} : Piles, ampoules, médicaments périmés, peinture.
    \item \zh{厨余垃圾} (chúyú lājī) — \textbf{Vert} : Épluchures, restes de repas, coquilles d'œufs (compost).
    \item \zh{其他垃圾} (qítā lājī) — \textbf{Gris/Noir} : Tout le reste (poussière, mégots, porcelaine cassée).
\end{itemize}
Ne pas trier expose à une amende (200 yuans pour les particuliers). Pékin, Shenzhen et d'autres métropoles ont suivi. Au Cameroun, le tri sélectif commence timidement (bacs jaunes/verts à Yaoundé), mais la filière de recyclage industrielle reste à développer. C'est un point de comparaison excellent pour un exposé oral !
\end{savaistu}

\begin{exercice}[Entraînement individuel : Expression orale enregistrée]
\textbf{Consigne :} Enregistrez-vous (téléphone / tablette) pour les deux tâches ci-dessous. Durée max 1 min 30 par tâche. Écoutez-vous, notez 2 points à améliorer, ré-enregistrez une fois.
\begin{enumerate}
    \item \textbf{Exposé (Continu) :} « Présentez en 1 min une action concrète pour améliorer l'hygiène dans votre quartier / lycée ». Utilisez : \zh{大家好}, \zh{今天我想谈谈}, \zh{首先/其次/最后}, \zh{我建议}, \zh{谢谢大家}. Mots-clés obligatoires : \zh{垃圾}, \zh{回收}, \zh{环境}, \zh{应该}.
    \item \textbf{Interaction (Simulée) :} Répondez oralement aux 3 questions d'un ami imaginaire (lisez la question, pausez, répondez).
    \begin{itemize}
        \item Q1 : \zh{你觉得得了疟疾怎么办？} (Que faire si on a le paludisme ?)
        \item Q2 : \zh{蚊帐贵不贵？买得到吗？} (La moustiquaire est-elle chère ? Trouvable ?)
        \item Q3 : \zh{除了蚊帐，还能怎么预防？} (À part la moustiquaire, comment prévenir ?)
    \end{itemize}
    Utilisez : \zh{我觉得}, \zh{要}, \zh{去医院}, \zh{吃药}, \zh{清理积水}, \zh{还有}.
\end{enumerate}
\textbf{Auto-évaluation (Grille) :} 1. Tons corrects (sur mots-clés) /5 — 2. Fluidité (pas de longs silences) /5 — 3. Structures (\zh{应该}/\zh{要}/\zh{因为/所以}) /5 — 4. Vocabulaire pertinent /5. Total /20.
\end{exercice}

\begin{corrige}[Exercice Entraînement individuel]
\textbf{Propositions de réponses types (niveau Terminale A / HSK 3-4) :}

\textbf{Tâche 1 (Exposé) :}
\zh{大家好！今天我想谈谈我们学校的卫生问题。首先，操场边有很多垃圾，比如塑料瓶子和塑料袋。其次，雨天积水，容易有蚊子，传播疟疾。所以，我建议大家每天放学前打扫一下，分类扔垃圾，把瓶子放进蓝色回收桶。保护环境是我们的责任。谢谢大家！}
(Pinyin : Dàjiā hǎo! Jīntiān wǒ xiǎng tántan wǒmen xuéxiào de wèishēng wèntí. Shǒuxiān, cāochǎng biān yǒu hěnduō lājī, bǐrúshuō sùliào píngzi hé sùliào dài. Qícì, yǔtiān jīshuǐ, róngyì yǒu wénzi, chuánbō nüèjí. Suǒyǐ, wǒ jiànyì dàjiā měitiān fàngxué qián dǎsǎo yíxià, fēnlèi rēng lājī, bǎ píngzi fàng jìn lánsè huíshōu tǒng. Bǎohù huánjìng shì wǒmen de zérèn. Xièxie dàjiā!)

\textbf{Tâche 2 (Interaction) :}
\begin{itemize}
    \item \textbf{Réponse Q1 :} \zh{我觉得得了疟疾，要马上去医院验血，吃医生给的药。不能自己买药乱吃。} (Wǒ juéde déle nüèjí, yào mǎshàng qù yīyuàn yànxiě, chī yīshēng gěi de yào. Bùnéng zìjǐ mǎi yào luàn chī.)
    \item \textbf{Réponse Q2 :} \zh{蚊帐不太贵，超市和药店都买得到。学校有时也发免费的。} (Wénzhàng bù tài guì, chāoshì hé yàodiàn dōu mǎi dé dào. Xuéxiào yǒushí yě fā miǎnfèi de.)
    \item \textbf{Réponse Q3 :} \zh{除了蚊帐，还要清理家里的积水，比如花盆底下、旧轮胎。还要关纱窗，穿长袖衣服。} (Chúle wénzhàng, hái yào qīnglǐ jiālǐ de jīshuǐ, bǐrúshuō huāpén dǐxià, jiù lúntāi. Hái yào guān shāchuāng, chuān chángxiù yīfu.)
\end{itemize}
\textbf{Points de vigilance pour la correction :} Vérifier \zh{疟疾} (nüèjí), \zh{验血} (yànxiě), \zh{买得到} (mǎi dé dào - complément de possibilité), \zh{清理} (qīnglǐ), \zh{纱窗} (shāchuāng - moustiquaire de fenêtre).
\end{corrige}

\coursec{Dialogue modèle 1 : exposé sur le recyclage}

Dans la section précédente, nous avons découvert le lexique oral du thème « environnement et santé » : \zh{回收} (huíshōu, recycler), \zh{垃圾} (lājī, déchets), \zh{环境} (huánjìng, environnement), \zh{污染} (wūrǎn, pollution). Nous allons maintenant mettre ce vocabulaire en action dans une situation de communication très fréquente au lycée : \cle{prendre la parole devant la classe} pour présenter un petit exposé. Le dialogue modèle ci-dessous se déroule dans une classe de Terminale à Yaoundé : un élève, A, présente un exposé sur le recyclage ; un camarade, B, lui pose une question.

\courssub{Écoute et lecture du dialogue modèle}

\begin{textebox}[Dialogue : un exposé sur le recyclage]
A : 大家好！今天我要谈一谈回收。\\
\emph{Dàjiā hǎo! Jīntiān wǒ yào tán yi tán huíshōu.}\\
« Bonjour à tous ! Aujourd'hui, je vais parler du recyclage. »\\[4pt]
B : 回收有什么好处？\\
\emph{Huíshōu yǒu shénme hǎochu?}\\
« Quels sont les avantages du recyclage ? »\\[4pt]
A : 回收可以保护环境，也可以减少垃圾。\\
\emph{Huíshōu kěyǐ bǎohù huánjìng, yě kěyǐ jiǎnshǎo lājī.}\\
« Le recyclage peut protéger l'environnement et aussi réduire les déchets. »\\[4pt]
A : 所以，我们应该回收瓶子和塑料。\\
\emph{Suǒyǐ, wǒmen yīnggāi huíshōu píngzi hé sùliào.}\\
« Donc, nous devrions recycler les bouteilles et le plastique. »
\end{textebox}

\textbf{Consignes d'écoute.} Le professeur lit le dialogue deux fois. Première écoute : les élèves écoutent sans lire, et repèrent le sujet de l'exposé. Deuxième écoute : les élèves suivent le texte des yeux et soulignent les mots connus. Ensuite, lecture chorale : toute la classe répète phrase par phrase, en veillant aux tons.

\begin{center}
\begin{tabularx}{\linewidth}{|X|X|X|X|}
\hline
\rowcolor{popblueL} Caractères & Pinyin & Nature & Traduction \\
\hline
大家 & dàjiā & pronom & tout le monde \\
\hline
谈一谈 & tán yi tán & verbe & parler un peu de, causer de \\
\hline
回收 & huíshōu & verbe / nom & recycler ; recyclage \\
\hline
好处 & hǎochu & nom & avantage, bienfait \\
\hline
保护 & bǎohù & verbe & protéger \\
\hline
环境 & huánjìng & nom & environnement \\
\hline
减少 & jiǎnshǎo & verbe & réduire, diminuer \\
\hline
垃圾 & lājī & nom & déchets, ordures \\
\hline
所以 & suǒyǐ & conjonction & donc, c'est pourquoi \\
\hline
应该 & yīnggāi & verbe modal & devoir (conseil) \\
\hline
瓶子 & píngzi & nom & bouteille \\
\hline
塑料 & sùliào & nom & plastique \\
\hline
\end{tabularx}
\end{center}

\courssub{Observation : comment est construit cet exposé ?}

Avant toute règle, observons. Que fait A, exactement, et dans quel ordre ?

\begin{enumerate}
\item Il \cle{salue} le public : \zh{大家好！}
\item Il \cle{annonce le sujet} : \zh{今天我要谈一谈回收。}
\item Il \cle{donne des arguments} (deux, reliés par \zh{也}) : \zh{回收可以保护环境，也可以减少垃圾。}
\item Il \cle{conclut par un conseil} introduit par \zh{所以} : \zh{我们应该回收瓶子和塑料。}
\end{enumerate}

\begin{methode}[Les trois parties d'un bon exposé oral en chinois]
Un exposé, même très court, suit toujours la même architecture :
\begin{enumerate}
\item \zh{开头} \emph{kāitóu} — l'\cle{ouverture} : saluer (\zh{大家好！}) puis annoncer le sujet (\zh{今天我要谈一谈……}).
\item \zh{中间} \emph{zhōngjiān} — le \cle{développement} : donner deux ou trois arguments, reliés par \zh{也} (aussi), \zh{而且} (de plus) ou \zh{因为……所以……} (parce que... donc...).
\item \zh{结尾} \emph{jiéwěi} — la \cle{conclusion} : résumer avec \zh{所以} et proposer une action avec \zh{我们应该……}, puis remercier : \zh{谢谢大家！} (Merci à tous !).
\end{enumerate}
Retenez la formule : \textbf{开头 — 中间 — 结尾}, ouverture — développement — conclusion.
\end{methode}

\begin{popfigure}
\begin{tikzpicture}[
  node distance=0.35cm,
  etape/.style={rectangle, rounded corners=3pt, draw=popblue, fill=popblueL, text width=3.1cm, align=center, minimum height=0.9cm, font=\small},
  fleche/.style={-{Stealth[length=2.5mm]}, thick, popdark}
]
\node[etape, align=center] (n1){1. Salutation\\ \zh{大家好！}};
\node[etape, below=of n1, align=center] (n2){2. Annonce du sujet\\ \zh{今天我要谈一谈……}};
\node[etape, below=of n2, align=center] (n3){3. Arguments (2--3)\\ \zh{可以……也可以……}};
\node[etape, below=of n3, align=center] (n4){4. Conclusion\\ \zh{所以，我们应该……}};
\node[etape, below=of n4, align=center] (n5){5. Remerciement\\ \zh{谢谢大家！}};
\draw[fleche] (n1) -- (n2);
\draw[fleche] (n2) -- (n3);
\draw[fleche] (n3) -- (n4);
\draw[fleche] (n4) -- (n5);
\end{tikzpicture}
\legende{La structure d'un exposé oral simple en chinois : cinq étapes à mémoriser.}
\end{popfigure}

\courssub{La formule d'ouverture : 今天我要谈一谈……}

\begin{textebox}[Formule à mémoriser]
今天我要谈一谈……\\
\emph{Jīntiān wǒ yào tán yi tán...}\\
« Aujourd'hui, je vais parler de... »\\[4pt]
Variantes : 今天我要谈一谈回收。(le recyclage) — 今天我要谈一谈疟疾。(le paludisme) — 今天我要谈一谈保护环境。(la protection de l'environnement)
\end{textebox}

Décomposons la formule : \zh{今天} (aujourd'hui) + \zh{我} (je) + \zh{要} (vouloir, aller — marque ici l'intention proche) + \zh{谈一谈} (parler un peu de). Le mot \zh{要} annonce une action future ou une intention, exactement comme « je vais parler » en français.

\begin{attention}[谈一谈 : la réduplication du verbe]
\zh{谈一谈} est la forme \cle{rédupliquée} du verbe \zh{谈} (parler, discuter). La structure est : verbe + \zh{一} + même verbe (\zh{看一看} regarder un coup, \zh{想一想} réfléchir un peu, \zh{谈一谈} parler un peu). Cette réduplication \textbf{adoucit le ton} et suggère une action brève et informelle : c'est la politesse modeste de l'orateur (« je vais juste vous parler un peu de... »).\\
\textbf{Pièges des francophones :}
\begin{itemize}
\item Ne dites jamais \zh{谈一个谈} : la réduplication se fait avec \zh{一}, pas avec le classificateur \zh{个}.
\item Ne traduisez pas \zh{一} : \zh{谈一谈} ne veut pas dire « parler une fois », mais « parler un peu ».
\item À l'oral, le \zh{一} de \zh{谈一谈} se prononce souvent léger, presque au ton neutre : \emph{tán yi tán}.
\end{itemize}
\end{attention}

\courssub{Deux verbes modaux à ne pas confondre : 可以 et 应该}

Observons deux phrases du dialogue :

\begin{itemize}
\item \zh{回收可以保护环境。} \emph{Huíshōu kěyǐ bǎohù huánjìng.} « Le recyclage \textbf{peut} protéger l'environnement. »
\item \zh{我们应该回收瓶子和塑料。} \emph{Wǒmen yīnggāi huíshōu píngzi hé sùliào.} « Nous \textbf{devrions} recycler les bouteilles et le plastique. »
\end{itemize}

\begin{propriete}[可以 kěyǐ et 应该 yīnggāi]
\cle{可以} exprime la \textbf{possibilité} ou la capacité objective : « pouvoir, il est possible de ». \cle{应该} exprime le \textbf{devoir, le conseil, ce qui est souhaitable} : « devoir (au sens de il faudrait) ». Les deux se placent \textbf{avant le verbe principal}, après le sujet : Sujet + 可以/应该 + Verbe.\\
\zh{我们可以回收。} « Nous pouvons recycler. » $\neq$ \zh{我们应该回收。} « Nous devrions recycler. »\\
Négations : \zh{不可以} (ne pas pouvoir, il est interdit de) ; \zh{不应该} (ne pas devoir, il ne faudrait pas).
\end{propriete}

\begin{center}
\begin{tabularx}{\linewidth}{|X|X|X|}
\hline
\rowcolor{popblueL} Structure & Exemple (caractères + pinyin) & Traduction \\
\hline
Sujet + 可以 + verbe & 回收可以减少污染。Huíshōu kěyǐ jiǎnshǎo wūrǎn. & Le recyclage peut réduire la pollution. \\
\hline
Sujet + 应该 + verbe & 我们应该保护环境。Wǒmen yīnggāi bǎohù huánjìng. & Nous devrions protéger l'environnement. \\
\hline
Sujet + 不应该 + verbe & 我们不应该扔垃圾。Wǒmen bù yīnggāi rēng lājī. & Nous ne devrions pas jeter de déchets. \\
\hline
\end{tabularx}
\end{center}

Dans un exposé, l'alternance des deux modaux est une stratégie efficace : on présente d'abord les \textbf{possibilités} (\zh{可以} : voilà ce que le recyclage peut faire), puis on conclut par le \textbf{devoir} (\zh{应该} : voilà ce que nous devons faire). C'est exactement le mouvement du dialogue modèle.

\begin{exemplebox}[Phrases utiles pour le développement]
\zh{回收可以减少污染。} \emph{Huíshōu kěyǐ jiǎnshǎo wūrǎn.} « Le recyclage peut réduire la pollution. »\\
\zh{我觉得回收很重要。} \emph{Wǒ juéde huíshōu hěn zhòngyào.} « Je pense que le recyclage est très important. »\\
\zh{回收也可以保护我们的城市。} \emph{Huíshōu yě kěyǐ bǎohù wǒmen de chéngshì.} « Le recyclage peut aussi protéger notre ville. »\\
\zh{在雅温得，我们应该回收塑料。} \emph{Zài Yǎwēndé, wǒmen yīnggāi huíshōu sùliào.} « À Yaoundé, nous devrions recycler le plastique. »
\end{exemplebox}

Remarquez \zh{我觉得……} (je pense que...) : c'est la formule la plus simple pour donner son avis. Le verbe \zh{觉得} introduit une opinion personnelle, comme « je trouve que » en français.

\courssub{Travail de mémorisation par substitution}

La \cle{substitution} est la technique reine de l'apprentissage oral : on garde le squelette de la phrase et on ne change qu'un élément. Reprenons la conclusion du dialogue :

\begin{center}
\textbf{我们应该回收 + [X] 。}\quad \emph{Wǒmen yīnggāi huíshōu [X].}
\end{center}

\begin{center}
\begin{tabularx}{\linewidth}{|X|X|X|}
\hline
\rowcolor{popblueL} Élément substitué & Phrase complète & Traduction \\
\hline
瓶子 píngzi (bouteilles) & 我们应该回收瓶子。 & Nous devrions recycler les bouteilles. \\
\hline
纸 zhǐ (papier) & 我们应该回收纸。 & Nous devrions recycler le papier. \\
\hline
玻璃 bōli (verre) & 我们应该回收玻璃。 & Nous devrions recycler le verre. \\
\hline
塑料 sùliào (plastique) & 我们应该回收塑料。 & Nous devrions recycler le plastique. \\
\hline
\end{tabularx}
\end{center}

On peut aussi substituer dans la phrase d'arguments : \zh{回收可以保护 + [环境 / 我们的城市 / 大自然]} (l'environnement / notre ville / la nature). Le professeur montre une image ou dit un mot en français, la classe produit la phrase entière à voix haute, puis un élève seul la répète. Exigez toujours la phrase \textbf{complète}, jamais le mot isolé : c'est la structure entière qui doit devenir automatique.

\begin{exoresolu}[Substitution guidée]
Énoncé. Remplacez l'élément souligné par le mot entre parenthèses et prononcez la phrase complète : 我们应该回收瓶子。(玻璃)
\tcblower
Solution. On remplace \zh{瓶子} (bouteilles) par \zh{玻璃} (verre) sans rien changer d'autre :\\
\zh{我们应该回收玻璃。} \emph{Wǒmen yīnggāi huíshōu bōli.} « Nous devrions recycler le verre. »\\
Attention à la prononciation de \zh{玻璃} : premier caractère au premier ton, deuxième au ton neutre (\emph{bōli}), comme \zh{瓶子} (\emph{píngzi}).
\end{exoresolu}

\courssub{Reformulation : à vous de parler}

Dernière étape : chaque élève \cle{reformule} l'exposé avec ses propres exemples. La reformulation n'est pas une récitation : le squelette (salutation, annonce, arguments, conclusion) reste, mais le contenu change. Modèle de production attendue :

\begin{textebox}[Exemple de reformulation par un élève de Douala]
大家好！今天我要谈一谈回收。\\
\emph{Dàjiā hǎo! Jīntiān wǒ yào tán yi tán huíshōu.}\\
回收可以减少垃圾，也可以保护我们的城市。\\
\emph{Huíshōu kěyǐ jiǎnshǎo lājī, yě kěyǐ bǎohù wǒmen de chéngshì.}\\
我觉得回收很重要。所以，我们应该回收纸和玻璃。\\
\emph{Wǒ juéde huíshōu hěn zhòngyào. Suǒyǐ, wǒmen yīnggāi huíshōu zhǐ hé bōli.}\\
谢谢大家！\\
\emph{Xièxie dàjiā!}\\
« Bonjour à tous ! Aujourd'hui, je vais parler du recyclage. Le recyclage peut réduire les déchets et aussi protéger notre ville. Je pense que le recyclage est très important. Donc, nous devrions recycler le papier et le verre. Merci à tous ! »
\end{textebox}

\begin{attention}[Erreurs fréquentes des francophones à l'oral]
\begin{itemize}
\item \textbf{Ordre des mots} : le complément de lieu se place avant le verbe. Dites \zh{在杜阿拉，我们回收塑料。} (\emph{Zài Dù'ālā, wǒmen huíshōu sùliào.}), jamais « 我们回收塑料在杜阿拉 ».
\item \textbf{Pas de « de » entre les verbes} : \zh{可以保护} se dit tel quel, sans mot-outil entre les deux verbes, contrairement au français « peut protéger » qui cache parfois « de » dans d'autres constructions.
\item \textbf{Tons de 可以} : citation \emph{kěyǐ} = 3\ieme{} ton + 3\ieme{} ton ; à l'oral, le premier 3\ieme{} ton devient un 2\ieme{} ton (changement de ton / sandhi) : prononcez \emph{kéyǐ}.
\item \textbf{好处} : le second caractère est au ton neutre (\emph{hǎochu}), ne le montez pas.
\end{itemize}
\end{attention}

\begin{experience}[Mini-exposés en chaîne]
Par groupes de quatre : l'élève 1 salue et annonce le sujet, l'élève 2 donne le premier argument avec \zh{可以}, l'élève 3 ajoute un argument avec \zh{也} ou donne son avis avec \zh{我觉得}, l'élève 4 conclut avec \zh{所以……应该} et remercie. Puis le groupe se lève et présente son exposé complet devant la classe. Les autres élèves écoutent et doivent pouvoir dire, en français, quels objets le groupe propose de recycler. Critères d'évaluation : respect des cinq étapes, exactitude des tons, phrase complète sans hésitation excessive.
\end{experience}

\begin{aretenir}[L'essentiel de la section]
\begin{itemize}
\item Un exposé oral simple suit cinq étapes : salutation, annonce du sujet, arguments, conclusion, remerciement (\zh{开头 — 中间 — 结尾}).
\item Formule d'ouverture : \zh{今天我要谈一谈……} ; la réduplication \zh{谈一谈} adoucit le ton.
\item \zh{可以} = possibilité ; \zh{应该} = devoir, conseil ; tous deux se placent avant le verbe principal.
\item \zh{我觉得……} pour donner son avis ; \zh{所以} pour conclure.
\item La mémorisation par substitution (瓶子 → 纸 → 玻璃) rend la structure automatique avant la reformulation libre.
\end{itemize}
\end{aretenir}

\coursec{Dialogue modèle 2 : débat sur le paludisme}

\courssub{Transition et mise en situation}

Après avoir travaillé l'exposé structuré sur le recyclage (section précédente), nous passons maintenant à l'\textbf{interaction orale spontanée} : le débat. Le paludisme (\zh{疟疾} \emph{nüèjí}) est un sujet de société majeur au Cameroun comme dans de nombreuses régions de Chine (notamment le Yunnan, le Hainan). Savoir exprimer son accord, son désaccord, justifier son avis et proposer des solutions en chinois est une compétence clé du niveau HSK 3-4 visé en Terminale.

\begin{textebox}[Dialogue modèle : Débat sur la prévention du paludisme]
\textbf{A :} 你觉得疟疾危险吗？\\
Nǐ juéde nüèjí wēixiǎn ma?\\
(Trouves-tu que le paludisme est dangereux ?)

\textbf{B :} 我觉得很危险，因为很多人得疟疾。\\
Wǒ juéde hěn wēixiǎn, yīnwèi hěn duō rén dé nüèjí.\\
(Je trouve que c'est très dangereux, car beaucoup de gens attrapent le paludisme.)

\textbf{A :} 我同意。我们应该怎么预防？\\
Wǒ tóngyì. Wǒmen yīnggāi zěnme yùfáng?\\
(Je suis d'accord. Comment devrions-nous le prévenir ?)

\textbf{B :} 我们应该用蚊帐，也应该看医生。\\
Wǒmen yīnggāi yòng wénzhàng, yě yīnggāi kàn yīshēng.\\
(Nous devrions utiliser des moustiquaires et aussi consulter un médecin.)

\textbf{A :} 你说得对！\\
Nǐ shuō de duì!\\
(Tu as raison !)
\end{textebox}

\courssub{Observation linguistique et analyse des structures}

Observons ce dialogue authentique. Il suit un schéma classique de l'interaction orale : \textbf{Question d'opinion} $\rightarrow$ \textbf{Avis justifié} $\rightarrow$ \textbf{Accord} $\rightarrow$ \textbf{Question de solution} $\rightarrow$ \textbf{Propositions} $\rightarrow$ \textbf{Validation}.

\begin{definition}[Vocabulaire clé du débat santé]
\begin{center}
\begin{tabularx}{\linewidth}{|X|X|X|X|}
\hline
\rowcolor{popblueL} \textbf{Caractères} & \textbf{Pinyin} & \textbf{Nature} & \textbf{Traduction} \\ \hline
疟疾 & nüèjí & nom & paludisme (maladie) \\ \hline
危险 & wēixiǎn & adj. & dangereux \\ \hline
觉得 & juéde & verbe & penser, trouver (avis subjectif) \\ \hline
因为 & yīnwèi & conj. & parce que \\ \hline
得 & dé & verbe & attraper (maladie), contracter \\ \hline
同意 & tóngyì & verbe & être d'accord \\ \hline
应该 & yīnggāi & verbe modal & devoir,应该 \\ \hline
怎么 & zěnme & adv. & comment, de quelle manière \\ \hline
预防 & yùfáng & verbe & prévenir (maladie) \\ \hline
蚊帐 & wénzhàng & nom & moustiquaire \\ \hline
治疗 & zhìliáo & verbe/nom & soigner, traitement \\ \hline
蚊子 & wénzi & nom & moustique \\ \hline
症状 & zhèngzhuàng & nom & symptôme \\ \hline
发烧 & fāshāo & verbe & avoir de la fièvre \\ \hline
\end{tabularx}
\end{center}
\end{definition}

\begin{propriete}[Structure : Exprimer son avis subjectif avec \cle{觉得}]
\textbf{Forme :} \textbf{Sujet + 觉得 + [Phrase d'opinion / Adjectif / Verbe]}
\begin{itemize}
    \item \zh{觉得} \emph{juéde} exprime un jugement personnel, une impression, moins fort que \zh{认为} \emph{rènwéi} (considérer, estimer après réflexion).
    \item La phrase qui suit peut être une phrase complète (Sujet + Verbe/Adj).
\end{itemize}

\begin{center}
\begin{tabularx}{\linewidth}{|X|X|X|}
\hline
\rowcolor{popgreenL} \textbf{Structure} & \textbf{Exemple (Caractères + Pinyin)} & \textbf{Traduction} \\ \hline
Suj + 觉得 + Adj & 我觉得很危险。 (Wǒ juéde hěn wēixiǎn.) & Je trouve que c'est très dangereux. \\ \hline
Suj + 觉得 + Suj + Verbe & 我觉得蚊帐很有用。 (Wǒ juéde wénzhàng hěn yǒuyòng.) & Je trouve que la moustiquaire est très utile. \\ \hline
Question : Suj + 觉得 + ... + 吗? & 你觉得疟疾危险吗？ (Nǐ juéde nüèjí wēixiǎn ma?) & Trouves-tu le paludisme dangereux ? \\ \hline
\end{tabularx}
\end{center}

\textbf{Comparaison FR/CN :} En français, « Je trouve que c'est... » ou « Je pense que... ». En chinois, \zh{觉得} place l'opinion \textbf{avant} le contenu de l'opinion, sans conjonction « que ».
\end{propriete}

\begin{propriete}[Structure : Justifier avec \cle{因为} ... (pas de \cle{所以} obligatoire à l'oral)]
\textbf{Forme :} \textbf{因为 + Cause/Raison , (Sujet) + Conséquence/Conclusion.}
\begin{itemize}
    \item À l'oral, \zh{所以} \emph{suǒyǐ} (donc) est souvent omis. La pause (virgule) marque la relation logique.
    \item \zh{得} \emph{dé} (2e ton) = « attraper/contracter » une maladie. \textbf{Ne pas confondre avec} \zh{得} \emph{de} (particule de complément de degré) ni \zh{得} \emph{děi} (devoir).
\end{itemize}

\begin{center}
\begin{tabularx}{\linewidth}{|X|X|X|}
\hline
\rowcolor{poporangeL} \textbf{Structure} & \textbf{Exemple} & \textbf{Traduction} \\ \hline
因为 + Phrase, Phrase. & 因为很多人得疟疾，(所以)很危险。 & Parce que beaucoup attrapent le paludisme, (donc) c'est dangereux. \\ \hline
因为 + Raisons multiples & 因为又热又潮湿，蚊子多。 & Parce qu'il fait chaud et humide, il y a beaucoup de moustiques. \\ \hline
\end{tabularx}
\end{center}
\end{propriete}

\begin{propriete}[Structures : Accord (\cle{同意}) et Désaccord (\cle{不同意})]
\textbf{Accord :} \zh{我同意。} \emph{Wǒ tóngyì.} / \zh{你说得对！} \emph{Nǐ shuō de duì!} (Tu as raison !) / \zh{没错。} \emph{Méi cuò.} (Exact.)
\textbf{Désaccord poli :} \zh{我不同意。} \emph{Wǒ bù tóngyì.} / \zh{我不太同意。} \emph{Wǒ bù tài tóngyì.} (Je ne suis pas tout à fait d'accord.) / \zh{恐怕不太对。} \emph{Kǒngpà bù tài duì.} (Je crains que ce ne soit pas tout à fait exact.)

\begin{center}
\begin{tabularx}{\linewidth}{|X|X|X|}
\hline
\rowcolor{poppurpleL} \textbf{Fonction} & \textbf{Formule chinoise (Pinyin)} & \textbf{Français} \\ \hline
Accord total & 我同意 / 你说得对 & Je suis d'accord / Tu as raison \\ \hline
Accord partiel & 我基本同意，但是…… & Je suis globalement d'accord, mais... \\ \hline
Désaccord direct & 我不同意，因为…… & Je ne suis pas d'accord, car... \\ \hline
Désaccord atténué & 我不太同意 / 我有不同看法 & Je ne suis pas trop d'accord / J'ai un autre avis \\ \hline
Demander avis & 你怎么看？ / 你有什么建议？ & Qu'en penses-tu ? / Quelle est ta suggestion ? \\ \hline
\end{tabularx}
\end{center}
\end{propriete}

\begin{propriete}[Structure : Proposer des solutions avec \cle{应该} + Verbe]
\textbf{Forme :} \textbf{Sujet + 应该 / 要 / 得 + Verbe (+ Objet).}
\begin{itemize}
    \item \zh{应该} \emph{yīnggāi} : devoir moral, conseil fort (recommandation standard).
    \item \zh{要} \emph{yào} : nécessité, impératif (plus fort, ordre ou contrainte forte).
    \item \zh{得} \emph{děi} (3e ton) : obligation contrainte, « n'avoir d'autre choix que » (oral, familier).
\end{itemize}

\begin{center}
\begin{tabularx}{\linewidth}{|X|X|X|}
\hline
\rowcolor{poppinkL} \textbf{Modal} & \textbf{Exemple} & \textbf{Nuance} \\ \hline
应该 & 我们应该用蚊帐。 & Conseil standard : « Nous devrions utiliser... » \\ \hline
要 & 发烧了要去医院。 & Nécessité vitale : « Il FAUT aller à l'hôpital. » \\ \hline
得 & 晚上得关窗户。 & Contrainte pratique : « Faut fermer les fenêtres le soir. » \\ \hline
\end{tabularx}
\end{center}
\end{propriete}

\courssub{Focus lexical : \cle{预防} vs \cle{治疗} — Prévenir vs Soigner}

\begin{definition}[Distinction cruciale en contexte sanitaire]
\begin{itemize}
    \item \zh{预防} \emph{yùfáng} (verbe) : \textbf{Prévenir}, agir \textbf{avant} la maladie. Ex: \zh{预防针} \emph{yùfángzhēn} (vaccin), \zh{预防措施} \emph{yùfáng cuòshī} (mesures de prévention).
    \item \zh{治疗} \emph{zhìliáo} (verbe/nom) : \textbf{Soigner}, \textbf{traiter}, agir \textbf{après} la déclaration de la maladie. Ex: \zh{治疗方案} \emph{zhìliáo fāng'àn} (protocole de traitement), \zh{及时治疗} \emph{jíshí zhìliáo} (soigner à temps).
\end{itemize}
\textbf{Mnémotechnique :} \zh{预} \emph{yù} = d'avance (pré-) ; \zh{治} \emph{zhì} = gouverner, gérer, ordonner (mettre de l'ordre dans le désordre de la maladie).
\end{definition}

\begin{exemplebox}[Phrases types avec 预防 et 治疗]
\begin{itemize}
    \item \zh{预防疟疾最好的办法是用蚊帐。} \emph{Yùfáng nüèjí zuì hǎo de bànfǎ shì yòng wénzhàng.} (Le meilleur moyen de prévenir le paludisme est d'utiliser une moustiquaire.)
    \item \zh{得了疟疾要马上去医院治疗。} \emph{Dé le nüèjí yào mǎshàng qù yīyuàn zhìliáo.} (Une fois qu'on a attrapé le paludisme, il faut aller à l'hôpital se faire soigner tout de suite.)
    \item \zh{预防胜于治疗。} \emph{Yùfáng shèng yú zhìliáo.} (La prévention vaut mieux que la guérison / Mieux vaut prévenir que guérir.)
\end{itemize}
\end{exemplebox}

\courssub{Variante du dialogue : Exprimer le désaccord et argumenter}

Le dialogue modèle montre un consensus. En débat réel, le désaccord est fréquent. Voici comment le structurer.

\begin{textebox}[Dialogue variante : Désaccord sur la méthode de prévention]
\textbf{A :} 你觉得吃药预防疟疾好不好？\\
Nǐ juéde chī yào yùfáng nüèjí hǎo bu hǎo?\\
(Qu'en penses-tu de prendre des médicaments pour prévenir le paludisme ?)

\textbf{B :} \textbf{我不同意。} 因为吃药有副作用，而且贵。\\
Wǒ \textbf{bù tóngyì}. Yīnwèi chī yào yǒu fùzuòyòng, érqiě guì.\\
\textbf{(Je ne suis pas d'accord.)} Car les médicaments ont des effets secondaires, et en plus c'est cher.

\textbf{A :} 那你有什么更好的建议？\\
Nà nǐ yǒu shénme gèng hǎo de jiànyi?\\
(Alors quelle est ta meilleure suggestion ?)

\textbf{B :} 我觉得用蚊帐又安全又便宜。\\
Wǒ juéde yòng wénzhàng yòu ānquán yòu piányi.\\
(Je trouve que la moustiquaire est à la fois sûre et bon marché.)

\textbf{A :} 有道理。那我们两种方法都用吧。\\
Yǒu dàoli. Nà wǒmen liǎng zhǒng fāngfǎ dōu yòng ba.\\
(C'est logique / Il y a du vrai. Alors utilisons les deux méthodes.)
\end{textebox}

\begin{aretenir}[Formules de relance et de concession pour le débat]
\begin{itemize}
    \item \zh{有道理。} \emph{Yǒu dàoli.} = « C'est logique / Il y a du vrai / Tu as un point. » (Concession polie).
    \item \zh{不过……} \emph{Búguò...} = « Cependant... / Mais... » (Introduit une nuance).
    \item \zh{另外……} \emph{Lìngwài...} = « De plus / En outre... » (Ajoute un argument).
    \item \zh{比如说……} \emph{Bǐrúshuō...} = « Par exemple... » (Illustre).
    \item \zh{总之……} \emph{Zǒngzhī...} = « En résumé / Bref... » (Conclut).
\end{itemize}
\end{aretenir}

\courssub{Travail de la prononciation : Les tons du débat}

La justesse des tons change le sens des mots-clés du débat. Travaillez ces paires minimales et ces phrases en binôme.

\begin{attention}[Pièges de prononciation fréquents pour francophones]
\begin{enumerate}
    \item \zh{疟疾} \emph{nüèjí} : \textbf{4e ton + 2e ton}. \zh{疟} \emph{nüè} (son \emph{üe} = \emph{ü} + \emph{e}, bouche arrondie) $\neq$ \zh{月} \emph{yuè}. \zh{疾} \emph{jí} (court, sec) $\neq$ \zh{鸡} \emph{jī} (poulet, 1er ton).
    \item \zh{得} \emph{dé} (attraper) vs \zh{得} \emph{de} (particule) vs \zh{得} \emph{děi} (devoir).
        \begin{itemize}
            \item \zh{得疟疾} \emph{dé nüèjí} (2e ton) = Attraper le paludisme.
            \item \zh{跑得快} \emph{pǎo de kuài} (ton neutre/léger) = Courir vite.
            \item \zh{我得去} \emph{wǒ děi qù} (3e ton) = Je dois y aller.
        \end{itemize}
    \item \zh{蚊帐} \emph{wénzhàng} : \zh{帐} \emph{zhàng} (4e ton, compte/tente) $\neq$ \zh{长} \emph{zhǎng} (3e ton, grandir) $\neq$ \zh{张} \emph{zhāng} (1er ton, classifier feuilles/billets).
    \item \zh{预防} \emph{yùfáng} : \zh{预} \emph{yù} (4e ton) $\neq$ \zh{雨} \emph{yǔ} (3e ton, pluie).
    \item \zh{治疗} \emph{zhìliáo} : \zh{治} \emph{zhì} (4e ton) $\neq$ \zh{之} \emph{zhī} (1er ton, particule).
\end{enumerate}
\end{attention}

\begin{methode}[Entraînement prosodique en 3 étapes]
\begin{enumerate}
    \item \textbf{Écoute-Modélisation :} L'enseignant (ou l'audio) lit le dialogue modèle. Les élèves identifient les groupes de souffle (pauses) et les mots forts (accent de phrase sur \zh{危险}, \zh{预防}, \zh{蚊帐}, \zh{医生}).
    \item \textbf{Chorale puis individuel :} Répétition en chœur pour la musique de la phrase, puis passage individuel pour corriger les tons lexicaux (\zh{nüèjí}, \zh{yùfáng}, \zh{zhìliáo}).
    \item \textbf{Substitution guidée :} Garder l'intonation du dialogue, changer le vocabulaire (ex: \zh{蚊帐} $\rightarrow$ \zh{驱蚊液} \emph{qūwényè} répulsif ; \zh{看医生} $\rightarrow$ \zh{吃药} \emph{chī yào}).
\end{enumerate}
\end{methode}

\courssub{Organigramme du débat : Structure visuelle de l'interaction}

\begin{popfigure}
\begin{tikzpicture}[
    node distance=1.2cm and 1.5cm,
    every node/.style={font=\small, align=center},
    box/.style={draw=popblue, thick, rounded corners=4pt, fill=popblueL, minimum width=3.5cm, minimum height=1.2cm},
    diamond/.style={draw=poporange, thick, aspect=2, fill=poporangeL, minimum width=3.5cm, minimum height=1.2cm},
    arrow/.style={-Stealth, thick, popdark},
    labelstyle/.style={fill=popcream, rounded corners, font=\tiny, inner sep=2pt}
]
% Nœuds
\node[box, align=center] (start){Question d'opinion\\ \zh{你觉得...吗？}};
\node[diamond, below of=start, align=center] (avis){Avis justifié\\ \zh{我觉得...，因为...}};
\node[diamond, below of=avis, align=center] (reaction){Réaction : Accord / Désaccord\\ \zh{我同意} / \zh{我不同意}};
\node[box, below left=1.5cm and 2cm of reaction, align=center] (accord){Accord + Question solution\\ \zh{我们应该怎么...？}};
\node[box, below right=1.5cm and 2cm of reaction, align=center] (desaccord){Désaccord + Contre-argument\\ \zh{因为...，所以...}};
\node[box, below=2.5cm of reaction, align=center] (propositions){Propositions de solutions\\ \zh{应该 / 要 / 得 + Verbe}};
\node[box, below of=propositions, align=center] (conclusion){Validation / Synthèse\\ \zh{你说得对 / 有道理 / 总之}};

% Flèches
\draw[arrow] (start) -- (avis);
\draw[arrow] (avis) -- (reaction);
\draw[arrow] (reaction) -| node[labelstyle, pos=0.5, left] {Oui} (accord);
\draw[arrow] (reaction) -| node[labelstyle, pos=0.5, right] {Non} (desaccord);
\draw[arrow] (accord) |- (propositions);
\draw[arrow] (desaccord) |- (propositions);
\draw[arrow] (propositions) -- (conclusion);
\end{tikzpicture}
\legende{Organigramme de la structure interactionnelle d'un débat oral (Question $\rightarrow$ Avis $\rightarrow$ Accord/Désaccord $\rightarrow$ Solutions $\rightarrow$ Conclusion).}
\end{popfigure}

\courssub{Exercice résolu : Transformer un accord en désaccord argumenté}

\begin{exoresolu}[Transformer l'échange]
\textbf{Consigne :} Partez du dialogue modèle (accord total). Réécrivez la réplique de B (ligne 4) et la conclusion de A (ligne 5) pour exprimer un \textbf{désaccord poli} : B préfère les répulsifs (\zh{驱蚊液} \emph{qūwényè}) aux moustiquaires (trop chaudes), A concède le point mais insiste sur le coût.

\textbf{Solution détaillée :}

\textbf{B (Désaccord + Justification) :}
\zh{我不同意。我觉得蚊帐太热了，睡觉不舒服。我们应该用驱蚊液，又凉快又方便。}
\emph{Wǒ bù tóngyì. Wǒ juéde wénzhàng tài rè le, shuìjiào bù shūfu. Wǒmen yīnggāi yòng qūwényè, yòu liángkuai yòu fāngbiàn.}
(Je ne suis pas d'accord. Je trouve la moustiquaire trop chaude, inconfortable pour dormir. Nous devrions utiliser du répulsif, c'est à la fois frais et pratique.)

\textbf{A (Concession partielle + Nouvel argument) :}
\zh{有道理，驱蚊液确实方便。但是蚊帐便宜，不用买很多瓶。两个方法一起用最好。}
\emph{Yǒu dàoli, qūwényè quèshí fāngbiàn. Dànshì wénzhàng piányi, bù yòng mǎi hěn duō píng. Liǎng gè fāngfǎ yīqǐ yòng zuì hǎo.}
(C'est vrai, le répulsif est effectivement pratique. Mais la moustiquaire est bon marché, pas besoin d'acheter beaucoup de flacons. Utiliser les deux méthodes ensemble est le mieux.)

\textbf{Analyse grammaticale :}
\begin{itemize}
    \item \zh{太...了} \emph{tài... le} : structure d'excès (trop chaud).
    \item \zh{又...又...} \emph{yòu... yòu...} : coordination de deux qualités positives (frais ET pratique).
    \item \zh{确实} \emph{quèshí} : adverbe « effectivement / vraiment » (renforce l'accord partiel).
    \item \zh{不用} \emph{bù yòng} : « pas besoin de » (absence de nécessité).
\end{itemize}
\end{exoresolu}

\courssub{Jeu de rôle guidé : Moustiquaire vs Médicaments préventifs}

\begin{experience}[Débat en binôme : « La meilleure prévention »]
\textbf{Contexte :} Deux élèves (ou groupes) débattent dans une classe de Terminale à Yaoundé ou Douala.
\textbf{Rôle A (Pro-Moustiquaire) :} Défendez l'usage de la moustiquaire imprégnée (\zh{驱蚊蚊帐} \emph{qūwén wénzhàng}) : gratuite (campagnes gov.), protection physique, pas d'effets secondaires, protège toute la famille.
\textbf{Rôle B (Pro-Médicaments/Autres) :} Défendez la chimio-prophylaxie (\zh{预防服药} \emph{yùfáng fúyào}) ou les répulsifs : moustiquaire déchirée/chaude, médicaments efficaces pour les voyages/zones à risque, rapide à prendre.
\textbf{Contraintes linguistiques obligatoires :}
\begin{enumerate}
    \item Utiliser \zh{我觉得} / \zh{我认为} pour l'avis.
    \item Utiliser \zh{因为...} pour justifier (au moins 2 arguments).
    \item Utiliser \zh{我不同意} / \zh{我不太同意} + \zh{因为...} pour contredire.
    \item Utiliser \zh{应该} / \zh{要} pour la recommandation finale.
    \item Conclure par \zh{有道理} / \zh{你说得对} / \zh{总之}...
\end{enumerate}
\textbf{Durée :} 3 min préparation (notes en pinyin/mots-clés autorisés), 2 min débat oral pur (sans notes), 1 min feedback pair (tons, structures).
\end{experience}

\courssub{Le saviez-vous ? : Paludisme, Cameroun et Chine}

\begin{savaistu}[Paludisme : Réalités camerounaises et coopération sino-africaine]
\begin{itemize}
    \item \textbf{Cameroun :} Le paludisme (\zh{疟疾} \emph{nüèjí}) reste la \textbf{première cause de consultation} (environ 30-40\% des consultations en structures de santé publiques) et la première cause de mortalité chez l'enfant de moins de 5 ans. Le Ministère de la Santé Publique (MINSANTE) mène des campagnes de \textbf{distribution gratuite de moustiquaires imprégnées à longue durée d'action (MILDA)} (\zh{长效驱蚊蚊帐} \emph{chángxiào qūwén wénzhàng}) tous les 3 ans (dernière en 2022/2023). Le traitement de première ligne est l'\textbf{Artémisinine} (\zh{青蒿素} \emph{qīnghāosù}) en thérapie combinée (ACT).
    \item \textbf{Chine :} La Chine a été certifiée \textbf{« exempte de paludisme »} par l'OMS en \textbf{juin 2021} (zéro cas autochtone pendant 4 ans consécutifs). C'est une victoire majeure après 70 ans de lutte. La découverte de l'\textbf{artémisinine} (\zh{青蒿素} \emph{qīnghāosù}) par \textbf{Tu Youyou} (\zh{屠呦呦} \emph{Tú Yōuyōu}), prix Nobel de médecine 2015, est issue de la pharmacopée chinoise traditionnelle (\zh{中医} \emph{zhōngyī}).
    \item \textbf{Coopération :} La Chine partage son expérience (\zh{经验} \emph{jīngyàn}) et ses technologies (diagnostic rapide, médicaments, moustiquaires) avec le Cameroun via des centres de coopération (ex: Centre de lutte antipaludique Chine-Cameroun à Yaoundé). Des experts chinois (\zh{专家} \emph{zhuānjiā}) forment le personnel soignant camerounais.
\end{itemize}
\textbf{Vocabulaire bonus :} \zh{根除} \emph{gēnchú} (éradiquer), \zh{流行病} \emph{liúxíngbìng} (épidémie), \zh{按蚊} \emph{ànwén} (anophèle / moustique vecteur), \zh{免疫} \emph{miǎnyì} (immunité).
\end{savaistu}

\courssub{Exercices d'application (Expression orale guidée)}

\begin{exercice}[Préparation d'un mini-débat structuré (5 min écrit / 3 min oral)]
\textbf{Sujet :} \zh{在喀麦隆，预防疟疾靠蚊帐还是靠药？} (Au Cameroun, la prévention du paludisme repose-t-elle sur la moustiquaire ou les médicaments ?)

\textbf{Étape 1 (Écrit - 5 min) :} Complétez cette fiche de préparation (mots-clés en caractères/pinyin).
\begin{itemize}
    \item Mon avis principal : \zh{我觉得...因为...}
    \item Argument 1 : \zh{第一，...}
    \item Argument 2 : \zh{第二，...}
    \item Contre-argument anticipé : \zh{有人说...但是我认为...}
    \item Ma conclusion / Proposition : \zh{总之，我们应该...}
\end{itemize}

\textbf{Étape 2 (Oral - Binôme) :} Débattez sans lire vos notes. L'enseignant évalue : \textbf{Tons (nüèjí, yùfáng, dé, zhàng)}, \textbf{Structures (因为/所以, 应该, 同意/不同意)}, \textbf{Fluidité (connecteurs : 另外, 比如说, 总之)}.
\end{exercice}

\begin{corrige}[Exercice : Mini-débat - Éléments de réponse attendus]
\textbf{Pro-Moustiquaire (Exemple) :}
\zh{我觉得蚊帐最好。因为第一，政府免费发，不花钱。第二，没有副作用，孕妇和小孩都能用。有人说热，但是现在的蚊帐很透气。总之，我们应该每晚睡蚊帐。}
\emph{Wǒ juéde wénzhàng zuì hǎo. Yīnwèi dì yī, zhèngfǔ miǎnfèi fā, bù huā qián. Dì èr, méiyǒu fùzuòyòng, yùnfù xiǎohái dōu néng yòng. Shuō de rén rè, dànshì xiànzài de wénzhàng hěn tòuqì. Zǒngzhī, wǒmen yīnggāi měi wǎn shuì wénzhàng.}

\textbf{Pro-Médicaments/Complément (Exemple) :}
\zh{我不太同意。蚊帐有用，但是不够。因为蚊子白天也咬人，比如登革热蚊子。而且蚊帐破了就没用。所以我们应该吃预防药，或者用驱蚊液。两个方法一起用最安全。}
\emph{Wǒ bù tài tóngyì. Wénzhàng yǒuyòng, dànshì bù gòu. Yīnwèi wénzi báitiān yě yǎo rén, bǐrúshuō dēngguèrè wénzi. Érqiě wénzhàng pò le jiù méiyòng. Suǒyǐ wǒmen yīnggāi chī yùfáng yào, huòzhě yòng qūwényè. Liǎng gè fāngfǎ yīqǐ yòng zuì ānquán.}
\end{corrige}

\coursec{Formules fonctionnelles de l'exposé et du débat}

Après avoir analysé les deux dialogues modèles — l'exposé structuré sur le recyclage et le débat interactif sur le paludisme —, nous identifions les « briques » linguistiques qui permettent de construire une prise de parole continue et interactive. En chinois, comme en français, l'oral codifié (exposé, débat) repose sur des \cle{formules figées} (chunks) qui signalent la structure du discours : ouverture, argumentation, relance, réaction, conclusion. Les maîtriser libère l'attention pour le fond (vocabulaire technique, justes idées) et garantit la cohérence, critère majeur à l'examen oral du baccalauréat.

\courssub{Ouvrir un exposé : saluer, annoncer le sujet, capter l'attention}

Observons l'entrée en matière du dialogue modèle 1 : \zh{大家好！今天我要谈一谈垃圾分类和回收。} (Dàjiā hǎo! Jīntiān wǒ yào tán yi tán lājī fēnlèi hé huíshōu.) « Bonjour à tous ! Aujourd'hui, je vais parler du tri et du recyclage des déchets. » La structure est immuable : salutation collective $\rightarrow$ marqueur temporel \cle{今天} (aujourd'hui) $\rightarrow$ verbe d'intention \cle{要} + verbe d'action \cle{谈一谈} (discuter un peu / présenter) $\rightarrow$ sujet.

\begin{definition}[Formules d'ouverture standards]
Les formules suivantes couvrent 90\% des situations d'exposé scolaire (HSK 3-4). Le choix du verbe après \zh{要} ou \cle{想} nuance le registre : \zh{谈一谈} (neutre, standard), \zh{介绍一下} (plus formel, « présenter »), \zh{分享一下} (partager une expérience/avis).
\end{definition}

\begin{center}
\begin{tabularx}{\linewidth}{|X|X|X|X|}
\hline
\rowcolor{popblueL} Caractères & Pinyin & Nature & Traduction \\
\hline
大家好 & Dàjiā hǎo & Locution & Bonjour à tous \\
各位同学好 & Gèwèi tóngxué hǎo & Locution & Bonjour chers camarades (plus formel) \\
今天我要谈一谈… & Jīntiān wǒ yào tán yi tán… & Phrase & Aujourd'hui, je vais parler de… \\
今天我想介绍一下… & Jīntiān wǒ xiǎng jièshào yí xià… & Phrase & Aujourd'hui, je voudrais présenter… \\
我想说的是… & Wǒ xiǎng shuō de shì… & Phrase & Ce que je veux dire, c'est… / Je veux parler de… \\
这个话题很重要。 & Zhège huàtí hěn zhòngyào. & Phrase & Ce sujet est important. \\
\hline
\end{tabularx}
\end{center}

\begin{exemplebox}[Variantes d'ouverture contextuelles]
\begin{itemize}
\item \zh{各位老师、同学们好！今天我想分享一下我对预防疟疾的看法。} (Gèwèi lǎoshī, tóngxuémen hǎo! Jīntiān wǒ xiǎng fēnxiǎng yí xià wǒ duì yùfáng nuèjí de kànfǎ.) « Bonjour professeurs, camarades ! Aujourd'hui, je voudrais partager mon point de vue sur la prévention du paludisme. » (Contexte formel, jury).
\item \zh{大家好！这个话题跟我们的健康息息相关。} (Dàjiā hǎo! Zhège huàtí gēn wǒmen de jiànkāng xīxīxiāngguān.) « Bonjour à tous ! Ce sujet est étroitement lié à notre santé. » (Accroche par l'intérêt personnel).
\end{itemize}
\end{exemplebox}

\courssub{Donner son avis et argumenter : les trois piliers de l'argumentation}

Dans le dialogue modèle 2 (débat paludisme), les intervenants utilisent trois marqueurs d'opinion principaux. Observons leurs nuances :

\begin{propriete}[Trois marqueurs d'opinion : 觉得 vs 认为 vs 在我看来]
\begin{itemize}
\item \zh{我觉得…} (Wǒ juéde…) « Je trouve que… / J'ai l'impression que… » \textbf{Registre :} courant, subjectif, basé sur le ressenti. \textbf{Idéal pour démarrer} car il engage peu la responsabilité factuelle.
\item \zh{我认为…} (Wǒ rènwéi…) « J'estime que… / Je considère que… » \textbf{Registre :} plus formel, implique une réflexion, un jugement raisonné. À privilégier pour l'argument central d'un exposé.
\item \zh{在我看来…} (Zài wǒ kànlái…) « À mon avis… / Selon moi… » \textbf{Registre :} marque explicitement la subjectivité, souvent en tête de phrase pour contraster avec d'autres avis.
\end{itemize}
\end{propriete}

L'argumentation exige une justification. Le chinois standard utilise deux structures majeures :

\begin{propriete}[Structures de justification : 因为…所以… et 第一…第二…]
\begin{enumerate}
\item \textbf{Causalité explicite :} \zh{因为} (yīnwèi) [Cause] \zh{，所以} (suǒyǐ) [Conséquence].
  \emph{Attention :} En chinois, les deux termes sont \textbf{obligatoires} dans la phrase complexe écrite/formelle (contrairement au français où « parce que » suffit souvent seul). À l'oral rapide, \zh{所以} peut être omis, mais pas \zh{因为}.
\item \textbf{Énumération ordonnée :} \zh{第一} (Dìyī) [Argument 1] \zh{；} \zh{第二} (Dì'èr) [Argument 2] \zh{。} (Possibilité d'ajouter \zh{第三}, \zh{最后}).
  \emph{Structure :} \zh{第一，} + Phrase complète \zh{；第二，} + Phrase complète.
\end{enumerate}
\end{propriete}

\begin{exemplebox}[Argumenter pour le recyclage (modèle réutilisable)]
\begin{itemize}
\item \zh{第一，回收保护环境；第二，回收减少垃圾。}
Dìyī, huíshōu bǎohù huánjìng; dì'èr, huíshōu jiǎnshǎo lājī.
Premièrement, le recyclage protège l'environnement ; deuxièmement, il réduit les déchets.
\item \zh{我认为我们应该分类垃圾，因为垃圾分类可以保护环境，所以很重要。}
Wǒ rènwéi wǒmen yīnggāi fēnlèi lājī, yīnwèi lājī fēnlèi kěyǐ bǎohù huánjìng, suǒyǐ hěn zhòngyào.
J'estime qu'on devrait trier les déchets, car le tri protège l'environnement, donc c'est important.
\end{itemize}
\end{exemplebox}

\begin{methode}[Stratégie « Filet de sécurité » pour l'examen oral]
\begin{enumerate}
\item \textbf{Démarrez par} \zh{我觉得…} (simple, naturel, gagne du temps de réflexion).
\item \textbf{Enchaînez immédiatement sur} \zh{因为…所以…} pour justifier : \zh{我觉得回收好，因为保护环境，所以重要。}
\item \textbf{Si vous avez deux idées claires}, basculez sur \zh{第一…第二…} : \zh{第一，保护环境；第二，省资源。}
\item \textbf{Évitez} les phrases trop complexes (subordonnées imbriquées) : l'oral note la fluidité et la correction tonale, pas la complexité syntaxique.
\end{enumerate}
\end{methode}

\begin{center}
\begin{tabularx}{\linewidth}{|X|X|X|X|}
\hline
\rowcolor{popgreenL} Caractères & Pinyin & Nature & Traduction \\
\hline
觉得 & juéde & Verbe & Trouver, avoir l'impression que \\
认为 & rènwéi & Verbe & Estimer, considérer que \\
在我看来 & zài wǒ kànlái & Locution & À mon avis, selon moi \\
因为 & yīnwèi & Conjonction & Parce que \\
所以 & suǒyǐ & Adverbe & Donc, par conséquent \\
第一 & dìyī & Numéral/Adv & Premièrement \\
第二 & dì'èr & Numéral/Adv & Deuxièmement \\
最后 & zuìhòu & Adverbe & Enfin, en dernier lieu \\
\hline
\end{tabularx}
\end{center}

\courssub{Relancer et interagir : la particule 呢 et les questions ouvertes}

Un débat n'est pas une suite de monologues. Le dialogue modèle 2 montre l'usage constant de \zh{你觉得呢？} pour passer la parole. En français, nous disons « Et toi ? », « Qu'en penses-tu ? ». En chinois, l'omission de \cle{呢} (ne) rend la question brutale ou incomplète sur le plan interactionnel.

\begin{propriete}[La particule interactionnelle 呢 (ne) en fin de question]
\zh{呢} signale une \textbf{relance attendue} ou un \textbf{retour symétrique}. Elle transforme une question factuelle en invitation à parler.
\begin{itemize}
\item \zh{你觉得呢？} (Nǐ juéde ne?) « Et toi, qu'en penses-tu ? » (Relance standard).
\item \zh{你有什么看法？} (Nǐ yǒu shénme kànfǎ?) « Quel est ton avis ? » (Plus formel, « regard/opinion »).
\item \zh{对吗？} (Duì ma?) « C'est juste, non ? / Tu es d'accord ? » (Vérification de consensus).
\item \zh{对吧？} (Duì ba?) « Hein ? / N'est-ce pas ? » (Très oral, recherche d'approbation).
\end{itemize}
\textbf{Règle d'or :} Dans un débat, \textbf{terminez toujours votre tour de parole par une relance} (\zh{你觉得呢？} ou \zh{对吧？}) pour ne pas « tuer » l'échange.
\end{propriete}

\begin{exemplebox}[Relances contextuelles]
\begin{itemize}
\item \zh{我觉得蚊帐很管用，你觉得呢？} (Wǒ juéde wénzhàng hěn guǎnyòng, nǐ juéde ne?) « Je trouve la moustiquaire très efficace, et toi ? »
\item \zh{清理积水能预防疟疾，对吧？} (Qīnglǐ jīshuǐ néng yùfáng nuèjí, duì ba?) « Éliminer l'eau stagnante prévient le paludisme, n'est-ce pas ? »
\end{itemize}
\end{exemplebox}

\courssub{Réagir : exprimer l'accord, le désaccord nuancé, valider}

Le dialogue modèle 2 montre des réactions courtes, percutantes. Le chinois oral privilégie les verbes d'état/opinion sans copule \zh{是} (shì).

\begin{attention}[Piège majeur : \textbf{我是同意} $\times$ $\rightarrow$ \textbf{我同意} $\checkmark$]
En français, « Je \textbf{suis} d'accord » utilise le verbe être + adjectif. En chinois, \zh{同意} (tóngyì) est un \textbf{verbe transitif} signifiant « approuver, être d'accord avec ». On ne met \textbf{jamais} \zh{是} devant.
\begin{itemize}
\item \textbf{Correct :} \zh{我同意。} (Wǒ tóngyì.) / \zh{我不同意。} (Wǒ bù tóngyì.)
\item \textbf{Correct (oral) :} \zh{你说得对。} (Nǐ shuō de duì.) « Tu as raison / Ce que tu dis est juste. »
\item \textbf{Correct (enthousiaste) :} \zh{这是一个好主意！} (Zhè shì yí ge hǎo zhǔyi.) « C'est une bonne idée ! »
\end{itemize}
\end{attention}

\begin{definition}[Échelle de réaction (du consensus au désaccord poli)]
\begin{itemize}
\item \zh{我完全同意。} (Wǒ wánquán tóngyì.) « Je suis tout à fait d'accord. »
\item \zh{我赞成。} (Wǒ zànchéng.) « J'approuve / Je suis pour. » (Registre formel, vote, réunion).
\item \zh{有道理。} (Yǒu dàoli.) « Il y a du vrai / C'est sensé. » (Validation partielle, très courant).
\item \zh{我不太同意。} (Wǒ bù tài tóngyì.) « Je ne suis pas trop d'accord. » (Atténuation par \zh{不太}).
\item \zh{我有不同的看法。} (Wǒ yǒu bùtóng de kànfǎ.) « J'ai un avis différent. » (Ouverture polie au désaccord).
\end{itemize}
\end{definition}

\courssub{Conclure : résumer, appeler à l'action, remercier}

La conclusion d'un exposé scolaire suit un rituel strict : \textbf{marqueur de synthèse} $\rightarrow$ \textbf{phrase-choc / recommandation} $\rightarrow$ \textbf{remerciement formulaire}.

\begin{propriete}[Marqueurs de conclusion : 总之 vs 最后]
\begin{itemize}
\item \zh{总之} (Zǒngzhī) « En résumé / En un mot / Bref ». Introduit la \textbf{synthèse logique} de l'argumentation. Place : en tête de phrase.
\item \zh{最后} (Zuìhòu) « Enfin / Pour finir ». Introduit le \textbf{dernier point} d'une énumération ou un \textbf{appel à l'action} final.
\end{itemize}
\end{propriete}

\begin{textebox}[Modèle de conclusion type (Exposé Paludisme)]
\zh{总之，预防疟疾靠大家。我们要用蚊帐，清理积水，及时看病。}
Zǒngzhī, yùfáng nuèjí kào dàjiā. Wǒmen yào yòng wénzhàng, qīnglǐ jīshuǐ, jíshí kànbìng.
\zh{最后，谢谢大家的聆听！}
Zuìhòu, xièxie dàjiā de tíngtīng!

« En résumé, la prévention du paludisme repose sur tous. Nous devons utiliser des moustiquaires, éliminer l'eau stagnante, consulter à temps. Enfin, merci à tous pour votre écoute ! »
\end{textebox}

\begin{center}
\begin{tabularx}{\linewidth}{|X|X|X|X|}
\hline
\rowcolor{poporangeL} Caractères & Pinyin & Nature & Traduction \\
\hline
总之 & zǒngzhī & Adverbe & En résumé, bref \\
最后 & zuìhòu & Adverbe & Enfin, pour finir \\
谢谢大家 & xièxie dàjiā & Locution & Merci à tous \\
谢谢大家的聆听 & xièxie dàjiā de tíngtīng & Locution & Merci pour votre écoute (formel) \\
我们应该… & wǒmen yīnggāi… & Structure & Nous devons / Il faut que nous… \\
呼吁大家… & hūyù dàjiā… & Verbe + Obj & Lancer un appel à tous pour… \\
\hline
\end{tabularx}
\end{center}

\courssub{Tableau récapitulatif et carte mentale des fonctions langagières}

Le tableau ci-dessous synthétise les formules par fonction communicative. La carte mentale (figure TikZ) visualise l'enchaînement logique d'une prise de parole complète.

\begin{center}
\begin{tabularx}{\linewidth}{|>{\bfseries}X|X|}
\hline
\rowcolor{poppurpleL} \textbf{Fonction} & \textbf{Formules clés (Chinois + Pinyin)} \\
\hline
\textbf{Ouvrir} & \zh{大家好！} (Dàjiā hǎo!) \\
& \zh{今天我要谈一谈…} (Jīntiān wǒ yào tán yi tán…) \\
& \zh{我想说的是…} (Wǒ xiǎng shuō de shì…) \\
\hline
\textbf{Donner son avis} & \zh{我觉得…} (Wǒ juéde…) \\
& \zh{我认为…} (Wǒ rènwéi…) \\
& \zh{在我看来…} (Zài wǒ kànlái…) \\
\hline
\textbf{Justifier / Argumenter} & \zh{因为…所以…} (Yīnwèi… suǒyǐ…) \\
& \zh{第一…，第二…} (Dìyī…, dì'èr…) \\
& \zh{另外…} (Lìngwài…) « De plus / En outre » \\
\hline
\textbf{Relancer} & \zh{你觉得呢？} (Nǐ juéde ne?) \\
& \zh{你有什么看法？} (Nǐ yǒu shénme kànfǎ?) \\
& \zh{对吗？ / 对吧？} (Duì ma? / Duì ba?) \\
\hline
\textbf{Réagir (Accord)} & \zh{我同意。} (Wǒ tóngyì.) \\
& \zh{你说得对。} (Nǐ shuō de duì.) \\
& \zh{这是一个好主意。} (Zhè shì yí ge hǎo zhǔyi.) \\
\hline
\textbf{Réagir (Désaccord)} & \zh{我不同意。} (Wǒ bù tóngyì.) \\
& \zh{我不太同意。} (Wǒ bù tài tóngyì.) \\
& \zh{我有不同的看法。} (Wǒ yǒu bùtóng de kànfǎ.) \\
\hline
\textbf{Conclure} & \zh{总之…} (Zǒngzhī…) \\
& \zh{最后…} (Zuìhòu…) \\
& \zh{谢谢大家！} (Xièxie dàjiā!) \\
\hline
\end{tabularx}
\end{center}

\begin{popfigure}
\begin{center}\tcbox[colback=popink,colframe=popink,coltext=white,arc=9pt,boxsep=4pt,left=6pt,right=6pt]{\bfseries\small Exposé / Débat}\end{center}
\vspace{-2pt}
\begin{tcbraster}[raster columns=3,raster equal height=rows,raster column skip=6pt,raster row skip=6pt,enhanced,blankest]
\begin{tcolorbox}[enhanced,colback=white,colframe=poporange,boxrule=0.9pt,arc=3pt,title={1. Ouvrir},fonttitle=\bfseries\scriptsize,coltitle=white,colbacktitle=poporange,left=4pt,right=4pt,top=2pt,bottom=2pt,boxsep=1pt,toptitle=1.5pt,bottomtitle=1.5pt,halign=left]
\begin{itemize}[nosep,leftmargin=*,itemsep=1pt,font=\scriptsize]
\item 大家好
\item 今天我要谈一谈…
\item 我想说的是…
\end{itemize}
\end{tcolorbox}
\begin{tcolorbox}[enhanced,colback=white,colframe=popgreen,boxrule=0.9pt,arc=3pt,title={2. Argumenter},fonttitle=\bfseries\scriptsize,coltitle=white,colbacktitle=popgreen,left=4pt,right=4pt,top=2pt,bottom=2pt,boxsep=1pt,toptitle=1.5pt,bottomtitle=1.5pt,halign=left]
\begin{itemize}[nosep,leftmargin=*,itemsep=1pt,font=\scriptsize]
\item 我觉得 / 我认为
\item 因为…所以…
\item 第一… 第二…
\end{itemize}
\end{tcolorbox}
\begin{tcolorbox}[enhanced,colback=white,colframe=poppurple,boxrule=0.9pt,arc=3pt,title={3. Relancer},fonttitle=\bfseries\scriptsize,coltitle=white,colbacktitle=poppurple,left=4pt,right=4pt,top=2pt,bottom=2pt,boxsep=1pt,toptitle=1.5pt,bottomtitle=1.5pt,halign=left]
\begin{itemize}[nosep,leftmargin=*,itemsep=1pt,font=\scriptsize]
\item 你觉得呢？
\item 你有什么看法？
\item 对吗？
\end{itemize}
\end{tcolorbox}
\begin{tcolorbox}[enhanced,colback=white,colframe=poppink,boxrule=0.9pt,arc=3pt,title={4. Réagir},fonttitle=\bfseries\scriptsize,coltitle=white,colbacktitle=poppink,left=4pt,right=4pt,top=2pt,bottom=2pt,boxsep=1pt,toptitle=1.5pt,bottomtitle=1.5pt,halign=left]
\begin{itemize}[nosep,leftmargin=*,itemsep=1pt,font=\scriptsize]
\item 我同意 / 你说得对
\item 我不同意 / 有不同看法
\item 好主意
\end{itemize}
\end{tcolorbox}
\begin{tcolorbox}[enhanced,colback=white,colframe=popgold,boxrule=0.9pt,arc=3pt,title={5. Conclure},fonttitle=\bfseries\scriptsize,coltitle=white,colbacktitle=popgold,left=4pt,right=4pt,top=2pt,bottom=2pt,boxsep=1pt,toptitle=1.5pt,bottomtitle=1.5pt,halign=left]
\begin{itemize}[nosep,leftmargin=*,itemsep=1pt,font=\scriptsize]
\item 总之…
\item 最后…
\item 谢谢大家！
\end{itemize}
\end{tcolorbox}
\end{tcbraster}

\legende{Carte mentale : structure fonctionnelle d'un exposé/débat oral (HSK 3-4).}
\end{popfigure}

\courssub{Exercice résolu : Tri des formules par fonction communicative}

\begin{exoresolu}[Classement de cartes « Exposé / Débat »]
\textbf{Énoncé :} Classez les 12 cartes ci-dessous dans le tableau à 5 colonnes selon leur fonction : \textbf{Ouvrir}, \textbf{Argumenter}, \textbf{Relancer}, \textbf{Réagir}, \textbf{Conclure}.

\begin{center}
\begin{tabularx}{\linewidth}{|X|X|X|X|X|}
\hline
\rowcolor{popblueL} \textbf{Ouvrir} & \textbf{Argumenter} & \textbf{Relancer} & \textbf{Réagir} & \textbf{Conclure} \\
\hline
\zh{大家好！} & \zh{我认为…} & \zh{你觉得呢？} & \zh{我同意。} & \zh{总之…} \\
\zh{今天我想介绍一下…} & \zh{因为…所以…} & \zh{你有什么看法？} & \zh{你说得对。} & \zh{谢谢大家！} \\
\zh{这个话题很重要。} & \zh{第一…第二…} & \zh{对吗？} & \zh{我有不同的看法。} & \zh{最后…} \\
\hline
\end{tabularx}
\end{center}

\tcblower

\textbf{Solution détaillée :}
\begin{enumerate}
\item \textbf{Ouvrir} : Formules de contact et d'annonce de sujet. \zh{大家好！} (salutation), \zh{今天我想介绍一下…} (annonce sujet), \zh{这个话题很重要。} (accroche).
\item \textbf{Argumenter} : Expression d'opinion + justification structurée. \zh{我认为…} (opinion formelle), \zh{因为…所以…} (causalité), \zh{第一…第二…} (énumération).
\item \textbf{Relancer} : Questions interactionnelles avec \zh{呢} ou ouvertes. \zh{你觉得呢？}, \zh{你有什么看法？}, \zh{对吗？}.
\item \textbf{Réagir} : Prise de position par rapport au locuteur précédent. \zh{我同意。} (accord), \zh{你说得对。} (validation), \zh{我有不同的看法。} (désaccord poli).
\item \textbf{Conclure} : Marqueurs de clôture + politesse de fin. \zh{总之…} (synthèse), \zh{最后…} (dernier point/appel), \zh{谢谢大家！} (congé).
\end{enumerate}
\end{exoresolu}

\courssub{Exercice d'application : Construction d'un mini-débat guidé}

\begin{exercice}[Préparation orale : « Le recyclage au lycée »]
\textbf{Consigne :} En binôme, préparez un échange de 2 minutes sur le thème : \zh{学校里应该推广垃圾分类吗？} (L'école devrait-elle généraliser le tri des déchets ?).
Utilisez \textbf{au minimum} : 1 formule d'ouverture, 2 arguments avec \zh{第一/第二} ou \zh{因为/所以}, 2 relances avec \zh{呢}, 1 réaction (accord ou désaccord nuancé), 1 conclusion avec \zh{总之}.

\textbf{Rôles :}
\begin{itemize}
\item \textbf{Élève A :} Lance le sujet, donne 2 arguments \textbf{POUR}.
\item \textbf{Élève B :} Relance, donne 1 argument \textbf{POUR} + 1 argument \textbf{CONTRE} (difficulté pratique), réagit aux arguments de A, conclut.
\end{itemize}

\textbf{Aide-mémoire vocabulaire :}
\zh{推广} (tuīguǎng, v. promouvoir/généraliser), \zh{垃圾桶} (lājītǒng, n. poubelle), \zh{麻烦} (máfan, adj. embêtant/compliqué), \zh{习惯} (xíguàn, n/v. habitude/s'habituer), \zh{长远} (chángyuǎn, adj. à long terme).
\end{exercice}

\begin{corrige}[Exercice : Mini-débat « Recyclage au lycée »]
\textbf{Proposition de script (Élève A / Élève B) :}

\textbf{A :} \zh{大家好！今天我想谈一谈学校垃圾分类。我觉得应该推广，因为保护环境，所以很重要。第一，分类减少污染；第二，变废为宝。你觉得呢？}
(Dàjiā hǎo! Jīntiān wǒ xiǎng tán yi tán xuéxiào lājī fēnlèi. Wǒ juéde yīnggāi tuīguǎng, yīnwèi bǎohù huánjìng, suǒyǐ hěn zhòngyào. Dìyī, fēnlèi jiǎnshǎo wūrǎn; dì'èr, biànfèiwéibǎo. Nǐ juéde ne?)

\textbf{B :} \zh{你有什么看法？我同意分类好，但是有点麻烦。因为垃圾桶不够，学生不习惯。对吗？}
(Nǐ yǒu shénme kànfǎ? Wǒ tóngyì fēnlèi hǎo, dànshì yǒudiǎn máfan. Yīnwèi lājītǒng bùgòu, xuésheng bù xíguàn. Duì ma?)

\textbf{A :} \zh{你说得对，现在有困难。但从长远看，习惯养成了就不麻烦了。这是一个好主意，我们要坚持。}
(Nǐ shuō de duì, xiànzài yǒu kùnnán. Dàn cóng chángyuǎn kàn, xíguàn yǎngchéng le jiù bù máfan le. Zhè shì yí ge hǎo zhǔyi, wǒmen yào jiānchí.)

\textbf{B :} \zh{总之，推广分类靠大家。最后，谢谢大家！}
(Zǒngzhī, tuīguǎng fēnlèi kào dàjiā. Zuìhòu, xièxie dàjiā!)

\textbf{Analyse pédagogique :} L'échange respecte la structure cible. Élève B utilise \zh{但是} (dànshì) pour la concession (argument contre), puis \zh{从长远看} (point de vue temporel) pour rebondir. La conclusion de B synthétise (\zh{总之}) et clôture (\zh{最后}).
\end{corrige}

\coursec{Travail des tons et de la prononciation}

\courssub{Transition depuis les formules fonctionnelles}

Après avoir acquis les \cle{formules fonctionnelles} pour structurer un exposé (\zh{开头}, \zh{正文}, \zh{结尾}) et animer un débat (\zh{表示赞成/反对}, \zh{补充观点}), il est indispensable de soigner la \cle{prononciation}. En chinois, un même enchaînement de sons (pinyin) peut avoir des sens radicalement différents selon les tons. Lors de l'examen oral, la justesse tonale est un critère majeur d'évaluation~: elle conditionne l'intelligibilité de votre message sur des thèmes techniques comme le recyclage (\zh{回收}) ou le paludisme (\zh{疟疾}). Cette section vous entraîne à discriminer, produire et corriger les tons les plus fréquents dans le vocabulaire de la leçon.

\courssub{Paires minimales à risque : discrimination auditive et articulatoire}

Commençons par l'observation de paires de mots qui ne diffèrent que par un seul ton ou une seule finale. L'écoute attentive et la répétition en miroir (écoute $\rightarrow$ répétition $\rightarrow$ auto-correction) sont la clé.

\begin{textebox}[Paires minimales du thème : Recyclage \& Paludisme]
\textbf{1. 情 vs 蚊 (même voyelle, tons différents)} \\
\zh{蚊子} (wénzi) --- moustique \quad vs \quad \zh{文字} (wénzì) --- caractère, écriture \\
\zh{问} (wèn) --- demander \quad vs \quad \zh{蚊} (wén) --- moustique (racine) \\

\textbf{2. 瓶 vs 病 (même consonne initiale, voyelles et tons différents)} \\
\zh{瓶子} (píngzi) --- bouteille \quad vs \quad \zh{病} (bìng) --- maladie \\
\emph{Attention :} \zh{病子} (bìngzi) n'est pas un mot standard (on dit \zh{病人} bìngrén pour « patient »). L'élève francophone confond souvent la nasale antérieure \emph{-in} [in] et la nasale postérieure \emph{-ing} [iŋ]. \\

\textbf{3. 收 vs 手 (même voyelle, tons différents)} \\
\zh{回收} (huíshōu) --- recycler \quad vs \zh{手} (shǒu) --- main \\
\zh{收拾} (shōushi) --- ranger \quad vs \zh{手术} (shǒushù) --- opération chirurgicale \\

\textbf{4. 防 vs 放 (même voyelle, tons différents)} \\
\zh{预防} (yùfáng) --- prévenir \quad vs \zh{放} (fàng) --- mettre, poser \\
\zh{防止} (fángzhǐ) --- empêcher \quad vs \zh{放心} (fàngxīn) --- être tranquille
\end{textebox}

\begin{propriete}[Rappel : Les quatre tons et le ton neutre]
Le chinois standard (putonghua) utilise quatre tons pleins et un ton léger (neutre). Chaque ton est un \cle{contour mélodique} sur la syllabe :
\begin{itemize}
    \item \textbf{1er ton (阴平 yīnpíng) :} Haut et plat (ˉ) --- \zh{mā} (maman), \zh{lājī} (déchet), \zh{fāshāo} (fièvre).
    \item \textbf{2e ton (阳平 yángpíng) :} Montant (ˊ) --- \zh{má} (chanvre), \zh{huíshōu} (recycler - 1re syllabe), \zh{yùfáng} (prévenir - 2e syllabe).
    \item \textbf{3e ton (上声 shǎngshēng) :} Descendant-remontant (ˇ) --- \zh{mǎ} (cheval), \zh{nüèjí} (paludisme - 1re syllabe), \zh{wǒ} (je).
    \item \textbf{4e ton (去声 qùshēng) :} Chutant brutalement (ˋ) --- \zh{mà} (gronder), \zh{nüèjí} (2e syllabe citation), \zh{bù} (non).
    \item \textbf{Ton neutre (轻声 qīngshēng) :} Court, sans contour propre, dépend du ton précédent --- \zh{ma} (particule), \zh{zi} (suffixe \zh{wénzi}, \zh{píngzi}).
\end{itemize}
\end{propriete}

\begin{popfigure}
\begin{tikzpicture}[scale=0.9, baseline=(current bounding box.center), every node/.style={font=\small}]
% Premier mot : nüèjí
\begin{scope}
% Axes
\draw[popline, thick, ->] (0,0) -- (5.5,0) node[right] {\small Temps};
\draw[popline, thick, ->] (0,0) -- (0,4.5) node[above] {\small Hauteur};
% Repères de hauteur
\foreach \y/\lab in {1/1, 2/2, 3/3, 4/4, 5/5} {
    \draw[popmuted, thin] (0,\y) -- (5.5,\y);
    \node[left, font=\tiny, popdark] at (0,\y) {\lab};
}
% Courbe 4e ton : nüè (5 -> 1)
\draw[popblue, very thick] (0.5,5) -- (1.5,1);
\node[popblue, above, font=\small\bfseries] at (1,5.3) {nüè (4e ton)};
% Courbe 1er ton : ji (5 -> 5) ou neutre (bas)
\draw[poporange, very thick] (2,4.5) -- (3,4.5);
\node[poporange, above, font=\small\bfseries] at (2.5,4.8) {ji (1er/neutre)};
% Etiquette mot
\node[align=center, font=\small, popdark] at (1.5,-0.6) {\zh{疟疾} \\ nüèjí};
\end{scope}

% Deuxième mot : huíshōu
\begin{scope}[xshift=7.5cm]
\draw[popline, thick, ->] (0,0) -- (5.5,0) node[right] {\small Temps};
\draw[popline, thick, ->] (0,0) -- (0,4.5) node[above] {\small Hauteur};
\foreach \y/\lab in {1/1, 2/2, 3/3, 4/4, 5/5} {
    \draw[popmuted, thin] (0,\y) -- (5.5,\y);
    \node[left, font=\tiny, popdark] at (0,\y) {\lab};
}
% Courbe 2e ton : huí (3 -> 5)
\draw[popgreen, very thick] (0.5,2) .. controls (1,3.5) .. (1.5,5);
\node[popgreen, above, font=\small\bfseries] at (1,5.3) {huí (2e ton)};
% Courbe 1er ton : shōu (5 -> 5)
\draw[poppurple, very thick] (2,4.5) -- (3,4.5);
\node[poppurple, above, font=\small\bfseries] at (2.5,4.8) {shōu (1er ton)};
% Etiquette mot
\node[align=center, font=\small, popdark] at (1.5,-0.6) {\zh{回收} \\ huíshōu};
\end{scope}
\end{tikzpicture}
\legende{Contours toniques de \zh{疟疾} (nüèjí : 4e + 1er ton citation / neutre à l'oral) et \zh{回收} (huíshōu : 2e + 1er ton). Le 4e ton chute de 5 à 1 ; le 2e ton monte de 3 à 5 ; le 1er ton reste haut (5).}
\end{popfigure}

\courssub{Tons des mots-clés du module : ancrage lexical}

Le tableau ci-dessous récapitule le vocabulaire technique de la leçon. Mémorisez le \textbf{contour tonique du mot entier} (mot = unité rythmique), pas seulement syllabe par syllabe.

\begin{center}
\begin{tabularx}{\linewidth}{|X|X|X|X|X|}
\hline
\rowcolor{popblueL}
\textbf{Caractères} & \textbf{Pinyin} & \textbf{Structure tonale} & \textbf{Nature} & \textbf{Traduction} \\
\hline
疟疾 & nüèjí & 4 + 1 (citation) / 4 + 0 (oral) & nom & paludisme \\
\hline
回收 & huíshōu & 2 + 1 & verbe & recycler, récupération \\
\hline
垃圾 & lājī & 1 + 1 & nom & déchets, ordures \\
\hline
预防 & yùfáng & 4 + 2 & verbe & prévenir \\
\hline
发烧 & fāshāo & 1 + 1 & verbe & avoir de la fièvre \\
\hline
蚊子 & wénzi & 2 + 0 & nom & moustique \\
\hline
瓶子 & píngzi & 2 + 0 & nom & bouteille \\
\hline
症状 & zhèngzhuàng & 4 + 4 & nom & symptôme \\
\hline
治疗 & zhìliáo & 4 + 2 & verbe & traiter, soigner \\
\hline
环境 & huánjìng & 2 + 4 & nom & environnement \\
\hline
\end{tabularx}
\end{center}

\begin{attention}[Note importante sur \zh{疟疾} nüèjí]
Dictionnaire standard : \zh{疟} nüè (4e ton), \zh{疾} jí (2e ton). Cependant, \zh{疾} est souvent prononcé en \textbf{ton neutre} (jí $\rightarrow$ ji) dans le mot composé \zh{疟疾} à l'oral rapide : \zh{nüèji} (4 + 0). Les deux sont acceptés, mais \textbf{nüèji} (4 + neutre) est plus naturel en exposé oral. Entraînez-vous aux deux versions.
\end{attention}

\courssub{Le sandhi du 3e ton : le « demi-3e ton »}

C'est la règle de phonétique contextuelle la plus fréquente et la plus oubliée par les francophones.

\begin{propriete}[Règle du sandhi du 3e ton (三声变调)]
Lorsque \textbf{deux 3e tons se suivent}, le \textbf{premier} 3e ton se prononce comme un \textbf{2e ton} (montant). C'est le « demi-3e ton » (半三声) : il monte mais ne redescend pas.
\medskip
\textbf{Schéma :} 3e ton + 3e ton $\rightarrow$ \textbf{2e ton} + 3e ton.
\end{propriete}

\begin{exemplebox}[Applications au vocabulaire de la leçon]
\begin{itemize}
    \item \zh{我觉得} : wǒ (3) + juéde (3/0) $\rightarrow$ \textbf{wó juéde} (2 + 0). \emph{« Je pense que... »} (formule d'exposé).
    \item \zh{很好} : hěn (3) + hǎo (3) $\rightarrow$ \textbf{hén hǎo} (2 + 3). \emph{« Très bien. »} (réaction en débat).
    \item \zh{可以} : kě (3) + yǐ (3) $\rightarrow$ \textbf{ké yǐ} (2 + 3). \emph{« Peut / C'est possible. »}
    \item \zh{马上} : mǎ (3) + shàng (4) $\rightarrow$ \textbf{pas de changement} (mǎ shàng). Le 3e ton reste 3e ton (demi-3e ton bas) devant un ton 1, 2 ou 4.
\end{itemize}
\end{exemplebox}

\begin{attention}[Piège francophone : le « 3e ton citation » vs « 3e ton contexte »]
En français, l'intonation porte sur la phrase (montante pour la question, descendante pour l'affirmation). En chinois, \textbf{chaque syllabe a son ton propre}. Ne laissez pas l'intonation française « écraser » le 3e ton.
\begin{itemize}
    \item \textbf{Erreur :} Dire \zh{wǒ juéde} avec un 3e ton complet (descente-remontée) sur \zh{wǒ}. Cela sonne hésitant, non natif.
    \item \textbf{Correct :} \zh{wó juéde} (montée nette sur \zh{wo}, puis ton neutre léger sur \zh{de}).
\end{itemize}
\emph{Astuce :} Visualisez une flèche $\nearrow$ pour le premier 3e ton devant un autre 3e ton.
\end{attention}

\courssub{Le sandhi de \zh{不} (bù) : adaptation tonale obligatoire}

Le mot \zh{不} (négation) est un 4e ton de base, mais il change de ton selon le ton de la syllabe qui suit. C'est une règle \textbf{systématique} (pas d'exception).

\begin{propriete}[Sandhi de \zh{不} bù]
\begin{center}
\begin{tabularx}{\linewidth}{|X|X|X|X|}
\hline
\rowcolor{popgreenL}
\textbf{Ton de la syllabe suivante} & \textbf{Ton de \zh{不}} & \textbf{Pinyin résultant} & \textbf{Exemple (thème leçon)} \\
\hline
1er ton (haut) & 4e ton (inchangé) & \zh{bù} & \zh{不发烧} (bù fāshāo) --- pas de fièvre \\
\hline
2e ton (montant) & 4e ton (inchangé) & \zh{bù} & \zh{不同意} (bù tóngyì) --- ne pas être d'accord \\
\hline
3e ton (desc-mont) & \textbf{2e ton (montant)} & \zh{bú} & \zh{不是} (bú shì) --- ce n'est pas \\
\hline
\textbf{4e ton (chutant)} & \textbf{2e ton (montant)} & \zh{bú} & \zh{不对} (bú duì), \zh{不要} (bú yào) \\
\hline
Ton neutre & 4e ton (souvent) & \zh{bù} & \zh{不好} (bù hǎo) *voir note \\
\hline
\end{tabularx}
\end{center}
\emph{Note :} Devant un adjectif monosyllabique 3e ton comme \zh{好} (hǎo), \zh{不} devient \zh{bú} (2e ton) $\rightarrow$ \zh{bú hǎo}. Devant un verbe 3e ton, idem.
\end{propriete}

\begin{exemplebox}[Exemples contrastés pour l'oral]
\begin{itemize}
    \item \zh{我不发烧。} Wǒ \textbf{bù} fāshāo. (Devant 1er ton $\rightarrow$ 4e ton). \emph{Je n'ai pas de fièvre.}
    \item \zh{这不对。} Zhè \textbf{bú} duì. (Devant 4e ton $\rightarrow$ 2e ton). \emph{C'est pas juste / incorrect.}
    \item \zh{我们不要垃圾。} Wǒmen \textbf{bú} yào lājī. (Devant 4e ton $\rightarrow$ 2e ton). \emph{Nous ne voulons pas de déchets.}
    \item \zh{这不是瓶子。} Zhè \textbf{bú} shì píngzi. (Devant 4e ton \zh{shì} $\rightarrow$ 2e ton). \emph{Ceci n'est pas une bouteille.}
    \item \zh{我不同意。} Wǒ \textbf{bù} tóngyì. (Devant 2e ton \zh{tóng} $\rightarrow$ 4e ton). \emph{Je ne suis pas d'accord.}
\end{itemize}
\end{exemplebox}

\courssub{Exercice de lecture à voix haute : vérification croisée des tons}

\begin{exoresolu}[Exercice de lecture à voix haute : vérification croisée des tons]
\textbf{Consigne :} Lisez les phrases ci-dessous à haute voix. Votre camarade coche la case correspondant au ton entendu pour les mots en gras. Inversez les rôles.

\begin{center}
\begin{tabularx}{\linewidth}{|X|c|c|c|c|c|}
\hline
\rowcolor{poporangeL}
\textbf{Phrase} & \textbf{Mot cible} & \textbf{1er} & \textbf{2e} & \textbf{3e} & \textbf{4e} \\
\hline
1. \zh{我们要回收瓶子。} (Wǒmen yào \textbf{huíshōu} \textbf{píngzi}.) & huí / shōu / píng &  /  /  & ✓ /  / ✓ &  /  /  &  / ✓ /  \\
\hline
2. \zh{蚊子传播疟疾。} (\textbf{Wénzi} chuánbō \textbf{nüèjí}.) & wén / nüè / ji &  /  / ✓* & ✓ /  /  &  /  /  &  / ✓ /  \\
\hline
3. \zh{我觉得预防很重要。} (\textbf{Wǒ juéde} \textbf{yùfáng} hěn zhòngyào.) & wó / jué / yù / fáng &  /  /  /  & ✓ /  /  / ✓ &  / ✓ /  /  &  /  / ✓ /  \\
\hline
4. \zh{我不发烧，也不要药。} (Wǒ \textbf{bù fāshāo}, yě \textbf{bú yào} yào.) & bù / fā / bú / yào &  / ✓ /  /  &  /  / ✓ /  &  /  /  /  & ✓ /  /  / ✓ \\
\hline
\end{tabularx}
\end{center}
\textbf{Corrigé détaillé :}
\begin{enumerate}
    \item \zh{huí} (2e), \zh{shōu} (1er), \zh{píng} (2e). \emph{Attention à ne pas dire \zh{huǐ} (3e) ni \zh{pǐng} (3e).}
    \item \zh{wén} (2e), \zh{nüè} (4e), \zh{ji} (1er ou neutre). \emph{Piège : \zh{wèn} (4e) = demander.}
    \item \zh{wó} (2e par sandhi), \zh{jué} (2e), \zh{yù} (4e), \zh{fáng} (2e).
    \item \zh{bù} (4e devant 1er ton \zh{fā}), \zh{bú} (2e devant 4e ton \zh{yào}).
\end{enumerate}
\end{exoresolu}

\courssub{Jeu pédagogique : « Ton ou pas ton ? » (听声调，举卡片)}

\begin{experience}[Activité de classe : Discrimination tonale rapide]
\textbf{Matériel :} 5 cartes par élève (ou doigts levés) : \textbf{1} (horizontal), \textbf{2} (montant), \textbf{3} (V), \textbf{4} (descendant), \textbf{0} (point/neutre).
\textbf{Déroulement :}
\begin{enumerate}
    \item L'enseignant prononce \textbf{une seule syllabe} ou un \textbf{mot de 2 syllabes} issu de la liste vocabulaire (ex: \zh{收}, \zh{疾}, \zh{预防}, \zh{瓶子}, \zh{不发烧}).
    \item Les élèves lèvent \textbf{immédiatement} la carte (ou les doigts) correspondant au(x) ton(s) entendu(s).
    \item L'enseignant scanne la salle : feedback instantané. Si erreur majoritaire $\rightarrow$ modélisation collective (chorale) $\rightarrow$ ré-essai individuel.
\end{enumerate}
\textbf{Variante « Piège » :} L'enseignant dit un mot \emph{proche} mais différent (ex: \zh{问} wèn au lieu de \zh{蚊} wén ; \zh{放} fàng au lieu de \zh{防} fáng). Les élèves doivent lever la carte du ton \textbf{entendu}, pas du mot \textbf{attendu}. Cela force l'écoute phonétique pure.
\end{experience}

\courssub{Méthode d'auto-correction : l'enregistrement personnel}

\begin{methode}[Enregistrer, comparer, corriger (3 étapes)]
\textbf{Étape 1 -- Enregistrement :} Utilisez le magnétophone de votre téléphone. Lisez le \textbf{Dialogue Modèle 1 (Exposé recyclage)} et le \textbf{Dialogue Modèle 2 (Débat paludisme)} de la section précédente. Parlez à vitesse normale, sans vous presser.
\textbf{Étape 2 -- Comparaison ciblée :} Réécoutez votre enregistrement \emph{en même temps} que le modèle audio (fourni dans l'application OnBuch+ ou lu par l'enseignant). Arrêtez sur chaque mot-clé du tableau ci-dessus (\zh{nüèjí}, \zh{huíshōu}, \zh{lājī}, \zh{yùfáng}, \zh{fāshāo}, \zh{wénzi}, \zh{píngzi}, \zh{bù fāshāo}, \zh{bú duì}).
\textbf{Étape 3 -- Répétition ciblée (Shadowing) :} Rejouez \textbf{uniquement les phrases où vos tons divergent}. Répétez 5 fois de suite en \emph{shadowing} (parler en même temps que l'audio, avec 0,5s de décalage). Enregistrez à nouveau pour valider.
\end{methode}

\begin{savaistu}[Les tons, « empreintes digitales » du mot]
Des études de psycholinguistique montrent que pour un locuteur natif chinois, le \textbf{contour tonal} est traité par le cerveau \textbf{avant} même les phonèmes (consonnes/voyelles). Si vous dites \zh{shuǐ} (eau, 3e ton) avec un 2e ton \zh{shuí} (qui), le Chinois entend d'abord « qui ? » avant de réaliser que le contexte impose « eau ». C'est pourquoi à l'oral, \textbf{un ton faux = un mot faux}, même si les voyelles/consonnes sont parfaites. Entraînez les tons comme des notes de musique : la justesse prime sur la vitesse.
\end{savaistu}

\courssub{Exercice de synthèse : Mini-exposé « Tons maîtrisés »}

\begin{exercice}[Production orale guidée : 2 minutes]
Préparez un mini-exposé de \textbf{8 phrases exactement} intégrant \textbf{au moins 8 mots-clés} du tableau de vocabulaire (section 3) et respectant \textbf{tous les sandhi} (3e ton, \zh{不}).
\textbf{Structure imposée :}
1. Ouverture (\zh{大家好，我说说...})
2. Fait : \zh{回收} + \zh{垃圾} / \zh{瓶子}
3. Problème : \zh{蚊子} + \zh{疟疾} / \zh{发烧}
4. Solution : \zh{预防} + \zh{环境}
5. Avis personnel : \zh{我觉得} + \zh{很重要} (sandhi 3e ton !)
6. Négation : \zh{不} + verbe 4e ton (\zh{不同意} / \zh{不要}) $\rightarrow$ \zh{bú} / \zh{bù} (attention au ton du verbe !)
7. Conclusion : \zh{谢谢大家。}

\textbf{Contrainte :} Écrivez le pinyin \textbf{avec les tons réels prononcés} (après sandhi) au-dessus des caractères dans votre brouillon.
\textbf{Exemple de début :}
\zh{大家好，我说说回收 (huíshōu) 和疟疾 (nüèjí)。}
\zh{回收 (huíshōu) 瓶子 (píngzi) 很重要。}
\zh{蚊子 (wénzi) 传播 疟疾 (nüèjí)，让人 发烧 (fāshāo)。}
\zh{我们 要 (yào) 预防 (yùfáng)，改善 环境 (huánjìng)。}
\zh{我 觉得 (wó juéde) 这 很 重要 (hén zhòngyào)。}
\zh{我 不 (bú) 同意 (tóngyì) 乱扔 垃圾 (lājī)。} \emph{(Note: 同意 2e ton $\rightarrow$ bù 4e ton. Si verbe 4e ton ex: 要 $\rightarrow$ bú yào)}
\zh{谢谢 大家 (dàjiā)！}
\end{exercice}

\begin{corrige}[Exercice Mini-exposé]
\textbf{Vérification des tons (clé pour l'enseignant / auto-évaluation) :}
\begin{itemize}
    \item \zh{回收} huí(2) shōu(1) --- \zh{疟疾} nüè(4) ji(1/0)
    \item \zh{瓶子} píng(2) zi(0) --- \zh{蚊子} wén(2) zi(0)
    \item \zh{发烧} fā(1) shāo(1) --- \zh{预防} yù(4) fáng(2)
    \item \zh{环境} huán(2) jìng(4) --- \zh{觉得} wó(2) juéde(0) [Sandhi 3+3/0]
    \item \zh{重要} hén(2) zhòng(4) yào(4) [Sandhi 3+4 $\rightarrow$ 2+4] --- \zh{不同意} bù(4) tóng(2) yì(4) [Sandhi \zh{不}+2e ton] \textbf{OU} \zh{不要} bú(2) yào(4) [Sandhi \zh{不}+4e ton]
    \item \zh{垃圾} lā(1) jī(1)
\end{itemize}
\textbf{Critères de réussite (grille d'évaluation orale) :}
\begin{center}
\begin{tabularx}{\linewidth}{|X|X|}
\hline
\rowcolor{poppurpleL}
\textbf{Critère} & \textbf{Indicateur de réussite} \\
\hline
Justesse tonale mots-clés & 0 erreur sur les 8 mots imposés (sandhi inclus) \\
\hline
Fluidité / Débit & Pas d'hésitation > 2s ; liaisons naturelles \\
\hline
Structures & 8 phrases, connecteurs logiques (\zh{因为/所以/但是}) utilisés \\
\hline
Prononciation segmentale & Initiales/finales claires (ex: \zh{nüè} vs \zh{nüe}, \zh{huí} vs \zh{huì}) \\
\hline
\end{tabularx}
\end{center}
\end{corrige}

\courssub{Tableau récapitulatif : Tons, Sandhi et Pièges pour l'examen oral}

\begin{center}
\begin{tabularx}{\linewidth}{|X|X|X|}
\hline
\rowcolor{popgoldL}
\textbf{Notion} & \textbf{Règle / Forme correcte} & \textbf{Erreur fréquente (Francophone)} \\
\hline
\textbf{3e ton + 3e ton} & 1er devient 2e ton (\zh{wó juéde}) & Garder 3e ton complet (\zh{wǒ juéde}) = son « chanté », non natif \\
\hline
\textbf{\zh{不} + 4e ton} & \zh{不} $\rightarrow$ \zh{bú} (2e ton) : \zh{bú duì, bú yào} & Laisser \zh{bù} (4e ton) : *\zh{bù duì} = accent fort, peu naturel \\
\hline
\textbf{\zh{不} + 1er/2e/3e ton} & \zh{不} reste \zh{bù} (4e ton) : \zh{bù fāshāo, bù tóngyì, bú shì} & Le passer en 2e ton par sur-généralisation : *\zh{bú fāshāo}, *\zh{bú tóngyì} \\
\hline
\textbf{\zh{一} yī (nombre)} & Seul/fin : 1er ; devant 4e : 2e (\zh{yí ge}) ; devant 1/2/3 : 4e (\zh{yì tiān}) & Oublier les changements ; toujours dire \zh{yī} \\
\hline
\textbf{Neutre \zh{zi} / \zh{de}} & Court, sans accent : \zh{wénzi, píngzi, juéde} & Allonger : *\zh{wén-zǐ}, *\zh{jué-dé} (ajoute une syllabe) \\
\hline
\textbf{Nasales -in / -ing} & \zh{píng} [piŋ] vs \zh{pín} [pin] ; \zh{jīn} vs \zh{jīng} & Confondre [in] et [iŋ] $\rightarrow$ \zh{瓶子} $\neq$ \zh{品子} \\
\hline
\end{tabularx}
\end{center}

\begin{attention}[Dernier rappel avant l'examen oral]
L'examinateur note : \textbf{1. Justesse des tons (lexicaux + sandhi)}, \textbf{2. Intelligibilité globale}, \textbf{3. Richesse lexicale (vocabulaire module)}, \textbf{4. Cohérence du discours}.
Un exposé avec un vocabulaire parfait mais des tons approximatifs (ex: *\zh{bù duì} au lieu de \zh{bú duì}, *\zh{wǒ juéde} au lieu de \zh{wó juéde}) perdra des points significatifs sur le critère 1. \textbf{Révisez les 10 phrases de l'exercice résolu chaque jour pendant 5 minutes jusqu'à l'examen.}
\end{attention}

\coursec{Jeux de rôle et débat guidé}

Après avoir consolidé la prononciation des tons et les formules fonctionnelles dans la section précédente, nous passons maintenant à la mise en situation réelle. L'objectif est de transformer les connaissances passives (vocabulaire, structures, intonation) en compétences actives de communication orale. Trois activités progressives sont proposées : deux jeux de rôle codifiés (interview, consultation) et un débat guidé à argumentation contradictoire. Chaque activité s'appuie sur une grille d'observation commune pour l'auto-évaluation et l'évaluation par les pairs.

\courssub{Jeu de rôle 1 : Interview — « Le recyclage dans notre lycée »}

\begin{methode}[Préparation de l'interview (5 minutes)]
\begin{enumerate}
\item \textbf{Rôles :} Élève A = Journaliste (记者 \zh{记者} \emph{jìzhě}) ; Élève B = Éco-délégué / Expert recyclage (环保代表 \zh{环保代表} \emph{huánbǎo dàibiǎo}).
\item \textbf{Journaliste :} Prépare 5 questions en utilisant les formules d'introduction (\zh{请问...}, \zh{我想知道...}) et de relance (\zh{然后呢？}, \zh{具体来说？}).
\item \textbf{Expert :} Prépare 3 arguments clés (tri, compost, réutilisation) avec des chiffres ou exemples concrets (ex. \zh{每周回收纸张十公斤}).
\item \textbf{Contrainte :} Notes autorisées (mots-clés uniquement), pas de phrases complètes écrites. Regarder son interlocuteur, parler distinctement.
\end{enumerate}
\end{methode}

\begin{textebox}[Modèle d'interview : Le recyclage au lycée de Yaoundé]
记者：你好，我是校园电台的记者。今天我们谈谈垃圾分类。请问，我们学校现在怎么回收垃圾？\\
Jìzhě: Nǐ hǎo, wǒ shì xiàoyuán diàntái de jìzhě. Jīntiān wǒmen tán tan lājī fēnlèi. Qǐngwèn, wǒmen xuéxiào xiànzài zěnme huíshōu lājī?\\
\textit{Journaliste : Bonjour, je suis journaliste à la radio du campus. Aujourd'hui nous parlons du tri des déchets. Comment notre école recycle-t-elle les déchets actuellement ?}\\
环保代表：你好。我们学校有三个颜色的垃圾桶：蓝色装纸，绿色装厨余，灰色装其他垃圾。每个教室门口都有。\\
Huánbǎo dàibiǎo: Nǐ hǎo. Wǒmen xuéxiào yǒu sān gè yánsè de lājītǒng: lánsè zhuāng zhǐ, lǜsè zhuāng chúyú, huīsè zhuāng qítā lājī. Měi gè jiàoshì ménkǒu dōu yǒu.\\
\textit{Éco-délégué : Bonjour. Notre école a trois poubelles de couleurs : bleue pour le papier, verte pour les déchets de cuisine, grise pour les autres déchets. Chaque porte de classe en a.}\\
记者：很好。学生们都遵守规则吗？\\
Jìzhě: Hěn hǎo. Xuéshengmen dōu zūnshǒu guīzé ma?\\
\textit{Journaliste : Très bien. Les élèves respectent-ils tous les règles ?}\\
环保代表：大部分同学做得好，但还有人把塑料瓶扔进纸篓。我们要加强宣传。\\
Huánbǎo dàibiǎo: Dàbùfen tóngxué zuò de hǎo, dàn hái yǒu rén bǎ sùliáopíng rēng jìn zhǐlóu. Wǒmen yào jiāqiáng xuānchuán.\\
\textit{Éco-délégué : La plupart font bien, mais certains jettent encore les bouteilles en plastique dans la corbeille à papier. Nous devons renforcer la sensibilisation.}\\
记者：具体有什么活动吗？\\
Jìzhě: Jùtǐ yǒu shénme huódòng ma?\\
\textit{Journaliste : Y a-t-il des activités concrètes ?}\\
环保代表：有。下个月我们办“变废为宝”比赛，用废纸、塑料瓶做手工。欢迎大家参加！\\
Huánbǎo dàibiǎo: Yǒu. Xià gè yuè wǒmen bàn “biànfèiwéibǎo” bǐsài, yòng fèizhǐ, sùliáopíng zuò shǒugōng. Huānyíng dàjiā cānjiā!\\
\textit{Éco-délégué : Oui. Le mois prochain nous organisons un concours "Transformer les déchets en trésors", faire de l'artisanat avec du papier usagé et des bouteilles. Bienvenue à tous !}\\
记者：谢谢您的介绍。祝活动成功！\\
Jìzhě: Xièxie nín de jièshào. Zhù huódòng chénggōng!\\
\textit{Journaliste : Merci pour votre présentation. Bonne chance pour l'activité !}
\end{textebox}

\begin{center}
\begin{tabularx}{\linewidth}{|X|X|X|X|}
\hline \rowcolor{popblueL} Caractères & Pinyin & Nature & Traduction \\ \hline
记者 & jìzhě & n. & journaliste \\
环保代表 & huánbǎo dàibiǎo & n. & représentant éco / délégué environnement \\
垃圾分类 & lājī fēnlèi & v./n. & tri des déchets / tri sélectif \\
垃圾桶 & lājītǒng & n. & poubelle \\
厨余 & chúyú & n. & déchets de cuisine (épluchures, restes) \\
遵守规则 & zūnshǒu guīzé & v. & respecter les règles \\
塑料瓶 & sùliáopíng & n. & bouteille en plastique \\
纸篓 & zhǐlóu & n. & corbeille à papier \\
宣传 & xuānchuán & v./n. & sensibiliser / sensibilisation / propagande \\
变废为宝 & biànfèiwéibǎo & idiome & transformer les déchets en trésors (recyclage créatif) \\
手工 & shǒugōng & n. & travaux manuels, artisanat \\
\hline
\end{tabularx}
\end{center}

\begin{exemplebox}[5 questions types pour le journaliste (avec formules fonctionnelles)]
\begin{enumerate}
\item \zh{请问，学校为什么要推行垃圾分类？} (Qǐngwèn, xuéxiào wèishénme yào tuīxíng lājī fēnlèi?) — \textit{Pourquoi l'école promeut-elle le tri ?}
\item \zh{我想知道，回收的纸和塑料最后去哪里？} (Wǒ xiǎng zhīdào, huíshōu de zhǐ hé sùliào zuìhòu qù nǎlǐ?) — \textit{Je voudrais savoir où finissent le papier et le plastique recyclés.}
\item \zh{据我所知，有些同学不愿意分类。您怎么看？} (Jùwǒ zhīdào, yǒuxiē tóngxué bù yuànyì fēnlèi. Nín zěnme kàn?) — \textit{D'après moi, certains élèves ne veulent pas trier. Quel est votre avis ?}
\item \zh{然后呢？有没有惩罚措施？} (Ránhòu ne? Yǒu méiyǒu chéngfá cuòshī?) — \textit{Et ensuite ? Y a-t-il des sanctions ?}
\item \zh{最后，给同学们一个建议吧。} (Zuìhòu, gěi tóngxuémen yí gè jiànyi ba.) — \textit{Enfin, un conseil pour les camarades.}
\end{enumerate}
\end{exemplebox}

\begin{attention}[Pièges à l'oral : Interview]
\begin{itemize}
\item \textbf{Ordre des mots :} En chinois, le complément de lieu/temps précède le verbe. Ne dites pas \zh{*我们回收垃圾在学校} mais \zh{我们在学校回收垃圾} (Wǒmen zài xuéxiào huíshōu lājī).
\item \textbf{Classificateurs :} Pour « bouteille », le classifieur générique \zh{个} (gè) est correct (\zh{一个瓶子} yí gè píngzi). Pour plus de précision, on utilise \zh{只} (zhǐ) pour le contenant vide (\zh{一只瓶子}) ou \zh{瓶} (píng) pour le contenu (\zh{一瓶水}). N'oubliez jamais le classifieur après le numéral.
\item \textbf{Ton de \zh{不} (bù) :} Devant un 4ème ton (\zh{是} shì, \zh{要} yào), \zh{不} devient 2ème ton (\zh{不\textbf{是}} bú shì, \zh{不\textbf{要}} bú yào). Ex: \zh{我不\textbf{是}不愿意} (Wǒ bú shì bù yuànyì = Ce n'est pas que je ne veuille pas).
\end{itemize}
\end{attention}

\courssub{Jeu de rôle 2 : Consultation médicale — « Symptômes de paludisme »}

\begin{methode}[Préparation de la consultation (5 minutes)]
\begin{enumerate}
\item \textbf{Rôles :} Élève A = Patient (病人 \zh{病人} \emph{bìngrén}) ; Élève B = Médecin (医生 \zh{医生} \emph{yīshēng}).
\item \textbf{Patient :} Décrit 3 symptômes minimum (\zh{发烧}, \zh{头疼}, \zh{浑身疼}, \zh{恶心}, \zh{怕冷}) et la durée (\zh{两天了}).
\item \textbf{Médecin :} Pose des questions (\zh{吃药了吗？}, \zh{查过血吗？}), donne un diagnostic probable (\zh{可能是疟疾}) et 3 conseils avec \zh{应该} / \zh{得} (\zh{多喝水}, \zh{好好休息}, \zh{按时吃药}).
\item \textbf{Contexte camerounais :} Évoquer la moustiquaire (蚊帐 \zh{蚊帐} \emph{wénzhàng}) et l'hôpital de district (区医院 \zh{区医院} \emph{qū yīyuàn}).
\end{enumerate}
\end{methode}

\begin{textebox}[Modèle de consultation : Paludisme à Douala]
病人：医生，我不舒服。我发烧了，头疼，浑身疼，还有点恶心。已经两天了。\\
Bìngrén: Yīshēng, wǒ bù shūfu. Wǒ fāshāo le, tóuténg, húnshēn téng, hái yǒu diǎn ěxīn. Yǐjīng liǎng tiān le.\\
\textit{Patient : Docteur, je ne vais pas bien. J'ai de la fièvre, mal à la tête, mal partout, un peu nauséeux. Ça fait déjà deux jours.}\\
医生：量过体温吗？多少度？\\
Yīshēng: Liáng guò tǐwēn ma? Duōshao dù?\\
\textit{Médecin : Avez-vous pris la température ? Combien de degrés ?}\\
病人：三十九度五。我吃了退烧药，但没用。\\
Bìngrén: Sānshíjiǔ dù wǔ. Wǒ chī le tuìshāo yào, dàn méi yòng.\\
\textit{Patient : 39,5°C. J'ai pris un antipyrétique, mais ça n'a pas marché.}\\
医生：有没有怕冷、出汗？\\
Yīshēng: Yǒu méiyǒu pà lěng, chū hàn?\\
\textit{Médecin : Frissons, transpiration ?}\\
病人：有。白天怕冷，晚上出很多汗。\\
Bìngrén: Yǒu. Bàitiān pà lěng, wǎnshàng chū hěn duō hàn.\\
\textit{Patient : Oui. Le jour j'ai froid, la nuit je transpire beaucoup.}\\
医生：很像疟疾。你应该马上去验血。确诊后，要按时吃抗疟药。平时应该多喝水，多休息，盖蚊帐睡觉。\\
Yīshēng: Hěn xiàng nüèji. Nǐ yīnggāi mǎshàng qù yànxiě. Quèzhěn hòu, yào ànshí chī kàngnüè yào. Píngshí yīnggāi duō hē shuǐ, duō xiūxi, gài wénzhàng shuìjiào.\\
\textit{Médecin : Ça ressemble beaucoup au paludisme. Vous devriez aller faire une prise de sang tout de suite. Une fois confirmé, il faut prendre les antipaludéens à l'heure. D'habitude, il faut boire beaucoup d'eau, bien se reposer, dormir sous moustiquaire.}\\
病人：谢谢医生。我现在就去验血。\\
Bìngrén: Xièxie yīshēng. Wǒ xiànzài jiù qù yànxiě.\\
\textit{Patient : Merci docteur. J'y vais maintenant pour la prise de sang.}
\end{textebox}

\begin{propriete}[Structure \zh{应该 yīnggāi} + Verbe (Devoir / Conseil fort)]
\begin{center}
\begin{tabularx}{\linewidth}{|X|X|X|}
\hline \rowcolor{popgreenL} Structure & Exemple (Caractères + Pinyin) & Traduction \\ \hline
Suj. + 应该 + V (+ Obj.) & \zh{你应该多喝水。} (Nǐ yīnggāi duō hē shuǐ.) & Tu dois / Tu devrais boire plus d'eau. \\
Suj. + 不应该 + V & \zh{病人不应该熬夜。} (Bìngrén bù yīnggāi áoyè.) & Le patient ne devrait pas veiller tard. \\
疑问句 : Suj. + 应该 + V + 吗？ & \zh{我应该吃什么药？} (Wǒ yīnggāi chī shénme yào?) & Quel médicament dois-je prendre ? \\
\hline
\end{tabularx}
\end{center}
\textbf{Emploi :} Exprime un devoir moral, un conseil médical ou une obligation logique. Plus fort que \zh{可以} (pouvoir), moins contraignant que \zh{必须} (devoir absolu).
\textbf{Négation :} \zh{不应该} (ne pas devoir / il ne faut pas).
\textbf{Comparatif français :} « Tu devrais / Il faut / Tu dois ». Attention : \zh{应该} ne se conjugue pas.
\end{propriete}

\begin{exemplebox}[Vocabulaire clé : Symptômes et Traitement du paludisme]
\begin{center}
\begin{tabularx}{\linewidth}{|X|X|X|X|}
\hline \rowcolor{poporangeL} Caractères & Pinyin & Nature & Traduction \\ \hline
疟疾 & nüèji & n. & paludisme / malaria \\
发烧 / 发热 & fāshāo / fārè & v. & avoir de la fièvre \\
头疼 / 头痛 & tóuténg / tóutòng & v. & mal à la tête \\
浑身疼 & húnshēn téng & v. & mal au corps / courbatures \\
恶心 & ěxīn & adj. & nauséeux \\
呕吐 & ǔtù & v. & vomir \\
怕冷 & pà lěng & v. & avoir froid / frissonner \\
出汗 & chū hàn & v. & transpirer \\
验血 & yànxiě & v. & faire une prise de sang (diagnostic) \\
抗疟药 & kàngnüè yào & n. & antipaludéen \\
蚊帐 & wénzhàng & n. & moustiquaire \\
按时 & ànshí & adv. & à l'heure / ponctuellement \\
\hline
\end{tabularx}
\end{center}
\end{exemplebox}

\begin{attention}[Erreurs fréquentes : Consultation médicale]
\begin{itemize}
\item \textbf{\zh{了} (le) et la durée :} Deux structures existent. 1) Action accomplie + durée à part : \zh{我发烧了，已经两天了} (dialogue). 2) Phrase unique « durée en cours » : \zh{我发烧两天了} (Wǒ fāshāo liǎng tiān le) — un seul \zh{了} final. Évitez \zh{*我发烧了两天了} (double \zh{了} lourd).
\item \textbf{\zh{有} vs \zh{是} :} Pour les symptômes, on utilise souvent \zh{有} (avoir) : \zh{有发烧}, \zh{有头疼}. Ne pas dire \zh{*我是发烧} (Je suis fièvre).
\item \textbf{Position de \zh{应该} :} Avant le verbe principal. \zh{你应该去医院} (Juste). \zh{*你去应该医院} (Faux).
\item \textbf{Prononciation \zh{疟} (nüè) :} 4ème ton, voyelle \zh{üe} (comme \zh{yuè} lune). Ne pas confondre avec \zh{虐} (nuè, maltraiter) même son, sens différent.
\end{itemize}
\end{attention}

\courssub{Débat guidé : « Faut-il interdire les bouteilles en plastique au lycée ? »}

\begin{popfigure}
\centering
\begin{tikzpicture}[
    node distance=1.2cm and 1.5cm,
    box/.style={draw=popblue, thick, rounded corners, fill=popblueL, text width=3.2cm, align=center, minimum height=1.8cm, font=\small},
    arrow/.style={-Stealth, thick, popdark},
    phase/.style={font=\bfseries\small, popblue, anchor=west}
]
% Nodes
\node[box, align=center] (prep){\textbf{1. Préparation (5 min)}\\\textit{Équipes Pour / Contre}\\\textit{3 arguments par élève}\\\textit{Mots-clés sur fiche}};
\node[box, right=of prep, align=center] (exposes){\textbf{2. Exposés (2 min/élève)}\\\textit{Équipe Pour : 1, 2, 3...}\\\textit{Équipe Contre : 1, 2, 3...}\\\textit{Formules : 我觉得...因为...}};
\node[box, right=of exposes, align=center] (reactions){\textbf{3. Réactions croisées (5 min)}\\\textit{Modérateur relance :}\\\textit{\zh{你觉得呢？} / \zh{有不同意见吗？}}\\\textit{Répliques courtes}};
\node[box, right=of reactions, align=center] (conclu){\textbf{4. Conclusion Modérateur (2 min)}\\\textit{Synthèse : \zh{总之...}}\\\textit{Vote main levée}\\\textit{Remarques prof}};

% Arrows
\draw[arrow] (prep.east) -- (exposes.west);
\draw[arrow] (exposes.east) -- (reactions.west);
\draw[arrow] (reactions.east) -- (conclu.west);

% Phase labels below
\node[phase] at (prep.south) [yshift=-0.8cm] {Phase 1};
\node[phase] at (exposes.south) [yshift=-0.8cm] {Phase 2};
\node[phase] at (reactions.south) [yshift=-0.8cm] {Phase 3};
\node[phase] at (conclu.south) [yshift=-0.8cm] {Phase 4};
\end{tikzpicture}
\legende{Frise du déroulement du débat guidé en classe (durée totale ~15-20 min)}
\end{popfigure}

\begin{methode}[Déroulement du débat guidé (15-20 minutes)]
\begin{enumerate}
\item \textbf{Constitution des équipes (1 min) :} Tirage au sort : Équipe A = \textbf{Pour} (赞成 \zh{赞成} \emph{zànchéng}), Équipe B = \textbf{Contre} (反对 \zh{反对} \emph{fǎnduì}). Un Modérateur (主持人 \zh{主持人} \emph{zhǔchírén}) désigné (élève ou prof).
\item \textbf{Préparation écrite silencieuse (5 min) :} Chaque élève note \textbf{3 arguments} en chinois (mots-clés + structure \zh{我觉得...因为...}). \textbf{Interdiction de lire son texte} pendant l'oral.
\item \textbf{Tour d'exposés (2 min par élève) :} Alternance Pour / Contre. L'élève se lève, regarde la classe, parle lentement. Ex: \zh{我觉得应该禁止塑料瓶，因为污染环境。} (Wǒ juéde yīnggāi jìnzhǐ sùliáopíng, yīnwèi wūrǎn huánjìng.)
\item \textbf{Phase de réactions croisées (5 min) :} Le modérateur ouvre la parole : \zh{你觉得呢？} (Nǐ juéde ne? — Et toi, qu'en penses-tu ?) / \zh{有不同意见吗？} (Yǒu bùtóng yìjiàn ma? — Y a-t-il un avis différent ?). Les élèves réagissent directement aux arguments adverses avec \zh{但是...} / \zh{可是...} / \zh{我不同意，因为...}.
\item \textbf{Conclusion du modérateur (2 min) :} Synthèse bilingue si besoin. Structure : \zh{总之，赞成方认为...，反对方认为...，建议...} (Zǒngzhī, zànchéng fāng rènwéi..., fǎnduì fāng rènwéi..., jiànyi...). Vote à main levée.
\end{enumerate}
\end{methode}

\begin{exemplebox}[Banque d'arguments types (Pour / Contre)]
\begin{center}
\begin{tabularx}{\linewidth}{|X|X|}
\hline \rowcolor{poppinkL} \textbf{Équipe POUR (赞成)} & \textbf{ÉQUIPE CONTRE (反对)} \\ \hline
\zh{1. 塑料瓶污染环境，分解要几百年。} (Sùliáopíng wūrǎn huánjìng, fēnjiě yào jǐbǎi nián.) & \zh{1. 禁了塑料瓶，喝水不方便，没地方装水。} (Jìn le sùliáopíng, hē shuǐ bù fāngbiàn, méi dìfang zhuāng shuǐ.) \\
\zh{2. 学校可以安装直饮水机，自带水杯。} (Xuéxiào kěyǐ ānzhuāng zhíyǐn shuǐjī, zìdài shuǐbēi.) & \zh{2. 直饮水机水质不一定好，维护贵。} (Zhíyǐn shuǐjī shuǐzhì bù yídìng hǎo, wéihù guì.) \\
\zh{3. 培养环保习惯，从小事做起。} (Péiyǎng huánbǎo xíguàn, cóng xiǎoshì zuòqǐ.) & \zh{3. 应该加强回收利用，而不是禁止。} (Yīnggāi jiāqiáng huíshōu lìyòng, ér bùshì jìnzhǐ.) \\
\zh{4. 减少垃圾，校园更干净美观。} (Jiǎnshǎo lājī, xiàoyuán gèng gānjìng měiguān.) & \zh{4. 很多同学买饮料，瓶子是必须的。} (Hěnduō tóngxué mǎi yǐnliào, píngzi shì bìxū de.) \\
\hline
\end{tabularx}
\end{center}
\end{exemplebox}

\begin{savaistu}[Structure comparative \zh{比} (bǐ) pour enrichir l'exposé]
La structure \zh{A 比 B + Adj.} (A est plus [Adj] que B) est très valorisée à l'oral pour comparer deux solutions.
\begin{itemize}
\item \zh{回收比扔垃圾好。} (Huíshōu bǐ rēng lājī hǎo.) — Recycler vaut mieux que jeter.
\item \zh{自带水杯比买塑料瓶便宜。} (Zìdài shuǐbēi bǐ mǎi sùliáopíng piányi.) — Apporter sa gourde est moins cher qu'acheter des bouteilles.
\item \zh{预防疟疾比治疗疟疾容易，也省钱。} (Yùfáng nüèji bǐ zhìliáo nüèji róngyì, yě shěngqián.) — Prévenir le paludisme est plus facile et moins cher que le soigner.
\end{itemize}
\textbf{Attention :} Pas de \zh{更} (gèng) ou \zh{很} (hěn) après l'adjectif dans la structure \zh{比}. \textbf{Faux :} \zh{*回收比扔垃圾更好}. \textbf{Juste :} \zh{回收比扔垃圾好得多} (beaucoup mieux).
\end{savaistu}

\courssub{Grille d'observation et débriefing collectif}

\begin{propriete}[Grille d'observation par les pairs (Grille d'évaluation orale)]
Chaque observateur (ou le professeur) coche pour chaque intervenant :
\begin{center}
\begin{tabularx}{\linewidth}{|X|c|c|c|c|}
\hline \rowcolor{poppurpleL} \textbf{Critère} & \textbf{Excellent (4)} & \textbf{Bien (3)} & \textbf{Passable (2)} & \textbf{Insuffisant (1)} \\ \hline
\textbf{Prononciation / Tons} & Tons justes, flux naturel & 1-2 erreurs de ton non génantes & Erreurs fréquentes, compréhension difficile & Incompréhensible \\ \hline
\textbf{Formules fonctionnelles} & Utilise \zh{我觉得/因为/但是/总之} à propos & Utilise 2-3 formules correctes & Utilise 1 formule, répétitive & Pas de formule, français \\ \hline
\textbf{Fluidité / Débit} & Parle sans lire, débit normal & Légères hésitations, ne lit pas & Hésitations longues, lit par moments & Lit entièrement, très lent \\ \hline
\textbf{Participation / Réactivité} & Réagit aux autres, relance & Répond si interpellé & Ne parle que lors de son tour & Muet hors tour imposé \\ \hline
\textbf{Vocabulaire thème} & Lexique précis (recyclage/paludisme) & Vocabulaire de base correct & Vocabulaire pauvre, anglicismes & Hors sujet \\ \hline
\end{tabularx}
\end{center}
\end{propriete}

\begin{methode}[Débriefing collectif (5-10 minutes)]
\begin{enumerate}
\item \textbf{Retour positif (Points forts) :} Le modérateur (ou prof) cite 2-3 réussites précises. Ex: \zh{张三的声调很标准，特别是“疟疾”读得很准。} / \zh{李四用了“而且”连接两个论点，很流畅。}
\item \textbf{Axes de progrès (Collectifs) :} 1 point de grammaire, 1 point de prononciation, 1 point de stratégie.
    \begin{itemize}
    \item \textit{Grammaire :} « Beaucoup ont oublié \zh{了} pour la durée (\zh{两天了}). Rappel : Verbe + Durée + \zh{了}. »
    \item \textit{Prononciation :} « Le ton neutre de \zh{吗} (ma) et \zh{呢} (ne) doit être léger et court. »
    \item \textit{Stratégie :} « Ne lisez pas ! Regardez le public. Utilisez \zh{第一，第二} pour structurer. »
    \end{itemize}
\item \textbf{Auto-évaluation :} Chaque élève note sur son carnet : « Ce que j'ai bien fait » / « Ce que je révise ce soir ».
\end{enumerate}
\end{methode}

\begin{exoresolu}[Analyse d'une prise de parole élève]
\textbf{Énoncé :} Un élève (Équipe Contre) dit : \zh{*我觉得不应该禁止塑料瓶，因为学生喜欢喝饮料，而且买水很方便，但是环境重要。}
\textbf{Analyse et correction :}
\begin{enumerate}
\item \textbf{Structure :} \zh{我觉得不应该...} (Bien). \zh{因为...而且...} (Bien, coordination).
\item \textbf{Problème logique / Connecteur :} \zh{但是} (mais) introduit une concession qui contredit l'argument principal. Ici, l'élève dit « les élèves aiment boire, c'est pratique, \textbf{mais} l'environnement est important ». Cela affaiblit sa position « Contre ».
\item \textbf{Correction :} Remplacer \zh{但是} par \zh{虽然...但是...} en début de phrase ou garder une argumentation unilatérale pour le tour de parole principal.
\item \textbf{Version améliorée :} \zh{虽然环境很重要，但是我觉得不应该完全禁止塑料瓶，因为同学们需要方便，而且学校回收做得不够好。} (Suīrán huánjìng hěn zhòngyào, dànshì wǒ juéde bù yīnggāi wánquán jìnzhǐ sùliáopíng, yīnwèi tóngxuémen xūyào fāngbiàn, érqiě xuéxiào huíshōu zuò de bù gòu hǎo.) — \textit{Bien que l'environnement soit important, je pense qu'on ne devrait pas interdire totalement les bouteilles, car les élèves ont besoin de commodité, et de plus le recyclage à l'école n'est pas assez bon.}
\end{enumerate}
\end{exoresolu}

\begin{exercice}[Entraînement personnel / Devoirs]
\begin{enumerate}
\item \textbf{Écrit (Préparation orale) :} Rédigez vos 3 arguments pour le débat « Interdire les bouteilles en plastique » (choisissez votre camp) en caractères + pinyin. Utilisez au moins : \zh{我觉得}, \zh{因为}, \zh{而且}, \zh{比}, \zh{应该}.
\item \textbf{Oral (Enregistrement) :} Enregistrez-vous (1 min 30) en simulant l'interview « Recyclage » (jouez les deux rôles en changeant de voix) ou la consultation « Paludisme ». Vérifiez vos tons sur \zh{疟疾}, \zh{蚊帐}, \zh{回收}, \zh{分类}.
\item \textbf{Traduction (Thème) :} Traduisez en chinois :
    \begin{itemize}
    \item Le médecin me dit de boire plus d'eau et de me reposer.
    \item Recycler le papier est plus facile que recycler le plastique.
    \item À mon avis, nous devrions installer plus de poubelles dans la cour.
    \end{itemize}
\end{enumerate}
\end{exercice}

\begin{corrige}[Exercice 3 - Thème]
\begin{enumerate}
\item \zh{医生说我应该多喝水，多休息。} (Yīshēng shuō wǒ yīnggāi duō hē shuǐ, duō xiūxi.)
\item \zh{回收纸比回收塑料容易。} (Huíshōu zhǐ bǐ huíshōu sùliáo róngyì.)
\item \zh{我觉得我们应该在操场安装更多垃圾桶。} (Wǒ juéde wǒmen yīnggāi zài cāochǎng ānzhuāng gèng duō lājītǒng.)
\end{enumerate}
\end{corrige}

\newpage

\coursec{Activité d'intégration}

\coursec{Leçon 15 — Expression orale : recyclage et paludisme}

\courssub{Objectifs de l'activité}
Mobiliser l'ensemble des ressources linguistiques et communicationnelles travaillées au cours de la leçon (lexique spécialisé, structures de l'exposé et du débat, formules fonctionnelles, travail prosodique) dans une situation de communication authentique, stimulante et évaluée par les pairs. L'élève doit être capable de :
\begin{itemize}
    \item Structurer et présenter un mini-exposé cohérent de 1 à 2 minutes sur l'un des deux thèmes au programme (recyclage au lycée ou lutte contre le paludisme dans le quartier), en respectant le schéma \textbf{ouverture -- arguments -- conclusion}.
    \item Interagir spontanément en tant que public : poser une question pertinente (\zh{你有什么看法？} \emph{Nǐ yǒu shénme kànfǎ?} « Quel est ton avis ? »), exprimer un accord ou un désaccord argumenté (\zh{我同意，因为……} / \zh{我不同意，因为……}).
    \item Maîtriser la prononciation des tons (notamment les tons 2/3, le \emph{sandhi} du 3e ton, les tons neutres sur les particules) et l'intonation énonciative/interrogative.
    \item Évaluer la prestation d'un camarade à l'aide d'une grille critériée (tons, formules, fluidité, richesse lexicale) et décerner le titre de \zh{最佳演讲者} \emph{zuìjiā yǎnjiǎngzhě} (meilleur orateur).
\end{itemize}

\courssub{Situation}
\begin{textebox}[Affiche officielle du Forum Santé-Environnement du Lycée Bilingue de Buea]
\zh{摆鲁双语高中 “健康与环境” 论坛} \\
\zh{时间：下周三下午 3 点} \\
\zh{地点：多媒体教室} \\
\zh{主题一：校园垃圾分类与回收行动} \\
\zh{主题二：社区防疟疾，从我做起} \\
\zh{形式：自由报名，分组演讲 + 现场问答} \\
\zh{评选：“最佳演讲者” 证书与奖品} \\
\zh{全体高三 (A) 同学踊跃参加！} \\
\emph{Báilǔ Shuāngyǔ Gāozhōng ``Jiànkāng yǔ Huánjìng'' Lùntán -- Shíjiān: Xià zhōusān xiàwǔ 3 diǎn -- Dìdiān: Duōméitǐ Jiàoshì -- Zhǔtí yī: Xiàoyuán Lājī Fēnlèi yǔ Huíshōu Xíngdòng -- Zhǔtí èr: Shèqū Fáng Nuèjí, Cóng Wǒ Zuòqǐ -- Xíngshì: Zìyóu Bàomíng, Fēnzǔ Yǎnjiǎng + Xiànchǎng Wèndá -- Píngxuǎn: ``Zuìjiā Yǎnjiǎngzhě'' Zhèngshū yǔ Jiǎngpǐn -- Quántǐ Gāosān (A) Tóngxué Yǒngyuè Cānjiā!} \\
« Forum "Santé et Environnement" du Lycée Bilingue de Buea -- Heure : Mercredi prochain 15h -- Lieu : Salle multimédia -- Thème 1 : Tri et recyclage des déchets au lycée -- Thème 2 : Prévention du paludisme dans la communauté, agissons -- Format : Inscription libre, exposés par groupes + questions-réponses en direct -- Prix : Certificat et lots pour le "Meilleur orateur" -- Tous les élèves de Terminale A sont vivement invités à participer ! »
\end{textebox}

\courssub{Consignes}
\begin{enumerate}
    \item \textbf{Préparation individuelle (15 min) :} Choisissez l'un des deux thèmes. Rédigez les notes de votre exposé (mots-clés, structure, 3-4 phrases clés en caractères) sur une fiche bristol. Utilisez obligatoirement : une formule d'ouverture (\zh{大家好，今天我要谈谈……}), au moins deux connecteurs logiques (\zh{首先} / \zh{其次} / \zh{另外} / \zh{然而}), une opinion personnelle (\zh{我认为} / \zh{我觉得}) et une formule de conclusion (\zh{总之，让我们一起行动吧！}).
    \item \textbf{Phase 1 -- Groupe A expose, Groupe B public (20 min) :} Les élèves du Groupe A passent au pupitre (1 par thème). Le Groupe B écoute activement, note une question ou une réaction sur un papier. À la fin de chaque exposé, 2 à 3 élèves du Groupe B interviennent : \zh{请问……，你有什么看法？} (Question) puis \zh{我同意/不同意你的观点，因为……} (Réaction).
    \item \textbf{Phase 2 -- Inversion des rôles (20 min) :} Le Groupe B présente ses exposés (sur l'autre thème de préférence). Le Groupe A joue le public et le jury.
    \item \textbf{Évaluation par les pairs (10 min) :} Chaque juré remplit la grille d'évaluation pour 3 orateurs (tons /2, formules /2, fluidité /2, vocabulaire /2, interaction /2). Calcul de la moyenne. Désignation du \zh{最佳演讲者}.
    \item \textbf{Bilan collectif (5 min) :} L'enseignant commente les prestations notables, corrige les erreurs systématiques de tons ou de structure entendues, félicite le lauréat.
\end{enumerate}

\courssub{Ressources à mobiliser}

\begin{center}
\begin{tabularx}{\linewidth}{|X|X|X|X|}
\hline
\rowcolor{popblueL} \textbf{Structure / Fonction} & \textbf{Formule type (Caractères + Pinyin)} & \textbf{Traduction} & \textbf{Point de vigilance (Tons / Grammaire)} \\
\hline
Ouverture & \zh{大家好，今天我要谈谈……} \emph{Dàjiā hǎo, jīntiān wǒ yào tántan……} & Bonjour à tous, aujourd'hui je vais parler de… & \zh{谈} 2e ton ; \zh{谈谈} reduplication (ton neutre sur 2e). \\
\hline
Énumérer arguments & \zh{首先，……；其次，……；另外，……} \emph{Shǒuxiān……; qícì……; lìngwài……} & Premièrement… ; Deuxièmement… ; En outre… & \zh{首} 3e ton $\rightarrow$ 2e ton devant \zh{先} 1er ton (\emph{shóuxiān}). \\
\hline
Donner opinion & \zh{我认为 / 我觉得……} \emph{Wǒ rènwéi / Wǒ juéde……} & Je pense que / Je trouve que… & \zh{认} 4e ton, \zh{为} 2e ton ; \zh{觉} 2e ton, \zh{得} ton neutre. \\
\hline
Contraste / Nuance & \zh{虽然……，但是……} \emph{Suīrán……, dànshì……} & Bien que… , cependant… & Structure figée. \zh{虽} 1er ton, \zh{然} 2e ton. \\
\hline
Conclusion & \zh{总之，让我们一起行动吧！} \emph{Zǒngzhī, ràng wǒmen yīqǐ xíngdòng ba!} & En bref, agissons ensemble ! & \zh{总} 3e ton $\rightarrow$ 2e ton devant \zh{之} 1er ton : \emph{zóngzhī}. \zh{吧} ton neutre. \\
\hline
Question public & \zh{请问，关于……，你有什么看法？} \emph{Qǐngwèn, guānyú……, nǐ yǒu shénme kànfǎ?} & Excusez-moi, concernant…, quel est votre avis? & \zh{请} 3e ton $\rightarrow$ 2e ton devant \zh{问} 4e ton : \emph{qíngwèn}. \zh{看} 4e ton, \zh{法} 3e ton. \\
\hline
Accord & \zh{我完全同意，因为……} \emph{Wǒ wánquán tóngyì, yīnwèi……} & Je suis tout à fait d'accord, car… & \zh{全} 2e ton, \zh{权} 2e ton ; \zh{同} 2e ton, \zh{意} 4e ton. \\
\hline
Désaccord poli & \zh{我有不同的看法，我觉得……} \emph{Wǒ yǒu bùtóng de kànfǎ, wǒ juéde……} & J'ai un avis différent, je trouve que… & \zh{不} 4e ton $\rightarrow$ 2e ton devant \zh{同} 2e ton : \emph{bútóng}. \\
\hline
\end{tabularx}
\end{center}

\begin{definition}[Lexique spécialisé -- Recyclage (回收 \emph{Huíshōu})]
\begin{center}
\begin{tabularx}{\linewidth}{|X|X|X|X|}
\hline
\rowcolor{popgreenL} \textbf{Caractères} & \textbf{Pinyin} & \textbf{Nature} & \textbf{Traduction} \\
\hline
\zh{垃圾分类} & lājī fēnlèi & N & Tri des déchets / tri sélectif \\
\zh{塑料瓶} & sùliào píng & N & Bouteille en plastique \\
\zh{回收箱} & huíshōu xiāng & N & Bac de recyclage / poubelle de tri \\
\zh{环保} & huánbǎo & Adj / V & Écologique / Protéger l'environnement \\
\zh{二手} & èrshǒu & Adj & D'occasion \\
\zh{减少} & jiǎnshǎo & V & Réduire / diminuer \\
\zh{污染} & wūrǎn & V / N & Polluer / pollution \\
\zh{投放} & tóufàng & V & Déposer (dans un bac) / jeter (tri) \\
\hline
\end{tabularx}
\end{center}
\end{definition}

\begin{definition}[Lexique spécialisé -- Paludisme (疟疾 \emph{Nuèjí})]
\begin{center}
\begin{tabularx}{\linewidth}{|X|X|X|X|}
\hline
\rowcolor{poporangeL} \textbf{Caractères} & \textbf{Pinyin} & \textbf{Nature} & \textbf{Traduction} \\
\hline
\zh{疟疾} & nuèjí & N & Paludisme (malaria) \\
\zh{蚊子} & wénzi & N & Moustique (ton neutre sur \zh{子}) \\
\zh{蚊帐} & wénzhàng & N & Moustiquaire (lit) \\
\zh{预防} & yùfáng & V & Prévenir / prévention \\
\zh{症状} & zhèngzhuàng & N & Symptôme \\
\zh{发烧} & fāshāo & VO & Avoir de la fièvre \\
\zh{就医} & jiùyī & V & Consulter un médecin / aller à l'hôpital \\
\zh{社区} & shèqū & N & Quartier / communauté \\
\zh{喷雾} & pēnwù & V / N & Pulvériser / spray (insecticide) \\
\zh{积水} & jīshuǐ & N / V & Eau stagnante / stagner (eau) \\
\hline
\end{tabularx}
\end{center}
\end{definition}

\courssub{Travail des tons et de la prononciation}
\begin{propriete}[Rappel : Tons critiques à surveiller pour l'oral]
\begin{itemize}
    \item \textbf{3e ton + 3e ton (Sandhi) :} \zh{很好} $\rightarrow$ \emph{hén hǎo} ; \zh{我也} $\rightarrow$ \emph{wó yě} ; \zh{可以} $\rightarrow$ \emph{kěyǐ} (pas de changement si isolé, mais en phrase \zh{我可以} $\rightarrow$ \emph{wó kěyǐ}).
    \item \textbf{\zh{不} (bù) :} 4e ton seul ou devant 1/2/3e ton (\zh{不是} \emph{bú shì}) ; 2e ton devant 4e ton (\zh{不同意} \emph{bútóngyì}).
    \item \textbf{\zh{一} (yī) :} 1er ton seul/chiffre/fin de mot ; 4e ton devant 1/2/3e ton (\zh{一个} \emph{yí ge}) ; 2e ton devant 4e ton (\zh{一定} \emph{yídìng}).
    \item \textbf{Tons neutres (轻声) :} Particules \zh{了 吗 吧 呢 的 地 得} ; suffixes \zh{子 们} ; 2e syllabe de reduplication verbale \zh{看看 听听 谈谈}.
    \item \textbf{Intonation interrogative :} Montée finale sur \zh{吗} / \zh{呢} / \zh{什么看法} ; chute sur question \zh{为什么} / \zh{怎么}.
\end{itemize}
\end{propriete}

\courssub{Dialogue modèle 1 : Exposé sur le recyclage}
\begin{exemplebox}[Mini-exposé : 校园垃圾分类，从我做起]
\zh{大家好，今天我要谈谈校园垃圾分类。} \\
\zh{首先，我们学校每天产生很多塑料瓶和纸张，如果不分类，会造成严重污染。} \\
\zh{其次，回收利用可以变废为宝，保护环境。} \\
\zh{另外，我觉得同学们应该养成好习惯，把瓶子投进蓝色的回收箱，纸张放进黄色的箱子。} \\
\zh{虽然改变习惯不容易，但是只要大家一起努力，校园就会更美丽。} \\
\zh{总之，让我们一起行动吧，从我做起，从现在做起！谢谢大家。} \\
\emph{Dàjiā hǎo, jīntiān wǒ yào tántan xiàoyuán lājī fēnlèi. -- Shǒuxiān, wǒmen xuéxiào měitiān chǎnshēng hěnduō sùliào píng hé zhǐzhāng, rúguǒ bù fēnlèi, huì zàochéng yánzhòng wūrǎn. -- Qícì, huíshōu lìyòng kěyǐ biànfèiwéibǎo, bǎohù huánjìng. -- Lìngwài, wǒ juéde tóngxuémen yīnggāi yǎngchéng hǎo xíguàn, bǎ píngzi tóu jìn lánsè de huíshōu xiāng, zhǐzhāng fàng jìn huángsè de xiāngzi. -- Suīrán gǎibiàn xíguàn bù róngyì, dànshì zhǐyào dàjiā yīqǐ nǔlì, xiàoyuán jiù huì gèng měilì. -- Zǒngzhī, ràng wǒmen yīqǐ xíngdòng ba, cóng wǒ zuòqǐ, cóng xiànzài zuòqǐ! Xièxie dàjiā.} \\
« Bonjour à tous, aujourd'hui je veux parler du tri des déchets au lycée. Premièrement, notre école produit chaque jour beaucoup de bouteilles en plastique et de papier ; si on ne trie pas, cela cause une pollution grave. Deuxièmement, le recyclage permet de transformer les déchets en ressources et de protéger l'environnement. En outre, je pense que les élèves doivent prendre de bonnes habitudes : déposer les bouteilles dans le bac de recyclage bleu, le papier dans le bac jaune. Bien que changer ses habitudes ne soit pas facile, tant que tout le monde fait des efforts ensemble, le campus deviendra plus beau. En bref, agissons ensemble, en commençant par moi, en commençant maintenant ! Merci à tous. »
\end{exemplebox}

\begin{exemplebox}[Échange type Questions-Réponses (Recyclage)]
\zh{观众甲：请问，你觉得学校现在的回收箱够不够用？} \\
\zh{演讲者：我觉得不太够用。因为人多，箱子很快满了，建议学校增加箱子，并且加强宣传。} \\
\zh{观众乙：我同意你的看法。另外，食堂的泡沫饭盒怎么处理？} \\
\zh{演讲者：这是个好问题。泡沫饭盒很难降解，最好自带饭盒，减少一次性用品。} \\
\emph{Guānzhòng jiǎ: Qǐngwèn, nǐ juéde xuéxiào xiànzài de huíshōu xiāng gòu bu gòuyòng? -- Yǎnjiǎngzhě: Wǒ juéde bù tài gòuyòng. Yīnwèi rén duō, xiāngzi hěn kuài mǎn le, jiànyì xuéxiào zēngjiā xiāngzi, bìngqiě jiāqiáng xuānchuán. -- Guānzhòng yǐ: Wǒ tóngyì nǐ de kànfǎ. Lìngwài, shítáng de pàomò fànhé zěnme chǔlǐ? -- Yǎnjiǎngzhě: Zhè shì gè hǎo wèntí. Pàomò fànhé hěn nán jiàngjiě, zuìhǎo zìdài fànhé, jiǎnshǎo yīcìxìng yòngpǐn.} \\
« Spectateur A : Excusez-moi, pensez-vous que les bacs de recyclage de l'école sont suffisants ? -- Orateur : Je trouve que non. Parce qu'il y a beaucoup de monde, les bacs se remplissent vite ; je suggère à l'école d'en ajouter et de renforcer la sensibilisation. -- Spectateur B : Je suis d'accord avec ton avis. Par ailleurs, comment traite-t-on les boîtes en polystyrène de la cantine ? -- Orateur : C'est une bonne question. Les boîtes en polystyrène se dégradent difficilement ; le mieux est d'apporter sa propre boîte repas, pour réduire les articles jetables. »
\end{exemplebox}

\courssub{Dialogue modèle 2 : Débat sur le paludisme}
\begin{exemplebox}[Mini-exposé : 社区防疟疾，从我做起]
\zh{大家好，今天我演讲的主题是社区预防疟疾。} \\
\zh{首先，疟疾是由蚊子传播的，雨季积水多，蚊子容易滋生。} \\
\zh{其次，主要症状是发烧、发冷、头痛，发现症状要及时就医。} \\
\zh{另外，预防最关键：晚上必须睡蚊帐，房间喷雾驱蚊，清理屋前屋后的积水。} \\
\zh{虽然蚊帐有点热，但是为了健康，这点不便值得。} \\
\zh{总之，防疟疾人人有责，让我们从我做起，共建无疟社区！谢谢大家。} \\
\emph{Dàjiā hǎo, jīntiān wǒ yǎnjiǎng de zhǔtí shì shèqū yùfáng nuèjí. -- Shǒuxiān, nuèjí shì yóu wénzi chuánbō de, yǔjì jīshuǐ duō, wénzi róngyì zīshēng. -- Qícì, zhǔyào zhèngzhuàng shì fāshāo, fālěng, tóutòng, fāxiàn zhèngzhuàng yào jíshí jiùyī. -- Lìngwài, yùfáng zuì guānjiàn: wǎnshang bìxū shuì wénzhàng, fángjiān pēnwù qūwén, qīnglǐ wūqián wūhòu de jīshuǐ. -- Suīrán wénzhàng yǒudiǎn rè, dànshì wèile jiànkāng, zhè diǎn bùbiàn zhíde. -- Zǒngzhī, fáng nuèjí rénrén yǒuzé, ràng wǒmen cóng wǒ zuòqǐ, gòngjiàn wú nuè shèqū! Xièxie dàjiā.} \\
« Bonjour à tous, mon exposé aujourd'hui porte sur la prévention du paludisme dans la communauté. Premièrement, le paludisme est transmis par les moustiques ; en saison des pluies, l'eau stagnante abonde, les moustiques prolifèrent facilement. Deuxièmement, les symptômes principaux sont la fièvre, les frissons, les maux de tête ; dès l'apparition des symptômes, il faut consulter d'urgence. En outre, la prévention est cruciale : dormir sous moustiquaire la nuit, pulvériser un répulsif dans la pièce, éliminer l'eau stagnante autour de la maison. Bien que la moustiquaire soit un peu chaude, pour la santé, ce désagrément en vaut la peine. En bref, prévenir le paludisme est le devoir de tous, agissons à partir de moi, construisons ensemble une communauté sans paludisme ! Merci à tous. »
\end{exemplebox}

\begin{exemplebox}[Échange type Questions-Réponses (Paludisme)]
\zh{观众甲：请问，社区里有人不爱睡蚊帐，怎么办？} \\
\zh{演讲者：可以宣传蚊帐的重要性，也可以推荐驱蚊手环或电蚊拍，多种方法结合。} \\
\zh{观众乙：我不同意只靠蚊帐。我觉得环境治理更根本，要填平坑洼，清理垃圾。} \\
\zh{演讲者：你的观点很有道理。环境治理确实是源头预防，需要社区和政府一起努力。} \\
\emph{Guānzhòng jiǎ: Qǐngwèn, shèqū lǐ yǒu rén bù ài shuì wénzhàng, zěnme bàn? -- Yǎnjiǎngzhě: Kěyǐ xuānchuán wénzhàng de zhòngyàoxìng, yě kěyǐ tuījiàn qūwén shǒuhuán huò diàn wénpāi, duōzhǒng fāngfǎ jiéhé. -- Guānzhòng yǐ: Wǒ bù tóngyì zhǐ kào wénzhàng. Wǒ juéde huánjìng zhìlǐ gèng gēnběn, yào tiánpíng kēngwā, qīnglǐ lājī. -- Yǎnjiǎngzhě: Nǐ de guāndiǎn hěn yǒu dàolǐ. Huánjìng zhìlǐ quèshì shì yuántóu yùfáng, xūyào shèqū hé zhèngfǔ yīqǐ nǔlì.} \\
« Spectateur A : Que faire si certains dans la communauté n'aiment pas dormir sous moustiquaire ? -- Orateur : On peut sensibiliser à l'importance de la moustiquaire, recommander aussi bracelets répulsifs ou raquettes électriques, combiner plusieurs méthodes. -- Spectateur B : Je ne suis pas d'accord pour compter seulement sur la moustiquaire. Je trouve que l'assainissement de l'environnement est plus fondamental : combler les trous, nettoyer les déchets. -- Orateur : Ton avis est très pertinent. L'assainissement est bien une prévention à la source, cela demande les efforts conjoints de la communauté et du gouvernement. »
\end{exemplebox}

\courssub{Proposition de production et corrigé commenté}
\begin{exoresolu}[Modèle d'exposé élève (Thème 1 : Recyclage) + Interaction + Commentaire linguistique]
\textbf{Énoncé :} Présenter un exposé de 1'30 sur le tri des déchets au Lycée Bilingue de Buea, puis répondre à une question du public.

\tcblower

\textbf{1. Exposé modèle (Thème 1) :}
\zh{大家好，今天我要谈谈校园垃圾分类。} \\
\emph{Dàjiā hǎo, jīntiān wǒ yào tántan xiàoyuán lājī fēnlèi.} \\
« Bonjour à tous, aujourd'hui je veux parler du tri des déchets au lycée. »

\zh{首先，我们学校每天产生很多塑料瓶和纸张，如果不分类，会造成严重污染。} \\
\emph{Shǒuxiān, wǒmen xuéxiào měitiān chǎnshēng hěnduō sùliào píng hé zhǐzhāng, rúguǒ bù fēnlèi, huì zàochéng yánzhòng wūrǎn.} \\
« Premièrement, notre école produit chaque jour beaucoup de bouteilles en plastique et de papier ; si on ne trie pas, cela cause une pollution grave. »

\zh{其次，回收利用可以变废为宝，保护环境。} \\
\emph{Qícì, huíshōu lìyòng kěyǐ biànfèiwéibǎo, bǎohù huánjìng.} \\
« Deuxièmement, le recyclage permet de transformer les déchets en ressources et de protéger l'environnement. »

\zh{另外，我觉得同学们应该养成好习惯，把瓶子投进蓝色的回收箱，纸张放进黄色的箱子。} \\
\emph{Lìngwài, wǒ juéde tóngxuémen yīnggāi yǎngchéng hǎo xíguàn, bǎ píngzi tóu jìn lánsè de huíshōu xiāng, zhǐzhāng fàng jìn huángsè de xiāngzi.} \\
« En outre, je pense que les élèves doivent prendre de bonnes habitudes : déposer les bouteilles dans le bac de recyclage bleu, le papier dans le bac jaune. »

\zh{虽然改变习惯不容易，但是只要大家一起努力，校园就会更美丽。} \\
\emph{Suīrán gǎibiàn xíguàn bù róngyì, dànshì zhǐyào dàjiā yīqǐ nǔlì, xiàoyuán jiù huì gèng měilì.} \\
« Bien que changer ses habitudes ne soit pas facile, tant que tout le monde fait des efforts ensemble, le campus deviendra plus beau. »

\zh{总之，让我们一起行动吧，从我做起，从现在做起！谢谢大家。} \\
\emph{Zǒngzhī, ràng wǒmen yīqǐ xíngdòng ba, cóng wǒ zuòqǐ, cóng xiànzài zuòqǐ! Xièxie dàjiā.} \\
« En bref, agissons ensemble, en commençant par moi, en commençant maintenant ! Merci à tous. »

\textbf{2. Interaction type (Public $\rightarrow$ Orateur) :}
\zh{观众甲：请问，你觉得学校现在的回收箱够不够用？} \\
\emph{Guānzhòng jiǎ: Qǐngwèn, nǐ juéde xuéxiào xiànzài de huíshōu xiāng gòu bu gòuyòng?} \\
« Excusez-moi, pensez-vous que les bacs de recyclage de l'école sont suffisants ? »

\zh{演讲者：我觉得不太够用。因为人多，箱子很快满了，建议学校增加箱子，并且加强宣传。} \\
\emph{Yǎnjiǎngzhě: Wǒ juéde bù tài gòuyòng. Yīnwèi rén duō, xiāngzi hěn kuài mǎn le, jiànyì xuéxiào zēngjiā xiāngzi, bìngqiě jiāqiáng xuānchuán.} \\
« Orateur : Je trouve que non. Parce qu'il y a beaucoup de monde, les bacs se remplissent vite ; je suggère à l'école d'en ajouter et de renforcer la sensibilisation. »

\zh{观众乙：我同意你的看法。另外，食堂的泡沫饭盒怎么处理？} \\
\emph{Guānzhòng yǐ: Wǒ tóngyì nǐ de kànfǎ. Lìngwài, shítáng de pàomò fànhé zěnme chǔlǐ?} \\
« Spectateur B : Je suis d'accord avec ton avis. Par ailleurs, comment traite-t-on les boîtes en polystyrène de la cantine ? »

\zh{演讲者：这是个好问题。泡沫饭盒很难降解，最好自带饭盒，减少一次性用品。} \\
\emph{Yǎnjiǎngzhě: Zhè shì gè hǎo wèntí. Pàomò fànhé hěn nán jiàngjiě, zuìhǎo zìdài fànhé, jiǎnshǎo yīcìxìng yòngpǐn.} \\
« Orateur : C'est une bonne question. Les boîtes en polystyrène se dégradent difficilement ; le mieux est d'apporter sa propre boîte repas, pour réduire les articles jetables. »

\textbf{3. Commentaire des choix de langue (Corrigé commenté) :}
\begin{itemize}
    \item \textbf{Structure macro :} Respect strict du schéma \emph{Ouverture (大家好…)} -- \emph{Corps argumenté (首先/其次/另外 + 论据)} -- \emph{Concession (虽然…但是…)} -- \emph{Conclusion mobilisatrice (总之…行动吧)}. Cela assure la cohérence attendue au niveau Terminale.
    \item \textbf{Grammaire / Structures :}
        \begin{itemize}
            \item \zh{如果不分类，会造成……} (Condition \zh{如果} + conséquence \zh{会}) : structure hypothétique réelle, très claire.
            \item \zh{变废为宝} (Chéngyǔ « transformer déchets en trésor ») : idiome de niveau HSK4, valorisé à l'oral.
            \item \zh{建议学校增加箱子} (Verbe \zh{建议} + COD + Verbe) : structure causative / suggestion correcte.
            \item \zh{最好自带饭盒} (\zh{最好} + Verbe) : conseil fort, naturel.
            \item \zh{把瓶子投进……} (Structure \zh{把} + objet + verbe + complément de lieu) : utilisation correcte de la préposition \zh{把} pour marquer la manipulation de l'objet.
        \end{itemize}
    \item \textbf{Tons \& Prosodie (Points d'attention pour l'élève) :}
        \begin{itemize}
            \item \zh{垃圾分类} : \emph{lājī fēnlèi} (1-1-1-4) -- bien distinguer \zh{分} (1) et \zh{类} (4).
            \item \zh{产生} \emph{chǎnshēng} (3-1) -- \zh{产} 3e ton isolé, mais en phrase souvent 2e ton par sandhi si suivi de 3e ton (ici \zh{生} 1er ton $\rightarrow$ reste 3e ou demi-3e).
            \item \zh{回收箱} \emph{huíshōu xiāng} (2-1-1) -- \zh{箱} 1er ton haut et stable.
            \item \zh{不太够用} \emph{bú tài gòuyòng} -- \zh{不} $\rightarrow$ 2e ton devant \zh{太} 4e ton ; \zh{够} 4e ton, \zh{用} 4e ton.
            \item \zh{降解} \emph{jiàngjiě} (4-3) -- \zh{解} 3e ton bas.
            \item \zh{蚊子} \emph{wénzi} -- \zh{子} ton neutre (léger, court), pas \emph{wénzǐ}.
        \end{itemize}
    \item \textbf{Interaction :} L'orateur utilise \zh{这是个好问题} (stratégie de gain de temps / politesse) avant de répondre. Le public utilise \zh{我同意你的看法} + \zh{另外} pour rebondir : excellente dynamique de débat.
\end{itemize}
\end{exoresolu}

\courssub{Bilan de l'activité}
\begin{aretenir}[Bilan compétences -- Forum Santé-Environnement]
\begin{itemize}
    \item \textbf{Production orale continue :} La majorité des élèves a réussi à tenir 1'30 avec une structure visible (mots-clés \zh{首先/其次/总之}). Les exposés sur le paludisme (\zh{疟疾}) ont nécessité plus de lexique médical (\zh{症状、就医、喷雾、积水}), ce qui a révélé des lacunes de vocabulaire à combler.
    \item \textbf{Interaction orale :} Les formules \zh{你有什么看法？} et \zh{我同意/不同意，因为……} sont désormais automatisées. Cependant, l'articulation \zh{因为……所以……} dans la réponse orale reste hésitante pour certains (inversion sujet/verbe ou omission de \zh{所以}).
    \item \textbf{Phonologie (Tons) :} Progrès nets sur le \emph{sandhi} du 3e ton (\zh{很好 \rightarrow hén hǎo}, \zh{我也 \rightarrow wó yě}) et sur \zh{不} (\zh{不同意 \rightarrow bútóngyì}). Persistent : le ton neutre sur \zh{子/们/了/吧} souvent prononcé comme un ton plein (ex: \zh{蚊子} \emph{wénzǐ} au lieu de \emph{wénzi}) ; l'intonation interrogative finale encore plate pour les questions en \zh{怎么/为什么}.
    \item \textbf{Évaluation par les pairs :} La grille a fonctionné comme un levier de méta-cognition. Les élèves juges ont été sévères mais justes sur le critère « Fluidité » (pauses longues, \zh{那个……那个……}). Le titre de \zh{最佳演讲者} a motivé l'engagement.
    \item \textbf{Prochaines étapes :} Réinvestir ces formules dans l'écrit (rédaction d'un article pour le journal du lycée en chinois) et préparer l'épreuve orale du Baccalauréat (sujets type « Environnement / Santé »).
\end{itemize}
\end{aretenir}

\newpage

\coursec{Méthodes et exercices résolus}

\coursec{Leçon 15 — Expression orale : recyclage et paludisme}

\courssub{Mise en route et lexique oral du thème}

\begin{definition}[Vocabulaire thématique : Recyclage (垃圾分类) et Paludisme (疟疾)]
\begin{center}\begin{tabularx}{\linewidth}{|X|X|X|X|}\hline \rowcolor{popblueL} \textbf{Caractères} & \textbf{Pinyin} & \textbf{Nature} & \textbf{Traduction} \\ \hline
垃圾分类 & lājī fēnlèi & VO / N & tri des déchets / tri sélectif \\ \hline
垃圾桶 & lājītǒng & N & poubelle \\ \hline
厨余垃圾 & chúyú lājī & N & déchets de cuisine (humides) \\ \hline
可回收物 & kěhuíshōu wù & N & recyclables \\ \hline
有害垃圾 & yǒuhài lājī & N & déchets dangereux \\ \hline
其他垃圾 & qítā lājī & N & autres déchets (secs/ultimes) \\ \hline
回收利用 & huíshōu lìyòng & VO & recycler / valoriser \\ \hline
压扁 & yā biǎn & V & aplatir (bouteille, canette) \\ \hline
洗干净 & xǐ gānjìng & V + Compl. rés. & laver proprement \\ \hline
疟疾 & nüèjí & N & paludisme (paludisme) \\ \hline
蚊子 & wénzi & N & moustique \\ \hline
蚊帐 & wénzhàng & N & moustiquaire \\ \hline
药蚊帐 & yào wénzhàng & N & moustiquaire imprégnée (MILD) \\ \hline
预防 & yùfáng & V & prévenir \\ \hline
驱蚊液 & qūwén yè & N & répulsif anti-moustiques \\ \hline
积水 & jīshuǐ & N & eau stagnante \\ \hline
繁殖 & fánzhí & V & proliférer, se reproduire \\ \hline
验血 & yànxiě & VO & faire une prise de sang \\ \hline
确诊 & quèzhěn & V & diagnostiquer (confirmer) \\ \hline
延误病情 & yánwù bìngqíng & VO & aggraver/retarder la prise en charge \\ \hline
发烧 & fāshāo & V & avoir de la fièvre \\ \hline
发冷 & fālěng & V & avoir des frissons \\ \hline
退烧药 & tuìshāo yào & N & antipyrétique \\ \hline
建议 & jiànyì & V / N & conseiller / conseil \\ \hline
配合 & pèihé & V & coopérer \\ \hline
\end{tabularx}\end{center}
\end{definition}

\courssub{Dialogue modèle 1 : Exposé sur le recyclage}

\begin{textebox}[Modèle d'exposé : 我家如何做垃圾分类]
大家好！今天我的话题是“我家如何做垃圾分类”。\\
Dàjiā hǎo! Jīntiān wǒ de huàtí shì “Wǒjiā rúhé zuò lājī fēnlèi”.\\
« Bonjour à tous ! Aujourd'hui mon sujet est "Comment ma famille trie les déchets". »

首先，我们家买了四个不同颜色的垃圾桶。\\
Shǒuxiān, wǒmen jiā mǎile sìge bùtóng yánsè de lājītǒng.\\
« D'abord, ma famille a acheté quatre poubelles de couleurs différentes. »

厨余垃圾扔绿色的，可回收物扔蓝色的，有害垃圾扔红色的，其他垃圾扔灰色的。\\
Chúyú lājī rēng lǜsè de, kěhuíshōu wù rēng lánsè de, yǒuhài lājī rēng hóngsè de, qítā lājī rēng huīsè de.\\
« Les déchets de cuisine vont dans la verte, les recyclables dans la bleue, les dangereux dans la rouge, les autres dans la grise. »

然后，全家人养成了分类的好习惯，而且经常提醒邻居也注意分类。\\
Ránhòu, quánjiārén yǎngchéngle fēnlèi de hǎo xíguàn, érqiě jīngcháng tíxíng línjū yě zhùyì fēnlèi.\\
« Ensuite, toute la famille a pris la bonne habitude de trier, et en plus on rappelle souvent aux voisins de faire attention au tri. »

最后，我认为垃圾分类虽然有点麻烦，但能保护环境，很有意义。\\
Zuìhòu, wǒ rènwéi lājī fēnlèi suīrán yǒudiǎn máfan, dàn néng bǎohù huánjìng, hěn yǒu yìyì.\\
« Enfin, je pense que le tri est un peu embêtant, mais ça protège l'environnement, c'est très significatif. »

谢谢大家！\\
Xièxie dàjiā!\\
« Merci à tous ! »
\end{textebox}

\courssub{Dialogue modèle 2 : Débat sur le paludisme}

\begin{textebox}[Dialogue de débat : 治疗疟疾，去医院还是自行用药？]
\textbf{A :} 最近雨季多，蚊子多，得疟疾的人也多。你觉得发烧了该怎么办？\\
Zuìjìn yǔjì duō, wénzi duō, dé nüèjí de rén yě duō. Nǐ juéde fāshāo le gāi zěnme bàn?\\
« Récemment, saison des pluies, beaucoup de moustiques, beaucoup de gens attrapent le palu. Tu penses qu'il faut faire quoi si on a de la fièvre ? »

\textbf{B :} 我觉得应该先去医院验血，确诊了再吃药。因为自己买药吃可能延误病情。\\
Wǒ juéde yīnggāi xiān qù yīyuàn yànxiě, quèzhěn le zài chī yào. Yīnwèi zìjǐ mǎi yào chī kěnéng yánwù bìngqíng.\\
« Je pense qu'il faut d'abord aller à l'hôpital faire une prise de sang, une fois diagnostiqué, prendre les médicaments. Parce que s'acheter soi-même des médocs peut retarder la guérison. »

\textbf{A :} 我不太同意。如果只是轻微发烧，先吃点退烧药、喝点水观察一下也可以。去医院太贵、太慢。\\
Wǒ bù tài tóngyì. Rúguǒ zhǐshì qīngwēi fāshāo, xiān chī diǎn tuìshāo yào, hē diǎn shuǐ guānchá yíxià yě kěyǐ. Qù yīyuàn tài guì, tài màn.\\
« Je ne suis pas trop d'accord. Si c'est juste une légère fièvre, prendre un antipyrétique, boire de l'eau et observer un peu ça va aussi. L'hôpital c'est trop cher, trop lent. »

\textbf{B :} 有道理。但是疟疾不像感冒，不及时治疗很危险。建议你：不管轻重，先去医院查一下比较保险。\\
Yǒu dàoli. Dànshì nüèjí bùxiàng gǎnmào, bù jíshí zhìliáo hěn wéixiǎn. Jiànyì nǐ: bùguǎn qīngzhòng, xiān qù yīyuàn chá yíxià bǐjiào bǎoxiǎn.\\
« C'est vrai / Tu as raison. Mais le palu n'est pas comme un rhume, pas soigné à temps c'est dangereux. Je te conseille : peu importe la gravité, va d'abord checker à l'hôpital, c'est plus sûr. »

\textbf{A :} 好吧，你说得有道理。那我们下周一起去社区卫生服务中心做个宣传吧？\\
Hǎoba, nǐ shuō de yǒu dàoli. Nà wǒmen xiàzhōu yīqǐ qù shèqū wèishēng fúwù zhōngxīn zuò ge xuānchuán ba?\\
« D'accord, tu as raison. Alors la semaine prochaine on va ensemble au centre de santé communautaire faire de la prévention ? »

\textbf{B :} 好主意！我们可以做海报，讲讲“防蚊、用药、早诊断”。\\
Hǎo zhǔyi! Wǒmen kěyǐ zuò hǎibào, jiǎngjiang “fáng wén, yòng yào, zǎo zhénduàn”.\\
« Bonne idée ! On peut faire des affiches, parler de "prévention moustiques, médicaments, diagnostic précoce". »
\end{textebox}

\courssub{Formules fonctionnelles de l'exposé et du débat}

\begin{aretenir}[Formules-outils pour l'oral (Exposé \& Débat)]
\begin{center}\begin{tabularx}{\linewidth}{|X|X|X|X|}\hline \rowcolor{popgreenL} \textbf{Fonction} & \textbf{Chinois} & \textbf{Pinyin} & \textbf{Français} \\ \hline
\multicolumn{4}{|c|}{\textbf{STRUCTURER UN EXPOSÉ}} \\ \hline
Introduction / Sujet & 大家好，今天我谈谈… & Dàjiā hǎo, jīntiān wǒ tántan… & Bonjour, aujourd'hui je parle de… \\ \hline
Premier argument & 首先 / 第一，… & Shǒuxiān / Dì yī,… & D'abord / Premièrement,… \\ \hline
Argument suivant & 其次 / 第二，… / 然后，… & Qícì / Dì èr,… / Ránhòu,… & Ensuite / Deuxièmement,… \\ \hline
Dernier argument & 最后 / 最重要的是，… & Zuìhòu / Zuì zhòngyào de shì,… & Enfin / Le plus important,… \\ \hline
Exemple personnel & 比如 / 举个例子，… & Bǐrú / Jǔ gè lìzi,… & Par exemple / Prenons un exemple,… \\ \hline
Conclusion & 总之 / 综上所述，… & Zǒngzhī / Zōngshàngshùshù,… & En bref / En conclusion,… \\ \hline
Appel à l'action & 让我们一起行动吧！ & Ràng wǒmen yīqǐ xíngdòng ba! & Agissons ensemble ! \\ \hline
\multicolumn{4}{|c|}{\textbf{EXPRIMER SON AVIS (DÉBAT)}} \\ \hline
Opinion personnelle & 我觉得 / 我认为 / 在我看来… & Wǒ juéde / Wǒ rènwéi / Zài wǒ kànlái… & Je pense que / À mon avis… \\ \hline
Accord total & 我完全同意。/ 赞成。 & Wǒ wánquán tóngyì. / Zànchéng. & Je suis tout à fait d'accord. / Pour. \\ \hline
Accord partiel & 有道理，但是… / 对，不过… & Yǒu dàoli, dànshì… / Duì, bùguò… & C'est vrai, mais… / Oui, cependant… \\ \hline
Désaccord poli & 我不太同意，因为… / 我怕不太对。 & Wǒ bù tài tóngyì, yīnwèi… / Wǒ pà bù tài duì. & Je ne suis pas trop d'accord, car… / Je crains que ce ne soit pas juste. \\ \hline
Argumentation & 因为…所以… / 原因是… & Yīnwèi… suǒyǐ… / Yuányīn shì… & Parce que… donc… / La raison est… \\ \hline
\multicolumn{4}{|c|}{\textbf{INTERAGIR ET RELANCER}} \\ \hline
Demander avis & 你怎么看？/ 你呢？ & Nǐ zěnme kàn? / Nǐ ne? & Qu'en penses-tu ? / Et toi ? \\ \hline
Demander précision & 这是什么意思？/ 能不能具体说说？ & Zhè shì shénme yìsi? / Néng bu néng jùtǐ shuōshuo? & Qu'est-ce que ça veut dire ? / Peux-tu préciser ? \\ \hline
Demander répéter & 请再说一遍。/ 我没听清。 & Qǐng zài shuō yíbiàn. / Wǒ méi tīngqīng. & Répétez s'il vous plaît. / Je n'ai pas bien entendu. \\ \hline
Céder la parole & 请说。/ 你先说。 & Qǐng shuō. / Nǐ xiān shuō. & Vas-y. / Toi d'abord. \\ \hline
\end{tabularx}\end{center}
\end{aretenir}

\courssub{Travail des tons et de la prononciation}

\begin{methode}[Maîtriser les tons critiques du thème « Environnement \& Santé »]
\textbf{Objectif :} Éviter les erreurs de tons qui changent le sens ou rendent le mot incompréhensible (critère éliminatoire au Bac oral).
\begin{enumerate}
    \item \textbf{Repérer les "mots-images" du thème :} Ce sont les mots-clés que le jury attend d'entendre correctement.
    \item \textbf{Travailler les paires minimales / confusions fréquentes :} Utiliser le tableau ci-dessous pour des exercices d'écoute discriminative (carte A/B) et de lecture en chaîne.
    \item \textbf{Appliquer le Sandhi (3-3 $\rightarrow$ 2-3) :} Systématiser sur \zh{保险}、\zh{好主意}、\zh{很危险}、\zh{很好}、\zh{马上}.
    \item \textbf{Intonation de la phrase :} Montée pour les questions (\zh{ma, ne, ba}), chute pour les affirmations. Particules finales en ton neutre.
\end{enumerate}
\end{methode}

\begin{propriete}[Tableau des tons pièges : Recyclage \& Paludisme]
\begin{center}\begin{tabularx}{\linewidth}{|X|X|X|X|}\hline \rowcolor{poporangeL} \textbf{Mot (Caractères)} & \textbf{Pinyin Standard} & \textbf{Sens} & \textbf{Piège tonal / Confusion fréquente} \\ \hline
垃圾 & lājī (1-1) & déchet & \textbf{Pas} \zh{làjǐ} (4-3) ni \zh{lājǐ} (1-3). 1er ton haut stable. \\ \hline
分类 & fēnlèi (1-4) & trier / classification & \zh{fēn} (1) haut ; \zh{lèi} (4) chute nette. Pas \zh{fènlèi} (4-4). \\ \hline
回收 & huíshōu (2-1) & recycler & \zh{huí} (2) monte ; \zh{shōu} (1) haut. \textbf{Pas} \zh{huǐshōu} (3-1 = détruire/regretter). \\ \hline
疟疾 & nüèjí (4-2) & paludisme & \zh{nüè} (4) chute forte sur \zh{üe} ; \zh{jí} (2) monte. \textbf{Mot difficile, à sur-apprendre}. \\ \hline
蚊子 & wénzi (2-neutre) & moustique & \zh{wén} (2) monte ; \zh{zi} ton neutre (court, bas). \textbf{Pas} \zh{wěnzi} (3). \\ \hline
蚊帐 & wénzhàng (2-4) & moustiquaire & \zh{zhàng} (4) chute. Différencier \zh{zhàng} (4, objet) vs \zh{zhǎng} (3, grandir/chef). \\ \hline
预防 & yùfáng (4-2) & prévenir & \zh{yù} (4) ; \zh{fáng} (2). \textbf{Pas} \zh{yǔfáng} (3-2 = pluie/langage). \\ \hline
确诊 & quèzhěn (4-3) & diagnostiquer & \zh{què} (4) ; \zh{zhěn} (3). \textbf{Pas} \zh{quēzhěn} (1-3 = manquer). \\ \hline
延误 & yánwù (2-4) & retarder & \zh{yán} (2) ; \zh{wù} (4). \textbf{Pas} \zh{yǎnwǔ} (3-3 = spectacle martial). \\ \hline
保险 & bǎoxiǎn (3-3) & sûr / assurance & \textbf{Règle 3-3 :} \zh{bǎo} $\rightarrow$ 2e ton (\zh{báo}) à l'oral ! \zh{báoxiǎn}. \\ \hline
\end{tabularx}\end{center}
\end{propriete}

\begin{exoresolu}[Exercices de discrimination et correction tonale]
\textbf{Énoncé :}
\begin{enumerate}
    \item \textbf{(Piège auditif)} Parmi les paires suivantes, laquelle a \textbf{exactement la même séquence tonale} (tons citation) ?
    \begin{enumerate}
        \item \zh{垃圾 (lājī)} -- \zh{蚊子 (wénzi)}
        \item \zh{回收 (huíshōu)} -- \zh{预防 (yùfáng)}
        \item \zh{分类 (fēnlèi)} -- \zh{蚊帐 (wénzhàng)}
        \item \zh{疟疾 (nüèjí)} -- \zh{确诊 (quèzhěn)}
    \end{enumerate}
    \item \textbf{Sandhi 3-3 :} Donnez la prononciation \textbf{réelle (à l'oral)} pour :
    \zh{保险} (bǎoxiǎn), \zh{好主意} (hǎo zhǔyi), \zh{很危险} (hěn wéixiǎn), \zh{马上} (mǎshàng).
    \item \textbf{Correction d'élève :} Corrigez les tons dans cette phrase : \zh{我觉得我们应该回收垃圾，预防疟疾。}
    \textit{Élève dit :} Wǒ juéde wǒmen yīnggāi \textbf{huǐshōu} lājī, \textbf{yǔfáng} \textbf{nüèjǐ}.
\end{enumerate}
\tcblower
\textbf{Solution détaillée :}
\textbf{1. Analyse des paires (Tons citation) :}
\begin{itemize}
    \item A : \zh{lājī} (1-1) vs \zh{wénzi} (2-neutre) $\rightarrow$ \textbf{Différents}.
    \item B : \zh{huíshōu} (2-1) vs \zh{yùfáng} (4-2) $\rightarrow$ \textbf{Différents}.
    \item C : \zh{fēnlèi} (1-4) vs \zh{wénzhàng} (2-4) $\rightarrow$ \textbf{Différents} (1ère syllabe : 1 vs 2).
    \item D : \zh{nüèjí} (4-2) vs \zh{quèzhěn} (4-3) $\rightarrow$ \textbf{Différents} (2ème syllabe : 2 vs 3).
    \item \textbf{Réponse :} \textbf{Aucune paire ne correspond.} C'est un test de vigilance.
\end{itemize}
\textbf{2. Sandhi 3-3 (Troisième ton + Troisième ton $\rightarrow$ 2e ton + 3e ton) :}
\begin{center}\begin{tabularx}{\linewidth}{|X|X|X|}\hline \rowcolor{popgreenL} Citation (Dictionnaire) & Prononciation Réelle (Sandhi) & Note \\ \hline
\zh{保险} (bǎoxiǎn) & \zh{báoxiǎn} (2-3) & 1er 3e ton $\rightarrow$ 2e ton \\ \hline
\zh{好主意} (hǎo zhǔyi) & \zh{háo zhǔyi} (2-3-4) & \zh{hǎo}(3)+\zh{zhǔ}(3)$\rightarrow$\zh{háo zhǔ} \\ \hline
\zh{很危险} (hěn wéixiǎn) & \zh{hén wéixiǎn} (2-2-3) & \zh{hěn}(3)+\zh{xiǎn}(3)$\rightarrow$\zh{hén} (\zh{wei} reste 2) \\ \hline
\zh{马上} (mǎshàng) & \zh{máshàng} (2-4) & \zh{mǎ}(3)+\zh{shàng}(4) $\rightarrow$ \zh{má} (règle 3-4 : 3e ton $\rightarrow$ 2e demi-ton haut) \\ \hline
\end{tabularx}\end{center}
\textbf{3. Correction de la phrase d'élève :}
\begin{itemize}
    \item \zh{huǐshōu} (3-1) $\rightarrow$ \textbf{\zh{huíshōu} (2-1)}. \textit{Erreur :} Confusion \zh{huí} (retourner) vs \zh{huǐ} (détruire/regretter).
    \item \zh{yǔfáng} (3-2) $\rightarrow$ \textbf{\zh{yùfáng} (4-2)}. \textit{Erreur :} Confusion \zh{yù} (prévenir) vs \zh{yǔ} (pluie/langage).
    \item \zh{nüèjǐ} (4-3) $\rightarrow$ \textbf{\zh{nüèjí} (4-2)}. \textit{Erreur :} \zh{jí} (maladie/urgent) est 2e ton, pas 3e ton (\zh{jǐ} = soi-même / combien).
\end{itemize}
\textbf{Phrase corrigée :} \zh{Wǒ juéde wǒmen yīnggāi huíshōu lājī, yùfáng nüèjí.}
\end{exoresolu}

\courssub{Jeux de rôle et débat guidé}

\begin{experience}[Jeu de rôle guidé : « Au point de tri sélectif du lycée »]
\textbf{Durée :} 10 min (Prépa 2 min + Jeu 3 min + Feedback 2 min + Restitution 3 min).
\textbf{Matériel :} Cartes "Rôle A / Rôle B", images de déchets (pile, bouteille sale, reste ndolé, papier, canette).
\begin{enumerate}
    \item \textbf{Distribution des rôles :}
    \begin{itemize}
        \item \textbf{Rôle A (Agent de tri) :} Objectif : Accueillir, identifier 3 erreurs dans le sac de B, donner la consigne exacte (couleur bac + geste : laver, aplatir), remercier.
        \item \textbf{Rôle B (Élève pressé) :} Objectif : Vider son sac vite fait, se faire corriger, s'excuser, demander où mettre la bouteille sale, promettre de faire mieux.
    \end{itemize}
    \item \textbf{Lexique de survie (affiché au tableau) :}
    \zh{这个属于…} (Ça appartient à…), \zh{请放进…} (Merci de mettre dans…), \zh{要洗干净/压扁} (Faut laver/aplatir), \zh{不好意思} (Désolé), \zh{下次注意} (Fait gaffe la prochaine fois), \zh{谢谢配合} (Merci coopération).
    \item \textbf{Déroulement :} Interdiction d'écrire. Si bloqué : \zh{请再说一遍}, \zh{这是什么意思？}, \zh{我没听清}.
    \item \textbf{Évaluation par les pairs (Grille simplifiée) :} Tons (mots-clés) / Vocabulaire précis (catégories) / Politesse (\zh{请/谢谢/不好意思}) / Tâche accomplie (3 erreurs corrigées).
\end{enumerate}
\end{experience}

\begin{exoresolu}[Simulation notée type Bac : « Au point de tri »]
\textbf{Énoncé :} Complétez le dialogue. A (Agent) corrige B sur 3 erreurs : pile (有害), bouteille plastique sale (可回收 mais sale), reste de ndolé (厨余). B s'excuse, demande pour la bouteille, s'engage.
\tcblower
\begin{textebox}[Jeu de rôle : 垃圾分类指导点]
\textbf{A (Agent) :} 同学，请把垃圾分类放好。这个黑色的袋子里装的是什么？\\
Tóngxué, qǐng bǎ lājī fēnlèi fàng hǎo. Zhège hēisè de dàizi lǐ zhuāng de shì shénme?\\
« Camarade, merci de bien trier les déchets. Qu'y a-t-il dans ce sac noir ? »

\textbf{B (Élève) :} 唔... 我有电池、喝完的塑料瓶、还有昨晚吃剩的炒面。\\
Wú... Wǒ yǒu diànchí, hē wán de sùliàopíng, hái yǒu zuówǎn chī shèng de chǎomiàn.\\
« Euh... j'ai des piles, une bouteille en plastique finie, et les nouilles sautées restantes d'hier soir. »

\textbf{A :} \textbf{不对。电池是有害垃圾，要放进红色桶。塑料瓶要洗干净、压扁，才能放进蓝色可回收桶。炒面是厨余垃圾，放绿色桶。}\\
Bùduì. Diànchí shì yǒuhài lājī, yào fàng jìn hóngsè tǒng. Sùliàopíng yào xǐ gānjìng, yā biǎn, cáinéng fàng jìn lánsè kěhuíshōu tǒng. Chǎomiàn shì chúyú lājī, fàng lǜsè tǒng.\\
« Non. Piles = déchets dangereux $\rightarrow$ bac rouge. Bouteille plastique : laver propre, aplatir, puis bac bleu recyclable. Nouilles = déchets cuisine $\rightarrow$ bac vert. »

\textbf{B :} \textbf{不好意思，我没注意。那这个脏塑料瓶我先拿去洗手池洗一下。下次我一定注意分类。}\\
Bùhǎoyìsi, wǒ méi zhùyì. Nà zhège zāng sùliàopíng wǒ xiān ná qù xǐshǒuchí xǐ yíxià. Xiàcì wǒ yídìng zhùyì fēnlèi.\\
« Désolé, je n'ai pas fait attention. Alors cette bouteille sale je vais la laver au lavabo. La prochaine fois je ferai gaffe au tri. »

\textbf{A :} \textbf{好，谢谢配合！记得洗手哦。}\\
Hǎo, xièxie pèihé! Jìde xǐshǒu o.\\
« Bien, merci de ta coopération ! N'oublie pas de te laver les mains. »
\end{textebox}
\textbf{Points clés validés (Grille Bac) :}
\begin{itemize}
    \item \textbf{Vocabulaire technique :} \zh{有害垃圾}, \zh{可回收物}, \zh{厨余垃圾}, \zh{压扁}, \zh{洗干净} (complément de résultat), \zh{才能} (seulement alors).
    \item \textbf{Structures :} \zh{要放进...} (obligation + direction), \zh{先...再...} (séquence), \zh{下次一定...} (engagement).
    \item \textbf{Politesse / Interaction :} \zh{不好意思}, \zh{谢谢配合}, \zh{哦} (particule finale adoucissante).
\end{itemize}
\end{exoresolu}

\begin{experience}[Débat guidé : « Paludisme : automédication vs Hôpital »]
\textbf{Consigne :} En groupes de 4. Deux "Pour l'hôpital systématique" (Rôle Médecin/Parent inquiet), Deux "Pour l'observation d'abord" (Rôle Étudiant/Parent occupé).
\textbf{Arguments imposés à caser :} \zh{验血}, \zh{确诊}, \zh{延误病情}, \zh{退烧药}, \zh{观察}, \zh{贵/慢}, \zh{不管轻重}, \zh{比较保险}.
\textbf{Arbitre :} Un 5e élève note les formules utilisées (Accord, Désaccord, Relance) et les tons des mots-clés (\zh{nüèjí, yùfáng, wénzi, wénzhàng}).
\textbf{Durée :} 5 min débat + 2 min débrief.
\end{experience}

\courssub{Expression écrite : Affiche et message}

\begin{methode}[Rédiger une affiche de prévention / un message court (50-80 caractères)]
\textbf{Objectif :} Produire un texte fonctionnel (affiche, SMS, WeChat) clair, visuel, grammaticalement correct.
\begin{enumerate}
    \item \textbf{Format \& Public :} Affiche scolaire = Ton impératif (\zh{请}, \zh{要}, \zh{别}), visuel, slogans 4 caractères. Message ami = Ton familier (\zh{吧}, \zh{呀}, \zh{哦}), particules modales.
    \item \textbf{Structures "clés en main" :}
    \begin{itemize}
        \item Titre : \zh{保护环境，从我做起！} / \zh{远离疟疾，健康生活！}
        \item Conseil + Résultat : \zh{请 + V + 好/干净，就能…} / \zh{为了…, 我们应该…}
        \item Interdiction : \zh{严禁 / 请勿 / 别 / 不要 + V}
        \item Signature : \zh{环保小组 发布} / \zh{卫生委员会} / \zh{联系电话：…}
    \end{itemize}
    \item \textbf{Ordre des mots canonique (Sujet - Temps - Lieu - Manière - Verbe - Complément) :}
    \zh{我们 (S) 在学校 (L) 每天 (T) 认真 (M) 分类垃圾 (VO).} $\checkmark$
    \item \textbf{Classificateurs obligatoires :} \zh{一顶/张 蚊帐}, \zh{一节/粒 电池}, \zh{一条 建议}, \zh{一张 海报}, \zh{一种 药}.
    \item \textbf{Relecture :} Caractères (clé, ordre traits), Ponctuation chinoise (， 。 ！), Pinyin mental.
\end{enumerate}
\end{methode}

\begin{exoresolu}[Rédaction d'une affiche : « Prévenir le paludisme, dès la lutte anti-moustiques »]
\textbf{Énoncé :} Affiche pour le hall du lycée (env. 60 caractères). Titre, 3 mesures (moustiquaire, répulsif, eau stagnante), slogan. Utiliser \zh{请}, \zh{不要}, \zh{要}.
\tcblower
\begin{textebox}[Affiche : 预防疟疾，从防蚊开始]
\textbf{预防疟疾，从防蚊开始}\\
Yùfáng nüèjí, cóng fáng wén kāishǐ\\
« Prévenir le paludisme, commence par la lutte anti-moustiques »

\textbf{晚上睡觉，请挂好蚊帐。}\\
Wǎnshang shuìjiào, qǐng guà hǎo wénzhàng.\\
« Le soir pour dormir, merci de bien mettre la moustiquaire. »

\textbf{外出涂抹驱蚊液，穿长袖衣服。}\\
Wàichū túmǒ qūwén yè, chuān chángxiù yīfu.\\
« En sortant, appliquez répulsif, portez manches longues. »

\textbf{清理积水，别让蚊子繁殖。}\\
Qīnglǐ jīshuǐ, bié ràng wénzi fánzhí.\\
« Éliminez l'eau stagnante, ne laissez pas les moustiques proliférer. »

\textbf{健康校园，你我共建！}\\
Jiànkāng xiàoyuán, nǐ wǒ gòngjiàn!\\
« Campus sain, ensemble on le construit ! »
\end{textebox}
\textbf{Justifications grammaticales :}
\begin{itemize}
    \item \zh{挂好} : \zh{好} = complément de résultat (action menée à bien).
    \item \zh{涂抹 驱蚊液} : Vocabulaire précis HSK 4.
    \item \zh{别让} : Impératif négatif + causatif (\zh{让} = faire en sorte que).
    \item \zh{繁殖} : Verbe intransitif (sujet = moustiques).
    \item \zh{你我共建} : Structure parallélique 4 caractères (style \zh{成语}), standard pour slogans.
\end{itemize}
\end{exoresolu}

\courssub{Traduction : Thème (Français $\rightarrow$ Chinois)}

\begin{methode}[Méthode du Thème : Français $\rightarrow$ Chinois]
\textbf{Objectif :} Éviter le calque syntaxique. Respecter la topicalisation (Thème-Rhème) et l'ordre canonique S - T - L - M - V - C.
\begin{enumerate}
    \item \textbf{Segmenter} en unités de sens.
    \item \textbf{Identifier le Sujet Chinois (Thème) :} Souvent le sujet du discours, pas toujours le sujet grammatical FR. Ex: \textit{« Ce déchet, on le recycle »} $\rightarrow$ \zh{这个垃圾, 回收利用就行。}
    \item \textbf{Placer les compléments :} \textbf{Temps} $\rightarrow$ \textbf{Lieu} $\rightarrow$ \textbf{Manière/Instrument} $\rightarrow$ \textbf{Verbe}.
    \item \textbf{Traiter les pièges grammaticaux :}
    \begin{itemize}
        \item \textbf{Passif / Focus agent :} \textit{« Le palu est transmis par les moustiques »} $\nrightarrow$ \zh{被} (trop formel) $\rightarrow$ \zh{蚊子传播疟疾} (Sujet agent) OU \zh{疟疾是蚊子传播的} (\zh{是...的} focus passé).
        \item \textbf{Comparatif :} \textit{« Trier est plus utile que jeter »} $\rightarrow$ \zh{分类比乱扔有用} (\zh{比} structure).
        \item \textbf{Aspect achevé / Expérience :} \textit{« J'ai déjà trié »} $\rightarrow$ \zh{我已经分类了} (\zh{已经...了}) OU \zh{我分类过} (\zh{过}).
        \item \textbf{Durée / Fréquence :} \textit{« Je dors sous moustiquaire chaque nuit »} $\rightarrow$ \zh{我每天晚上都睡蚊帐} (\zh{都} + Adv fréquence avant V).
        \item \textbf{Structure \zh{把} :} Action de manipulation/placement sur objet connu. \textit{« Jette la poubelle dans le bac »} $\rightarrow$ \zh{把垃圾扔进垃圾桶}.
    \end{itemize}
    \item \textbf{Vérifier classificateurs :} \textit{« Trois types de poubelles »} $\rightarrow$ \zh{三种垃圾桶} (pas \zh{三个}).
\end{enumerate}
\end{methode}

\begin{exoresolu}[Thème : 5 phrases à traduire (Niveau Bac)]
\textbf{Énoncé :} Traduisez en chinois (caractères + pinyin). Attention : ordre mots, classificateurs, \zh{比}, \zh{是...的}, \zh{了/过}, \zh{把}.
\begin{enumerate}
    \item Hier, j'ai trié tous les déchets de la cuisine.
    \item Le recyclage est plus utile que de jeter tout ensemble.
    \item Ce moustique a été attrapé par mon petit frère.
    \item Nous devons dormir sous une moustiquaire imprégnée chaque nuit.
    \item Le médecin m'a conseillé de prendre ce médicament pendant trois jours.
\end{enumerate}
\tcblower
\textbf{Solution détaillée :}
\begin{enumerate}
    \item \textbf{Hier, j'ai trié tous les déchets de la cuisine.}
    \zh{昨天，我把厨房的垃圾全部分类了。}\\
    Zuótiān, wǒ \textbf{bǎ} chúfáng de lājī quánbù fēnlèi \textbf{le}.\\
    \textit{Analyse :} \zh{把} structure (objet défini \zh{厨房的垃圾} + manipulation/résultat \zh{分类了}). \zh{全都} avant V. \zh{了} final = action achevée passée.
    \item \textbf{Le recyclage est plus utile que de jeter tout ensemble.}
    \zh{回收利用比乱扔垃圾有用多了。}\\
    Huíshōu lìyòng \textbf{bǐ} luàn rēng lājī yǒuyòng \textbf{duō le}.\\
    \textit{Analyse :} \zh{比} structure. \zh{有用} (adj) + \zh{多了} (beaucoup plus). \zh{乱扔} (jeter n'importe comment) = verbe composé.
    \item \textbf{Ce moustique a été attrapé par mon petit frère.}
    \zh{这只蚊子是我弟弟抓住的。}\\
    Zhè \textbf{zhī} wénzi \textbf{shì} wǒ dìdi zhuāzhù \textbf{de}.\\
    \textit{Analyse :} Classificateur \zh{只} (insectes). Structure \zh{是...的} focus sur agent (\zh{我弟弟}) action passée accomplie. \zh{抓住} (V + compl. résultat).
    \item \textbf{Nous devons dormir sous une moustiquaire imprégnée chaque nuit.}
    \zh{我们每天晚上都得在药蚊帐里睡觉。}\\
    Wǒmen \textbf{měitiān wǎnshang dōu} děi \textbf{zài yào wénzhàng lǐ} shuìjiào.\\
    \textit{Analyse :} Ordre T-L : \zh{每天晚上} (T) + \zh{都} (portée) + \zh{得} (devoir fort). Locatif \zh{在...里} placé AVANT verbe \zh{睡觉} (verbe statique/déplacement volontaire). \zh{药蚊帐} (MILD).
    \item \textbf{Le médecin m'a conseillé de prendre ce médicament pendant trois jours.}
    \zh{医生建议我吃这个药，吃三天。}\\
    Yīshēng jiànyì wǒ chī zhège yào, chī \textbf{sān tiān}.\\
    \textit{Analyse :} \zh{建议} + Objet (\zh{我}) + Prop. Durée \zh{三天} placée APRÈS le verbe (répétition \zh{吃}) pour clarté. Structure standard : V + Durée (+ Objet) ou V + Objet, V + Durée.
\end{enumerate}
\end{exoresolu}

\courssub{Éléments socioculturels : Cameroun / Chine}

\begin{savaistu}[Gestion des déchets et Paludisme : Regards croisés Cameroun - Chine]
\textbf{1. Tri des déchets (垃圾分类) :}
\begin{itemize}
    \item \textbf{Cameroun (Yaoundé, Douala) :} Déploiement progressif du tri à la source (4 flux : \zh{可回收物} bleu, \zh{厨余垃圾} vert, \zh{有害垃圾} rouge, \zh{其他垃圾} gris) par HYSACAM / municipalités. Défis : sensibilisation, infrastructures de collecte séparée, secteur informel (trieurs à la décharge de Nkolfoulou / PK10). Le \zh{ndolé} (feuilles de bitterleaf) et restes de \zh{manioc}/\zh{plantain} forment l'essentiel des \zh{厨余垃圾}.
    \item \textbf{Chine (Shanghai, Pékin, Shenzhen) :} Tri obligatoire strict depuis 2019 (Shanghai) / 2021 (national). 4 catégories similaires mais noms légèrement différents : \zh{可回收物} (bleu), \zh{有害垃圾} (rouge), \zh{厨余垃圾/湿垃圾} (vert), \zh{干垃圾/其他垃圾} (gris). Amendes pour particuliers (200 CNY) et entreprises (50 000 CNY). Applications \zh{小程序} (WeChat/Alipay) pour scanner le code déchet. Taux de tri > 90\% à Shanghai.
    \item \textbf{Point commun :} Responsabilité individuelle (\zh{人人有责}) et éducation scolaire (cours \zh{垃圾分类} dès le primaire).
\end{itemize}

\textbf{2. Paludisme (疟疾) :}
\begin{itemize}
    \item \textbf{Cameroun :} Endémie stable, 1ère cause de consultation (30-40\%). \zh{Anopheles gambiae} vecteur principal. Résistance aux insecticides (pyréthrinoïdes) et à la chloroquine (d'où ACT : \zh{青蒿素} \zh{联合疗法}). Automédication fréquente (\zh{自行用药}) : \zh{街边药店} (pharmacies de rue), tisanes traditionnelles (\zh{草药}). Campagnes \zh{MILD} (Moustiquaires Imprégnées de Longue Durée) gratuites (distribution massive 2023).
    \item \textbf{Chine :} \textbf{Certifiée "sans paludisme" par l'OMS en juin 2021} (0 cas autochtone depuis 2017). Stratégie "1-3-7" : Déclaration 1j, Enquête 3j, Réponse foyer 7j. Utilisation historique de la \zh{青蒿素} (Artémisine, Prix Nobel Tu Youyou 2015). Aujourd'hui : surveillance aux frontières (Yunnan) pour cas importés (\zh{输入性病例}).
    \item \textbf{Coopération :} Équipes médicales chinoises au Cameroun (hôpital de l'Amitié, centres de santé). Partage d'expérience sur la \zh{青蒿素} et la gestion des cas.
\end{itemize}
\end{savaistu}

\courssub{Exercices d'entraînement (Travail personnel / Devoir)}

\begin{exercice}[Exercice 1 : Discrimination tonale et Sandhi]
\begin{enumerate}
    \item Entourez le mot dont le ton est \textbf{incorrect} dans chaque groupe :
    \begin{enumerate}
        \item \zh{垃圾} \zh{分类} \zh{回收} \zh{利用}
        \item \zh{疟疾} \zh{蚊子} \zh{蚊帐} \zh{预防}
        \item \zh{确诊} \zh{延误} \zh{保险} \zh{建议}
    \end{enumerate}
    \item Appliquez le Sandhi 3-3 et écrivez le pinyin \textbf{réel (oral)} :
    \zh{很好} (hěn hǎo), \zh{马上} (mǎ shàng), \zh{哪里} (nǎ lǐ), \zh{可以} (kě yǐ).
\end{enumerate}
\end{exercice}

\begin{exercice}[Exercice 2 : Compléter le débat (Expression orale écrite)]
\textbf{Contexte :} A et B discutent de la meilleure prévention contre le palu. Complétez les répliques de B (caractères + pinyin). Utilisez : \zh{我觉得}, \zh{因为...所以...}, \zh{建议}, \zh{不管...都...}.
\begin{textebox}[À compléter]
\textbf{A :} 听说雨季快到了，怎么预防疟疾最有效？
\textbf{B :} \underline{\hspace{8cm}} (Mon avis : moustiquaire + répulsif + eau stagnante).
\textbf{A :} 可是挂蚊帐很热，涂驱蚊液麻烦。
\textbf{B :} \underline{\hspace{8cm}} (Désaccord poli : santé > confort. Parce que palu dangereux).
\textbf{A :} 那如果已经被咬了怎么办？
\textbf{B :} \underline{\hspace{8cm}} (Conseil : peu importe symptômes, aller hôpital vérifier).
\end{textebox}
\end{exercice}

\begin{exercice}[Exercice 3 : Traduction Thème (Français $\rightarrow$ Chinois)]
Traduisez en chinois (caractères + pinyin).
\begin{enumerate}
    \item Pour protéger l'environnement, nous devons trier les déchets chaque jour.
    \item Cette bouteille en plastique est sale, il faut la laver et l'aplatir avant de la recycler.
    \item Le petit frère a attrapé le paludisme parce qu'il n'a pas dormi sous la moustiquaire.
    \item Le médecin lui a conseillé de prendre le médicament pendant cinq jours.
\end{enumerate}
\end{exercice}

\begin{corrige}[Exercices 1, 2, 3]
\textbf{Exercice 1 :}
\begin{enumerate}
    \item \zh{利用} (lìyòng, 4-4) -- Les autres : \zh{lājī}(1-1), \zh{fēnlèi}(1-4), \zh{huíshōu}(2-1). \textit{Note :} Pas d'erreur de ton \textit{citation} ici, mais \zh{利用} est 4-4. Si la question porte sur le Sandhi : \zh{回收利用} $\rightarrow$ \zh{huíshōu lìyòng} (pas de sandhi 3-3).
    \item \zh{蚊子} (wénzi, 2-neutre) -- Les autres : \zh{nüèjí}(4-2), \zh{wénzhàng}(2-4), \zh{yùfáng}(4-2). \zh{蚊子} est le seul avec ton neutre.
    \item \zh{建议} (jiànyì, 4-4) -- Les autres : \zh{quèzhěn}(4-3), \zh{yánwù}(2-4), \zh{bǎoxiǎn}(3-3$\rightarrow$2-3).
    \item \textbf{Sandhi 3-3 réel :}
    \zh{很好} $\rightarrow$ \zh{hén hǎo} (2-3) ; \zh{马上} $\rightarrow$ \zh{máshàng} (2-4, règle 3-4) ; \zh{哪里} $\rightarrow$ \zh{nálǐ} (2-3, \zh{nǎ} 3e ton isol$\rightarrow$2e devant 3e) ; \zh{可以} $\rightarrow$ \zh{kěyǐ} (3-3, \zh{kě} 3e ton isol$\rightarrow$2e devant 3e ? \textbf{Non !} \zh{kě} est 3e ton mais \zh{yǐ} est 3e ton. \zh{kěyǐ} $\rightarrow$ \zh{kéyǐ} (2-3)). \textit{Correction :} \zh{kěyǐ} se prononce \zh{kéyǐ} (2-3).
\end{enumerate}

\textbf{Exercice 2 (Proposition) :}
\begin{textebox}[Corrigé débat]
\textbf{B :} \textbf{我觉得最有效的是：挂蚊帐、涂驱蚊液、清理积水。}\\
Wǒ juéde zuì yǒuxiào de shì: guà wénzhàng, tú qūwén yè, qīnglǐ jīshuǐ.
\textbf{B :} \textbf{我不太同意。因为疟疾很危险，所以健康比舒服重要。不管多麻烦，都要预防。}\\
Wǒ bù tài tóngyì. Yīnwèi nüèjí hěn wéixiǎn, suǒyǐ jiànkāng bǐ shūfu zhòngyào. Bùguǎn duō máfan, dōu yào yùfáng.
\textbf{B :} \textbf{建议你：不管有没有症状，都去医院验血确诊，比较保险。}\\
Jiànyì nǐ: bùguǎn yǒu méiyǒu zhèngzhuàng, dōu qù yīyuàn yànxiě quèzhěn, bǐjiào bǎoxiǎn.
\end{textebox}

\textbf{Exercice 3 :}
\begin{enumerate}
    \item \zh{为了保护环境，我们每天都要分类垃圾。} (Wèile bǎohù huánjìng, wǒmen měitiān dōu yào fēnlèi lājī.)
    \item \zh{这个塑料瓶很脏，得先洗干净、压扁，才能回收。} (Zhège sùliàopíng hěn zāng, děi xiān xǐ gānjìng, yā biǎn, cáinéng huíshōu.)
    \item \zh{弟弟因为没睡蚊帐，得了疟疾。} (Dìdi yīnwèi méi shuì wénzhàng, déle nüèjí.) \textit{ou} \zh{是因为没睡蚊帐得的疟疾。} (\zh{是...的} focus cause).
    \item \zh{医生建议他吃这个药，吃五天。} (Yīshēng jiànyì tā chī zhège yào, chī wǔ tiān.)
\end{enumerate}
\end{corrige}

\courssub{Erreurs fréquentes et pièges du Bac (Leçon 15)}

\begin{attention}[Les 5 erreurs qui coûtent des points au Bac (Leçon 15)]
\begin{enumerate}
    \item \textbf{Tons : Confusion \zh{nüèjí} (疟疾) vs \zh{nüèjǐ} / \zh{yùfáng} (预防) vs \zh{yǔfáng} (雨/语).}
    Le mot \zh{疟疾} (4-2) et \zh{预防} (4-2) sont des "mots-images" du thème. Une erreur de ton sur la 2e syllabe (\zh{jí} vs \zh{jǐ}, \zh{fáng} vs \zh{fǎng}) rend le mot incompréhensible. \textbf{Réflexe :} S'entraîner sur paires minimales \zh{jí/jǐ}, \zh{fáng/fǎng} avant l'oral.
    \item \textbf{Ordre des mots : Complément de lieu AVANT le verbe (sauf verbes de placement avec \zh{把}).}
    \textit{Faux :} \zh{我睡在蚊帐里觉。} / \zh{我扔垃圾在垃圾桶里。} \\
    \textit{Vrai :} \zh{我睡蚊帐。} / \zh{我在蚊帐里睡觉。} (Lieu statique avant V). \\
    \textit{Vrai (Action placement) :} \zh{我把垃圾扔在垃圾桶里。} / \zh{把垃圾放进绿色桶。} (\zh{把} + Objet + \zh{在/进} + Lieu + V). \textbf{Règle :} Lieu \zh{在...} avant V pour position statique (\zh{我在学校学习}) ; \zh{把} + Objet + \zh{在/进} + Lieu + V pour action de déplacement/manipulation.
    \item \textbf{Classificateurs : \zh{个} "fourre-tout" vs classificateurs spécifiques (niveau HSK 4 requis).}
    \textit{Faux :} \zh{一个蚊帐, 一个垃圾, 一个建议, 一个电池。} \\
    \textit{Vrai :} \zh{一顶/一张 蚊帐}, \zh{一堆/一袋 垃圾}, \zh{一条 建议}, \zh{一节/一粒 电池}, \zh{一只 蚊子}. Au Bac, l'usage correct du classificateur montre le niveau HSK 4.
    \item \textbf{Mots-outils : \zh{了} (aspect perfectif/changement d'état) vs \zh{过} (expérience) vs \zh{着} (état durable).}
    \textit{Contexte :} « J'ai déjà trié les déchets. » $\rightarrow$ \zh{我已经分类\cle{了} / 分类\cle{过}了。} (Action achevée / Expérience passée). \\
    « La poubelle est pleine (maintenant). » $\rightarrow$ \zh{垃圾桶满\cle{了}。} (Changement d'état). \\
    « Il porte un masque (en ce moment). » $\rightarrow$ \zh{他戴\cle{着}口罩。} (État durable). Ne pas écrire \zh{戴了} (action ponctuelle passée).
    \item \textbf{Calque français : "Éviter les moustiques" $\nrightarrow$ \zh{避免蚊子}.}
    \textit{Faux :} \zh{我们要避免蚊子。} (On évite les moustiques eux-mêmes ? Sens bizarre). \\
    \textit{Vrai :} \zh{我们要防蚊 / 防止被蚊子叮咬 / 驱蚊。} (Prévenir les moustiques / Éviter de se faire piquer / Chasser les moustiques). \textbf{Collocations VO :} \zh{防蚊}, \zh{灭蚊}, \zh{驱蚊}, \zh{挂蚊帐}. Ne pas traduire "éviter/prévenir" par \zh{避免/预防} + [nom concret animé] sans verbe d'action spécifique.
\end{enumerate}
\end{attention}

\newpage

\coursec{Exercices}

\courssub{Je vérifie mes connaissances}

\begin{exercice}[QCM : vocabulaire et structures du thème]
Choisissez la bonne réponse parmi les trois propositions. Une seule réponse est correcte.
\begin{enumerate}
    \item Comment dit-on « \textbf{moustiquaire} » en chinois ?
    \begin{enumerate}
        \item \zh{蚊子} (wénzi)
        \item \zh{蚊帐} (wénzhàng)
        \item \zh{蚊香} (wénxiāng)
    \end{enumerate}
    \item Quelle structure permet d'exprimer « \textbf{il faut / on doit} » pour donner un conseil ?
    \begin{enumerate}
        \item \zh{可以} (kěyǐ)
        \item \zh{应该} (yīnggāi)
        \item \zh{想要} (xiǎng yào)
    \end{enumerate}
    \item Le mot \zh{回收} (huíshōu) signifie :
    \begin{enumerate}
        \item polluer
        \item recycler / récupérer
        \item jeter
    \end{enumerate}
    \item Pour introduire un argument dans un exposé, on utilise souvent :
    \begin{enumerate}
        \item \zh{总之} (zǒngzhī)
        \item \zh{首先} (shǒuxiān)
        \item \zh{但是} (dànshì)
    \end{enumerate}
\end{enumerate}
\end{exercice}

\begin{exercice}[Vrai ou Faux : justifiez en chinois ou en français]
Déterminez si les affirmations suivantes sont vraies (V) ou fausses (F). Justifiez votre réponse par une phrase ou un mot-clé du cours.
\begin{enumerate}
    \item \zh{疟疾} (nuèjí) se transmet principalement par l'eau sale.
    \item \zh{垃圾分类} (lājī fēnlèi) est une mesure efficace pour faciliter le recyclage.
    \item Dans un débat, \zh{我不同意} (wǒ bù tóngyì) s'utilise pour approuver l'opinion de l'autre.
    \item Le symptôme \zh{发烧} (fāshāo) est fréquent chez les malades du paludisme.
    \item \zh{塑料袋} (sùliào dài) appartient à la catégorie des déchets recyclables (\zh{可回收物}) dans tous les systèmes de tri.
\end{enumerate}
\end{exercice}

\begin{exercice}[Vocabulaire : associez le caractère, le pinyin et le sens]
Reliez chaque élément de la colonne A (caractères) à son pinyin (colonne B) et à sa traduction (colonne C). Écrivez les numéros correspondants (ex: 1 - b - iii).

\begin{center}
\begin{tabularx}{\linewidth}{|X|X|X|X|}
\hline
\rowcolor{popblueL} \textbf{N°} & \textbf{Colonne A : Caractères} & \textbf{Colonne B : Pinyin} & \textbf{Colonne C : Traduction} \\
\hline
1 & \zh{预防} & a. fāshāo & i. moustiquaire \\
\hline
2 & \zh{蚊帐} & b. wénzhàng & ii. fièvre \\
\hline
3 & \zh{发烧} & c. yùfáng & iii. prévention \\
\hline
4 & \zh{污染} & d. wūrǎn & iv. pollution \\
\hline
5 & \zh{保护} & e. bǎohù & v. protéger / protection \\
\hline
\end{tabularx}
\end{center}
\end{exercice}

\begin{exercice}[Pinyin et tons : écoutez et identifiez]
Votre professeur lit les dix mots ci-dessous. Écrivez le pinyin complet avec les marques de tons (ex: \emph{huíshōu}) et entourez le ton du premier caractère (1, 2, 3 ou 4).
\begin{enumerate}
    \item \zh{疟疾} \hfill Pinyin : \underline{\hspace{4cm}} \hfill Ton 1er car. : 1 / 2 / 3 / 4
    \item \zh{回收} \hfill Pinyin : \underline{\hspace{4cm}} \hfill Ton 1er car. : 1 / 2 / 3 / 4
    \item \zh{垃圾} \hfill Pinyin : \underline{\hspace{4cm}} \hfill Ton 1er car. : 1 / 2 / 3 / 4
    \item \zh{预防} \hfill Pinyin : \underline{\hspace{4cm}} \hfill Ton 1er car. : 1 / 2 / 3 / 4
    \item \zh{发烧} \hfill Pinyin : \underline{\hspace{4cm}} \hfill Ton 1er car. : 1 / 2 / 3 / 4
    \item \zh{蚊帐} \hfill Pinyin : \underline{\hspace{4cm}} \hfill Ton 1er car. : 1 / 2 / 3 / 4
    \item \zh{塑料} \hfill Pinyin : \underline{\hspace{4cm}} \hfill Ton 1er car. : 1 / 2 / 3 / 4
    \item \zh{污染} \hfill Pinyin : \underline{\hspace{4cm}} \hfill Ton 1er car. : 1 / 2 / 3 / 4
    \item \zh{保护} \hfill Pinyin : \underline{\hspace{4cm}} \hfill Ton 1er car. : 1 / 2 / 3 / 4
    \item \zh{医院} \hfill Pinyin : \underline{\hspace{4cm}} \hfill Ton 1er car. : 1 / 2 / 3 / 4
\end{enumerate}
\end{exercice}

\courssub{Je m'entraîne}

\begin{exercice}[Traduction : du français vers le chinois (Thème)]
Traduisez les phrases suivantes en chinois (caractères + pinyin). Utilisez le vocabulaire de la leçon.
\begin{enumerate}
    \item Nous devons trier les ordures pour protéger l'environnement.
    \item Le paludisme se prévient en dormant sous une moustiquaire.
    \item À mon avis, le recyclage du plastique est très important.
    \item Le médecin me conseille de boire plus d'eau et de me reposer.
    \item Premièrement, il faut réduire les déchets ; deuxièmement, il faut recycler.
\end{enumerate}
\end{exercice}

\begin{exercice}[Transformation de phrases : structures de l'exposé et du débat]
Transformez les phrases selon l'indication entre parenthèses. Écrivez le résultat en caractères et en pinyin.
\begin{enumerate}
    \item \zh{我们要保护环境。} (Ajoutez « \textbf{我认为} » au début pour exprimer l'opinion)
    \item \zh{回收垃圾很有用。} (Remplacez « \zh{很有用} » par « \zh{非常重要} »)
    \item \zh{他发烧了。} (Ajoutez « \zh{因为} » + « \zh{他得了疟疾} » au début pour donner la cause)
    \item \zh{你应该睡蚊帐。} (Transformez en question avec « \zh{你觉得...怎么样？} »)
    \item \zh{塑料污染很严重。} (Ajoutez « \zh{总之} » au début pour conclure)
\end{enumerate}
\end{exercice}

\begin{exercice}[Texte à trous : dialogue à l'hôpital (Paludisme)]
Complétez le dialogue entre un médecin (\zh{医生}) et un patient (\zh{病人}) avec les mots manquants (caractères ou pinyin). Les premiers caractères sont donnés comme indice.

\begin{textebox}[Dialogue à l'hôpital]
\textbf{医生}：你怎么了？\quad
\textbf{病人}：医生，我头疼，还 f\_\_\_\_\_ (fièvre)。\quad
\textbf{医生}：你有没有其他的症状？\quad
\textbf{病人}：我 y\_\_\_\_\_ (aussi) froid, je transpire.\quad
\textbf{医生}：可能是 \_\_\_\_\_ (paludisme). Il faut faire un test de sang.\quad
\textbf{病人}：Ah ? Je dois \_\_\_\_\_ (rester) à l'hôpital ?\quad
\textbf{医生}：Oui, tu \_\_\_\_\_ (dois) prendre des médicaments et bien te reposer. Et à la maison, tu \_\_\_\_\_ (dois) dormir sous une moustiquaire.\quad
\textbf{病人}：D'accord, merci docteur.
\end{textebox}
\end{exercice}

\begin{exercice}[Remise en ordre : structure d'un exposé sur le recyclage]
Remettez les phrases suivantes dans l'ordre logique pour former un exposé cohérent (Salutation $\rightarrow$ Annonce du sujet $\rightarrow$ Arguments $\rightarrow$ Conclusion $\rightarrow$ Remerciement). Numérotez-les de 1 à 6.

\begin{enumerate}
    \item \zh{总之，保护环境从我做起。}
    \item \zh{大家好，今天我的话题是垃圾回收。}
    \item \zh{谢谢大家！}
    \item \zh{首先，回收可以减少污染。}
    \item \zh{其次，回收能节约资源。}
    \item \zh{各位同学，各位老师，大家早上好。}
\end{enumerate}
\end{exercice}

\courssub{Je me prépare au Bac}

\begin{exercice}[Type Bac — Compréhension orale/écrite : Exposé sur l'environnement]
Lisez le texte original ci-dessous (caractères + pinyin), puis répondez aux questions en français.

\begin{textebox}[Exposé : 我们应该怎么保护环境？]
各位同学，大家好。\\
今天，我想谈谈环境保护。\\
首先，我们要减少使用塑料袋，因为塑料污染很严重，分解需要几百年。\\
其次，我们应该做好垃圾分类。可回收物，比如纸、塑料、金属，要放进蓝色垃圾桶；厨余垃圾放进绿色垃圾桶。\\
再次，节约水电也是保护环境。离开房间要关灯，洗手要关水龙头。\\
总之，保护环境人人有责，从我做起，从小事做起。\\
谢谢大家！
\end{textebox}

\textbf{Questions de compréhension :}
\begin{enumerate}
    \item Quel est le sujet principal de l'exposé ?
    \item Citez deux arguments développés par l'orateur pour protéger l'environnement (en français).
    \item Selon le texte, où doit-on jeter les déchets de cuisine (\zh{厨余垃圾}) ?
    \item Quelle formule fonctionnelle l'orateur utilise-t-il pour conclure ?
\end{enumerate}

\textbf{Questions de langue :}
\begin{enumerate}
    \setcounter{enumi}{4}
    \item Relevez dans le texte deux connecteurs logiques (mots de liaison) et donnez leur sens en français.
    \item Trouvez dans le texte le mot qui signifie « \textbf{robinet} » et écrivez son pinyin.
    \item Transformez la phrase « \zh{我们要减少使用塑料袋} » en phrase négative avec « \zh{不要} ».
\end{enumerate}
\end{exercice}

\begin{exercice}[Type Bac — Expression orale guidée : Exposé et Débat]
\textbf{Consigne :} Préparez un exposé oral d'\textbf{environ 1 minute} (environ 150--200 caractères chinois) sur le thème : \zh{我们应该怎么预防疟疾？} (Comment devons-nous prévenir le paludisme ?).

\textbf{Contraintes obligatoires :}
\begin{itemize}
    \item Utilisez au moins \textbf{trois formules fonctionnelles} différentes (ex: \zh{首先}、\zh{其次}、\zh{总之}、\zh{我认为}、\zh{建议}...).
    \item Utilisez au moins \textbf{deux connecteurs logiques} (ex: \zh{因为...所以...}、\zh{虽然...但是...}、\zh{如果...就...}).
    \item Intégrez au moins \textbf{cinq mots de vocabulaire} de la leçon (ex: \zh{蚊帐}、\zh{蚊子}、\zh{发烧}、\zh{医院}、\zh{预防}、\zh{药}、\zh{环境}).
    \item Structure : Salutation $\rightarrow$ Annonce du sujet $\rightarrow$ 2--3 arguments $\rightarrow$ Conclusion $\rightarrow$ Remerciement.
\end{itemize}

\textbf{Étape 1 (Écrite - 10 min) :} Rédigez le texte de votre exposé en caractères + pinyin sur votre copie.
\textbf{Étape 2 (Orale - 1 min) :} Présentez votre exposé devant la classe sans lire, en soignant les tons et l'intonation.
\end{exercice}

\begin{exercice}[Type Bac — Thème et Version : Recyclage et Santé]
\textbf{A. Thème (Français $\rightarrow$ Chinois) :} Traduisez le paragraphe suivant en chinois (caractères + pinyin). \textit{Attendu : environ 80--100 caractères.}

« Au Cameroun, la pollution par les déchets plastiques est un grand problème. Pour protéger notre environnement, le gouvernement encourage le tri des ordures. Chaque élève doit participer : il faut mettre les bouteilles en plastique dans la poubelle bleue et les déchets de cuisine dans la poubelle verte. Si tout le monde fait un effort, notre ville sera plus propre. »

\textbf{B. Version (Chinois $\rightarrow$ Français) :} Traduisez le dialogue suivant en français naturel.

\begin{textebox}[Dialogue : Débat sur l'interdiction du plastique]
\textbf{记者}：同学，你觉得应该禁止塑料袋吗？
\textbf{学生}：我同意。因为塑料袋很难分解，污染土壤和河流。
\textbf{记者}：但是塑料袋很方便，便宜。
\textbf{学生}：虽然方便，但我们可以用布袋。布袋可以重复使用，环保。
\textbf{记者}：你的建议很好。总之，保护环境需要大家努力。
\textbf{学生}：对，从我做起，少用塑料袋。
\end{textebox}
\end{exercice}

\newpage

\coursec{Corrigés des exercices}

\courssub{Je vérifie mes connaissances}

\begin{corrige}[Exercice 1 — QCM : vocabulaire et structures du thème]
\textbf{1. Réponse 2 :} \zh{蚊帐} (wénzhàng) = moustiquaire. \\
\zh{蚊子} (wénzi) = moustique (l'insecte). \zh{蚊香} (wénxiāng) = spirale anti-moustiques.

\textbf{2. Réponse 2 :} \zh{应该} (yīnggāi) = devoir / devrait (obligation morale, conseil). \\
\zh{可以} (kěyǐ) = pouvoir (permission, possibilité). \zh{想要} (xiǎng yào) = vouloir.

\textbf{3. Réponse 2 :} \zh{回收} (huíshōu) = recycler / récupérer. \\
\zh{污染} (wūrǎn) = polluer. \zh{扔} (rēng) = jeter.

\textbf{4. Réponse 2 :} \zh{首先} (shǒuxiān) = premièrement / tout d'abord (introduit le premier argument). \\
\zh{总之} (zǒngzhī) = en conclusion. \zh{但是} (dànshì) = mais (opposition).
\end{corrige}

\begin{corrige}[Exercice 2 — Vrai ou Faux : justifiez]
\textbf{1. Faux (F).} Justification : \zh{疟疾} (paludisme) se transmet par la piqûre de \zh{蚊子} (moustiques anophèles femelles), pas par l'eau sale.

\textbf{2. Vrai (V).} Justification : \zh{垃圾分类} (tri sélectif) sépare les \zh{可回收物} (déchets recyclables) des autres déchets, ce qui facilite le \zh{回收} (recyclage).

\textbf{3. Faux (F).} Justification : \zh{我不同意} (wǒ bù tóngyì) signifie « Je ne suis pas d'accord » (désaccord). Pour approuver : \zh{我同意} (wǒ tóngyì).

\textbf{4. Vrai (V).} Justification : \zh{发烧} (fāshāo, fièvre) est le symptôme principal et typique du \zh{疟疾} (accès fébriles).

\textbf{5. Faux (F).} Justification : La classification varie selon les villes (ex: Shanghai, Pékin). Souvent, les \zh{塑料袋} sales sont des \zh{干垃圾} (autres déchets) ou \zh{其他垃圾}, seuls les plastiques propres sont \zh{可回收物}. Dire « dans tous les systèmes » est faux.
\end{corrige}

\begin{corrige}[Exercice 3 — Vocabulaire : associez]
\begin{center}
\begin{tabularx}{\linewidth}{|X|X|X|X|}
\hline
\rowcolor{popblueL} \textbf{N°} & \textbf{Caractères} & \textbf{Pinyin} & \textbf{Traduction} \\
\hline
1 & \zh{预防} & c. yùfáng & iii. prévention \\
\hline
2 & \zh{蚊帐} & b. wénzhàng & i. moustiquaire \\
\hline
3 & \zh{发烧} & a. fāshāo & ii. fièvre \\
\hline
4 & \zh{污染} & d. wūrǎn & iv. pollution \\
\hline
5 & \zh{保护} & e. bǎohù & v. protéger / protection \\
\hline
\end{tabularx}
\end{center}
\textbf{Correspondances :} 1 - c - iii \quad 2 - b - i \quad 3 - a - ii \quad 4 - d - iv \quad 5 - e - v
\end{corrige}

\begin{corrige}[Exercice 4 — Pinyin et tons : écoutez et identifiez]
\begin{center}
\begin{tabularx}{\linewidth}{|X|X|X|}
\hline
\rowcolor{popblueL} \textbf{Mot} & \textbf{Pinyin (tons)} & \textbf{Ton 1er car.} \\
\hline
\zh{疟疾} & nuèjí & 4 \\
\hline
\zh{回收} & huíshōu & 2 \\
\hline
\zh{垃圾} & lājī & 1 \\
\hline
\zh{预防} & yùfáng & 4 \\
\hline
\zh{发烧} & fāshāo & 1 \\
\hline
\zh{蚊帐} & wénzhàng & 2 \\
\hline
\zh{塑料} & sùliào & 4 \\
\hline
\zh{污染} & wūrǎn & 1 \\
\hline
\zh{保护} & bǎohù & 3 \\
\hline
\zh{医院} & yīyuàn & 1 \\
\hline
\end{tabularx}
\end{center}
\emph{Note :} Pour \zh{疟疾}, le premier caractère \zh{疟} se prononce \emph{nuè} (4e ton). Pour \zh{垃圾}, \zh{垃} est 1er ton (\emph{lā}), \zh{圾} est ton neutre (souvent noté 1er ton en citation). Pour \zh{医院}, \zh{医} est 1er ton (\emph{yī}).
\end{corrige}

\courssub{Je m'entraîne}

\begin{corrige}[Exercice 5 — Traduction : du français vers le chinois (Thème)]
\textbf{1.} 我们应该分类垃圾，保护环境。 (Wǒmen yīnggāi fēnlèi lājī, bǎohù huánjìng.) \\
\emph{Variante :} 我们得把垃圾分类，以保护环境。

\textbf{2.} 预防疟疾，要睡蚊帐。 (Yùfáng nuèjí, yào shuì wénzhàng.) \\
\emph{Structure :} Sujet (thème) + 要 + Verbe.

\textbf{3.} 我认为回收塑料非常重要。 (Wǒ rènwéi huíshōu sùliào fēicháng zhòngyào.) \\
\emph{Variante :} 我觉得塑料回收很重要。

\textbf{4.} 医生建议我多喝水，好好休息。 (Yīshēng jiànyi wǒ duō hē shuǐ, hǎohao xiūxi.) \\
\emph{Grammaire :} \zh{建议} + Objet + Verbe. \zh{好好} + Verbe = bien faire qqch.

\textbf{5.} 首先，要减少垃圾；其次，要回收利用。 (Shǒuxiān, yào jiǎnshǎo lājī; qícì, yào huíshōu lìyòng.) \\
\emph{Connecteurs :} \zh{首先}... \zh{其次}... \zh{回收利用} = recycler et réutiliser.
\end{corrige}

\begin{corrige}[Exercice 6 — Transformation de phrases]
\textbf{1.} \zh{我认为我们要保护环境。} (Wǒ rènwéi wǒmen yào bǎohù huánjìng.) \\
\emph{Structure :} \zh{我认为} (Je pense que) + Phrase affirmative.

\textbf{2.} \zh{回收垃圾非常重要。} (Huíshōu lājī fēicháng zhòngyào.) \\
\emph{Remplacement :} \zh{很有用} $\rightarrow$ \zh{非常重要}.

\textbf{3.} \zh{因为他得了疟疾，所以他发烧了。} (Yīnwèi tā déle nuèjí, suǒyǐ tā fāshāo le.) \\
\emph{Structure cause-effet :} \zh{因为}... \zh{所以}... \zh{得了} = attraper (maladie). \zh{了} = aspect accompli.

\textbf{4.} \zh{你觉得睡蚊帐怎么样？} (Nǐ juéde shuì wénzhàng zěnmeyàng?) \\
\emph{Structure avis :} \zh{你觉得} + Sujet/Action + \zh{怎么样？} (Supprimer \zh{应该}).

\textbf{5.} \zh{总之，塑料污染很严重。} (Zǒngzhī, sùliào wūrǎn hěn yánzhòng.) \\
\emph{Marque de conclusion :} \zh{总之} en tête de phrase.
\end{corrige}

\begin{corrige}[Exercice 7 — Texte à trous : dialogue à l'hôpital]
\begin{textebox}[Dialogue complété]
\textbf{医生}：你怎么了？\quad
\textbf{病人}：医生，我头疼，还 \zh{发烧} (fāshāo)。\quad
\textbf{医生}：你有没有其他的症状？\quad
\textbf{病人}：我 \zh{也} (yě) froid, je transpire.\quad
\textbf{医生}：可能是 \zh{疟疾} (nuèjí). Il faut faire un test de sang.\quad
\textbf{病人}：Ah ? Je dois \zh{住院} (zhùyuàn) / \zh{留院} (liúyuàn) ?\quad
\textbf{医生}：Oui, tu \zh{应该} (yīnggāi) / \zh{得} (děi) / \zh{必须} (bìxū) prendre des médicaments et bien te reposer. Et à la maison, tu \zh{应该} (yīnggāi) / \zh{得} (děi) dormir sous une moustiquaire.\quad
\textbf{病人}：D'accord, merci docteur.
\end{textebox}
\emph{Vocabulaire clé :} \zh{头疼} (tóuténg, mal de tête), \zh{症状} (zhèngzhuàng, symptôme), \zh{冷} (lěng, avoir froid), \zh{出汗} (chūhàn, transpirer), \zh{验血} (yànxiě, test sanguin), \zh{吃药} (chī yào, prendre médicaments), \zh{休息} (xiūxi, se reposer).
\end{corrige}

\begin{corrige}[Exercice 8 — Remise en ordre : structure d'un exposé]
Ordre logique : \textbf{6 $\rightarrow$ 2 $\rightarrow$ 4 $\rightarrow$ 5 $\rightarrow$ 1 $\rightarrow$ 3}

\begin{enumerate}
    \item \zh{各位同学，各位老师，大家早上好。} (Salutation formelle)
    \item \zh{大家好，今天我的话题是垃圾回收。} (Annonce du sujet)
    \item \zh{首先，回收可以减少污染。} (Argument 1 : \zh{首先})
    \item \zh{其次，回收能节约资源。} (Argument 2 : \zh{其次})
    \item \zh{总之，保护环境从我做起。} (Conclusion : \zh{总之})
    \item \zh{谢谢大家！} (Remerciement)
\end{enumerate}
\end{corrige}

\courssub{Je me prépare au Bac}

\begin{corrige}[Exercice 9 — Type Bac : Compréhension orale/écrite]
\textbf{Questions de compréhension :}
\begin{enumerate}
    \item \textbf{Sujet :} Comment protéger l'environnement / La protection de l'environnement (\zh{环境保护}).
    \item \textbf{Deux arguments (au choix) :}
        \begin{itemize}
            \item Réduire l'utilisation des sacs en plastique (\zh{减少使用塑料袋}) car pollution grave et décomposition longue.
            \item Faire le tri sélectif (\zh{垃圾分类}) : poubelle bleue (recyclables) / verte (cuisine).
            \item Économiser l'eau et l'électricité (\zh{节约水电}) : éteindre lumière/robinet.
        \end{itemize}
    \item \textbf{Déchets de cuisine (\zh{厨余垃圾}) :} Dans la poubelle verte (\zh{绿色垃圾桶}).
    \item \textbf{Formule de conclusion :} \zh{总之} (En conclusion / En somme).
\end{enumerate}

\textbf{Questions de langue :}
\begin{enumerate}
    \setcounter{enumi}{4}
    \item \textbf{Connecteurs logiques (exemples) :}
        \begin{itemize}
            \item \zh{首先} (shǒuxiān) : Premièrement / Tout d'abord.
            \item \zh{其次} (qícì) : Deuxièmement / Ensuite.
            \item \zh{再次} (zàicì) : Troisièmement / De plus.
            \item \zh{因为}... (yīnwèi) : Parce que / Comme (cause).
            \item \zh{所以} (suǒyǐ) : Donc / Par conséquent (conséquence).
        \end{itemize}
    \item \textbf{Robinet :} \zh{水龙头} (shuǐlóngtóu).
    \item \textbf{Phrase négative :} \zh{我们不要使用塑料袋。} (Wǒmen bùyào shǐyòng sùliào dài.) \\
    \emph{Attention :} \zh{不要} (bùyào) = ne pas / il ne faut pas (impératif négatif / conseil). Ne pas confondre avec \zh{不} (bù) simple.
\end{enumerate}

\emph{Barème indicatif (sur 10) :} Compréhension 1 (1 pt), 2 (2 pts), 3 (1 pt), 4 (1 pt) = 5 pts. Langue 5 (2 pts), 6 (1 pt), 7 (2 pts) = 5 pts.
\end{corrige}

\begin{corrige}[Exercice 10 — Type Bac : Expression orale guidée (Modèle de rédaction)]
\textbf{Proposition de texte d'exposé (environ 170 caractères) :}

\begin{textebox}[Exposé modèle : 我们应该怎么预防疟疾？]
各位同学，大家好。\\
今天，我的话题是我们应该怎么预防疟疾。\\
\textbf{我认为} 预防疟疾非常重要。\\
\textbf{首先}，我们要睡 \zh{蚊帐} (wénzhàng)，\textbf{因为} \zh{蚊子} (wénzi) 传播疟疾，\textbf{所以} 蚊帐能防叮咬。\\
\textbf{其次}，要清理 \zh{环境} (huánjìng)，不要有积水，\textbf{如果} 环境干净，\textbf{就} 少蚊子。\\
\textbf{再次}，\textbf{如果} \zh{发烧} (fāshāo)，要马上去 \zh{医院} (yīyuàn) 吃 \zh{药} (yào)。\\
\textbf{总之}，预防疟疾人人有责，从我做起。\\
谢谢大家！
\end{textebox}

\textbf{Vérification des contraintes :}
\begin{itemize}
    \item \textbf{Formules fonctionnelles (4) :} \zh{我认为}、\zh{首先}、\zh{其次}、\zh{总之} (aussi \zh{今天我的话题是...}).
    \item \textbf{Connecteurs logiques (3) :} \zh{因为...所以...}、\zh{如果...就...} (x2).
    \item \textbf{Vocabulaire leçon (7) :} \zh{蚊帐}、\zh{蚊子}、\zh{发烧}、\zh{医院}、\zh{预防} (titre/texte)、\zh{药}、\zh{环境}.
    \item \textbf{Structure :} Salutation $\rightarrow$ Sujet $\rightarrow$ 3 arguments $\rightarrow$ Conclusion $\rightarrow$ Remerciement. ✓
\end{itemize}

\textbf{Points de vigilance pour l'oral :}
\begin{itemize}
    \item Tons : \zh{疟疾} (nuèjí, 4-2), \zh{蚊帐} (wénzhàng, 2-4), \zh{蚊子} (wénzi, 2-ton neutre), \zh{发烧} (fāshāo, 1-1), \zh{医院} (yīyuàn, 1-4), \zh{环境} (huánjìng, 2-4).
    \item Liaison : \zh{因为}... \zh{所以}... ne pas faire de pause longue entre les deux.
    \item Intonation : Montante sur \zh{怎么样？}, descendante sur affirmations.
\end{itemize}
\end{corrige}

\begin{corrige}[Exercice 11 — Type Bac : Thème et Version]
\textbf{A. Thème (Français $\rightarrow$ Chinois)}

\begin{textebox}[Traduction proposée]
在喀麦隆，塑料垃圾污染是一个大问题。\\
为了保护环境，政府鼓励垃圾分类。\\
每个学生都要参加：塑料瓶要放进蓝色垃圾桶，厨余垃圾放进绿色垃圾桶。\\
如果大家都努力，我们的城市会更干净。
\end{textebox}

\textbf{Pinyin :} \\
Zài Kāmàilóng, sùliào lājī wūrǎn shì yí gè dà wèntí. \\
Wèile bǎohù huánjìng, zhèngfǔ gǔlì lājī fēnlèi. \\
Měi gè xuéshēng dōu yào cānjiā: sùliào píng yào fàng jìn lán sè lājī tǒng, chúyú lājī fàng jìn lǜ sè lājī tǒng. \\
Rúguǒ dàjiā dōu nǔlì, wǒmen de chéngshì huì gèng gānjìng.

\textbf{Points grammaticaux attendus :}
\begin{itemize}
    \item \zh{在喀麦隆} (Localisateur en tête).
    \item \zh{是一个大问题} (Jugement avec \zh{是}).
    \item \zh{为了}... \zh{政府鼓励}... (Finalité + Sujet action).
    \item \zh{每个学生都要参加} (Quantificateur \zh{每} + \zh{都}).
    \item \zh{要放进} (Obligation \zh{要} + Verbe directionnel \zh{放进}).
    \item \zh{如果}... \zh{就}... (Conditionnel hypothétique).
\end{itemize}

\textbf{B. Version (Chinois $\rightarrow$ Français)}

\begin{textebox}[Traduction française]
\textbf{Journaliste} : Camarade (collègue/étudiant), penses-tu qu'il faille interdire les sacs en plastique ?
\textbf{Étudiant} : Je suis d'accord. Parce que les sacs en plastique sont très difficiles à décomposer, ils polluent le sol et les rivières.
\textbf{Journaliste} : Mais les sacs en plastique sont pratiques, et bon marché.
\textbf{Étudiant} : Bien qu'ils soient pratiques, nous pouvons utiliser des sacs en tissu. Les sacs en tissu peuvent être réutilisés, c'est écologique.
\textbf{Journaliste} : Ta suggestion est bonne. En conclusion, protéger l'environnement nécessite les efforts de tous.
\textbf{Étudiant} : Exact, en commençant par moi, utiliser moins de sacs en plastique.
\end{textebox}

\textbf{Notes de traduction :}
\begin{itemize}
    \item \zh{同学} : « Camarade », « jeune homme/jeune fille », « étudiant » selon contexte. Ici « Camarade » ou « Jeune homme » (journaliste s'adresse à un élève).
    \item \zh{禁止} (jìnzhǐ) : Interdire (officiel).
    \item \zh{分解} (fēnjiě) : Se décomposer (chimiquement/biologiquement).
    \item \zh{土壤和河流} (tǔráng hé héliú) : Le sol et les rivières / cours d'eau.
    \item \zh{虽然}... \zh{但/但是}... : Bien que... mais... (Concession). \zh{虽然方便} (ellipse du sujet « ils »).
    \item \zh{布袋} (bùdài) : Sac en tissu / tote bag.
    \item \zh{重复使用} (chóngfù shǐyòng) : Réutiliser / usage répété.
    \item \zh{环保} (huánbǎo) : Écologique / protection de l'environnement (abréviation de \zh{环境保护}).
    \item \zh{从我做起} (cóng wǒ qǐ) : En commençant par moi / à partir de moi (slogan standard).
    \item \zh{少用} (shǎo yòng) : Utiliser moins / réduire l'usage.
\end{itemize}

\emph{Barème indicatif (sur 10) :} Thème (5 pts) : Vocabulaire juste (1), Structures \zh{为了/要/如果/就} (2), Caractères/Pinyin corrects (2). Version (5 pts) : Sens général fidélité (2), Vocabulaire précis \zh{分解/重复使用/环保/从我做起} (2), Français naturel (1).
\end{corrige}

\newpage

\coursec{Fiche bilan}

\courssub{Fiche Bilan — Leçon 15 : Expression orale — Recyclage et Paludisme}

\begin{popfigure}
\centering
\begin{center}\tcbox[colback=popgreen,colframe=popgreen,coltext=white,arc=9pt,boxsep=4pt,left=6pt,right=6pt]{\bfseries\small Leçon 15 Recyclage \& Paludisme}\end{center}
\vspace{-2pt}
\begin{tcbraster}[raster columns=3,raster equal height=rows,raster column skip=6pt,raster row skip=6pt,enhanced,blankest]
\begin{tcolorbox}[enhanced,colback=popgreen,colframe=popgreen,coltext=white,boxrule=0.9pt,arc=3pt,left=4pt,right=4pt,top=3pt,bottom=3pt,boxsep=1pt,halign=left]\bfseries\scriptsize Lexique Recyclage\end{tcolorbox}
\begin{tcolorbox}[enhanced,colback=poporange,colframe=poporange,coltext=white,boxrule=0.9pt,arc=3pt,left=4pt,right=4pt,top=3pt,bottom=3pt,boxsep=1pt,halign=left]\bfseries\scriptsize Lexique Paludisme\end{tcolorbox}
\begin{tcolorbox}[enhanced,colback=poppurple,colframe=poppurple,coltext=white,boxrule=0.9pt,arc=3pt,left=4pt,right=4pt,top=3pt,bottom=3pt,boxsep=1pt,halign=left]\bfseries\scriptsize Dialogue Exposé\end{tcolorbox}
\begin{tcolorbox}[enhanced,colback=poppink,colframe=poppink,coltext=white,boxrule=0.9pt,arc=3pt,left=4pt,right=4pt,top=3pt,bottom=3pt,boxsep=1pt,halign=left]\bfseries\scriptsize Dialogue Débat\end{tcolorbox}
\begin{tcolorbox}[enhanced,colback=popgold,colframe=popgold,coltext=white,boxrule=0.9pt,arc=3pt,left=4pt,right=4pt,top=3pt,bottom=3pt,boxsep=1pt,halign=left]\bfseries\scriptsize Formules Fonctionnelles\end{tcolorbox}
\begin{tcolorbox}[enhanced,colback=popdark,colframe=popdark,coltext=white,boxrule=0.9pt,arc=3pt,left=4pt,right=4pt,top=3pt,bottom=3pt,boxsep=1pt,halign=left]\bfseries\scriptsize Tons \& Prononciation\end{tcolorbox}
\begin{tcolorbox}[enhanced,colback=popmuted,colframe=popmuted,coltext=white,boxrule=0.9pt,arc=3pt,left=4pt,right=4pt,top=3pt,bottom=3pt,boxsep=1pt,halign=left]\bfseries\scriptsize Jeux de Rôle Débat Guidé\end{tcolorbox}
\end{tcbraster}

\legende{Carte mentale de la Leçon 15 : axes lexicaux, communicatifs et phonologiques}
\end{popfigure}

\begin{bilan}

\paragraph{Lexique clé — Recyclage (回收 huíshōu)}

\begin{center}
\begin{tabularx}{\linewidth}{|X|X|X|X|}
\hline
\rowcolor{popblueL} \textbf{Caractères} & \textbf{Pinyin} & \textbf{Nature} & \textbf{Traduction} \\ \hline
回收 & huíshōu & verbe & recycler, récupérer \\ \hline
垃圾分类 & lājī fēnlèi & nom & tri des déchets \\ \hline
可回收物 & kě huíshōu wù & nom & matériaux recyclables \\ \hline
有害垃圾 & yǒuhài lājī & nom & déchets dangereux \\ \hline
厨余垃圾 & chúyú lājī & nom & déchets de cuisine (organiques) \\ \hline
其他垃圾 & qítā lājī & nom & autres déchets (résidus) \\ \hline
环保 & huánbǎo & adj/verbe & protection de l'environnement / écolo \\ \hline
减少 & jiǎnshǎo & verbe & réduire \\ \hline
再利用 & zài lìyòng & verbe & réutiliser \\ \hline
填埋场 & tiánmái chǎng & nom & décharge (enfouissement) \\ \hline
\end{tabularx}
\end{center}

\paragraph{Lexique clé — Paludisme (疟疾 nüèjí)}

\begin{center}
\begin{tabularx}{\linewidth}{|X|X|X|X|}
\hline
\rowcolor{poporangeL} \textbf{Caractères} & \textbf{Pinyin} & \textbf{Nature} & \textbf{Traduction} \\ \hline
疟疾 & nüèjí & nom & paludisme (malaria) \\ \hline
蚊子 & wénzi & nom & moustique \\ \hline
叮咬 & dīngyǎo & verbe & piquer (insecte) \\ \hline
传播 & chuánbō & verbe & transmettre, propager \\ \hline
症状 & zhèngzhuàng & nom & symptôme \\ \hline
发烧 & fāshāo & verbe & avoir de la fièvre \\ \hline
发冷 & fālěng & verbe & avoir des frissons \\ \hline
头痛 & tóutòng & nom/verbe & mal de tête / avoir mal à la tête \\ \hline
蚊帐 & wénzhàng & nom & moustiquaire \\ \hline
预防 & yùfáng & verbe/nom & prévention / prévenir \\ \hline
就医 & jiùyī & verbe & consulter un médecin \\ \hline
药物 & yàowù & nom & médicament \\ \hline
\end{tabularx}
\end{center}

\paragraph{Structures grammaticales et formules fonctionnelles (à connaître par cœur)}

\begin{center}
\begin{tabularx}{\linewidth}{|X|X|X|}
\hline
\rowcolor{poppurpleL} \textbf{Structure / Fonction} & \textbf{Exemple (Caractères + Pinyin)} & \textbf{Traduction} \\ \hline
Présenter un sujet : \zh{今天我要谈谈…} & 今天我要谈谈垃圾分类的重要性。\newline (Jīntiān wǒ yào tántan lājī fēnlèi de zhòngyàoxìng.) & Aujourd'hui, je vais parler de l'importance du tri des déchets. \\ \hline
Structurer : \zh{首先…，其次…，最后…} & 首先，我们要分类；其次，要减少垃圾；最后，要回收。\newline (Shǒuxiān, wǒmen yào fēnlèi; qícì, yào jiǎnshǎo lājī; zuìhòu, yào huíshōu.) & Premièrement, il faut trier ; deuxièmement, réduire les déchets ; enfin, recycler. \\ \hline
Donner son avis : \zh{我认为… / 我觉得…} & 我认为预防疟疾最重要的方法是睡蚊帐。\newline (Wǒ rènwéi yùfáng nüèjí zuì zhòngyào de fāngfǎ shì shuì wénzhàng.) & Je pense que la méthode la plus importante pour prévenir le paludisme est de dormir sous moustiquaire. \\ \hline
Exprimer accord : \zh{我同意… / 赞成} & 我同意你的看法，环保从我做起。\newline (Wǒ tóngyì nǐ de kànfǎ, huánbǎo cóng wǒ zuòqǐ.) & Je suis d'accord avec ton avis, la protection de l'environnement commence par moi. \\ \hline
Exprimer désaccord poli : \zh{我不太同意… / 恐怕…} & 我不太同意，光靠药物不够，还要消灭蚊子。\newline (Wǒ bù tài tóngyì, guāng kào yàowù bùgòu, hái yào xiāomiè wénzi.) & Je ne suis pas tout à fait d'accord, compter seulement sur les médicaments ne suffit pas, il faut aussi éradiquer les moustiques. \\ \hline
Relancer / Demander avis : \zh{你怎么看？ / 你有什么建议？} & 关于垃圾分类，你有什么建议？\newline (Guānyú lājī fēnlèi, nǐ yǒu shénme jiànyì?) & Concernant le tri des déchets, quelle est ta suggestion ? \\ \hline
Insister sur action : \zh{应该 / 必须 / 要} & 我们必须保护环境，不能乱扔垃圾。\newline (Wǒmen bìxū bǎohù huánjìng, bùnéng luàn rēng lājī.) & Nous devons protéger l'environnement, ne pas jeter les déchets n'importe où. \\ \hline
Conséquence : \zh{否则 / 要不然} & 要是不预防，否则就会得疟疾。\newline (Yàoshi bù yùfáng, fǒuzé jiù huì dé nüèjí.) & Si on ne prévient pas, sinon on attrapera le paludisme. \\ \hline
\end{tabularx}
\end{center}

\paragraph{Méthodes express}

\begin{itemize}
    \item \textbf{Préparer un exposé (3 min) :} 1) Choisir un angle précis (ex: \zh{校园垃圾分类} tri à l'école). 2) Écrire 3 mots-clés en chinois (pas de phrases entières). 3) Structurer : \zh{开头} (accroche) — \zh{中间} (2-3 arguments avec \zh{首先/其次}) — \zh{结尾} (appel à l'action \zh{让我们一起行动吧}). 4) S'entraîner à voix haute en surveillant les tons (3e ton, ton neutre).
    \item \textbf{Participer à un débat :} Écouter le mot-clé (\zh{关键词}) de l'interlocuteur. Utiliser \zh{我觉得…因为…} (Je pense que... parce que...). Ne pas rester silencieux : \zh{我不太懂，能再说一遍吗？} (Je ne comprends pas bien, pouvez-vous répéter ?).
    \item \textbf{Travailler les tons :} Isoler les mots pièges : \zh{回收 huíshōu (2-1)}, \zh{疟疾 nüèjí (4-2)}, \zh{蚊帐 wénzhàng (2-4)}, \zh{分类 fēnlèi (1-4)}. Lire les dialogues modèles en exagérant les contours tonaux (main qui monte/descend).
\end{itemize}

\end{bilan}

\begin{aretenir}[Checklist avant le Bac — Leçon 15]
$\square$ Je sais épeler et prononcer les 20 mots de vocabulaire (recyclage + paludisme) avec les tons justes. \\
$\square$ Je sais écrire les caractères clés : \zh{回收、分类、疟疾、蚊帐、预防、环保} (ordre des traits respecté). \\
$\square$ Je peux présenter un exposé de 1 à 2 minutes structuré (\zh{今天我要谈谈… 首先… 最后…}). \\
$\square$ Je maîtrise les formules pour donner mon avis (\zh{我认为…}) et exprimer (dés)accord (\zh{我同意/我不同意}). \\
$\square$ Je sais relancer mon camarade : \zh{你怎么看？ / 你有什么建议？}. \\
$\square$ Je distingue les paires tonales pièges : \zh{huíshōu (2-1)} vs \zh{huǐshòu (3-4)} ; \zh{nüèjí (4-2)} vs \zh{nièjí}. \\
$\square$ Je connais les 4 catégories de déchets en chinois (\zh{可回收物、有害垃圾、厨余垃圾、其他垃圾}). \\
$\square$ Je peux citer 3 symptômes du paludisme et 3 mesures de prévention en chinois. \\
$\square$ Je suis capable de participer à un débat guidé de 5 minutes sans passer au français. \\
$\square$ J'ai révisé les dialogues modèles (Exposé recyclage + Débat paludisme) par cœur. \\
$\square$ Je peux traduire un court paragraphe français $\rightarrow$ chinois sur l'environnement ou la santé (thème).
\end{aretenir}

\findecours

\end{document}
